\documentclass[a4paper,12pt]{article}

% ==================== PACKAGES ====================
\usepackage[french]{babel}      % Pour la langue française
\usepackage[T1]{fontenc}        % Encodage des polices
\usepackage[utf8]{inputenc}     % Encodage des caractères (accents)
\usepackage{geometry}           % Pour gérer les marges
\usepackage{amsmath, amssymb}   % Pour les mathématiques
\usepackage{enumitem}           % Pour personnaliser les listes
\usepackage{titlesec}           % Pour personnaliser les titres
\usepackage{xcolor}             % Pour les couleurs
\usepackage{hyperref}           % Pour les liens internes
\usepackage{amsthm}             % Pour les environnements théorèmes
\usepackage{graphicx}           % Pour les images (optionnel)
\usepackage{array}              % Pour les tableaux avancés
\usepackage{booktabs}           % Pour de meilleurs tableaux
\usepackage{eurosym}     % Pour le symbole euro
% Configuration des listes pour éviter les débordements
\setlist{leftmargin=*, itemsep=2pt, parsep=2pt}

% ==================== ENVIRONNEMENTS PERSONNALISÉS ====================
\newtheorem{definition}{Définition}[section]
\newtheorem{exemple}{Exemple}[section]

% ==================== MISE EN PAGE ====================
\geometry{
    top=2.5cm,
    bottom=2.5cm,
    left=2.5cm,
    right=2.5cm
}

% Personnalisation des titres de sections
\titleformat{\section}
    {\normalfont\Large\bfseries\color{blue!70!black}}
    {\thesection}{1em}{}

\titleformat{\subsection}
    {\normalfont\large\bfseries\color{blue!50!black}}
    {\thesubsection}{1em}{}

\titleformat{\subsubsection}
    {\normalfont\normalsize\bfseries\color{blue!40!black}}
    {\thesubsubsection}{1em}{}

% ==================== DOCUMENT ====================
\begin{document}

% ==================== PAGE DE TITRE ====================
\begin{titlepage}
    \centering
    \vspace*{2cm}
    
    {\Huge \textbf{Étude sur le Chronométrage}}\\[0.5cm]
    {\Large Méthodologie et Applications}\\[1.5cm]
    
    \vfill
    
    \begin{minipage}{0.7\textwidth}
        \centering
        \rule{\textwidth}{0.5pt}\\[0.5cm]
        \textbf{Cours de Mesure du Travail et Facteur Humain}\\
        \vspace{0.3cm}
        \textit{Document de cours - Année académique 2025-2026}
    \end{minipage}
    
    \vfill
    
    {\large \today}
\end{titlepage}

% ==================== SOMMAIRE ====================
\tableofcontents
\newpage

% ==================== PARTIE 1 : INTRODUCTION ====================
\part*{Partie 1 : Introduction}
\addcontentsline{toc}{part}{Partie 1 : Introduction}

\section{Définition et objectifs de la mesure du travail}

\subsection{Définition}

La mesure du travail est une technique fondamentale de l'organisation industrielle. Elle consiste à déterminer le temps nécessaire à un opérateur qualifié pour accomplir une tâche donnée, selon une méthode définie et à un niveau de performance déterminé.

\begin{definition}
La \textbf{mesure du travail} est l'application de techniques visant à déterminer le temps que demande à un ouvrier qualifié l'exécution d'une tâche donnée, à un niveau de rendement bien défini.
\end{definition}

Cette définition soulève trois questions essentielles :
\begin{itemize}
    \item \textbf{Faire quoi ?} (définition précise de la tâche)
    \item \textbf{Avec quoi ?} (moyens et équipements utilisés)
    \item \textbf{Comment ?} (méthode opératoire)
\end{itemize}

\subsection{Objectifs de la mesure du travail}

La mesure du travail poursuit plusieurs objectifs stratégiques et opérationnels :

\begin{itemize}
    \item \textbf{Définir les prix et les coûts} : établir les coûts de revient des produits
    \item \textbf{Motiver les opérateurs} : fixer des objectifs réalistes et équitables
    \item \textbf{Comparer les possibilités de conception} : évaluer différentes solutions techniques
    \item \textbf{Établir les plannings} : planifier la production
    \item \textbf{Contrôler la stabilisation des postes de travail} : suivre la performance
    \item \textbf{Fixer le temps de fabrication} : déterminer le temps nécessaire à l'élaboration d'un produit
    \item \textbf{Étudier, diminuer et éliminer les temps improductifs}
    \item \textbf{Apprécier la performance} : évaluer l'efficacité des opérateurs et des équipes
\end{itemize}

\subsection{Productivité industrielle}

La productivité est le rapport du produit obtenu (production) aux ressources utilisées pour l'obtenir :

\[
\text{Productivité} = \frac{\text{Production}}{\text{Ressources utilisées}}
\]

La mesure de la productivité peut être de trois types :

\begin{itemize}
    \item \textbf{Partielle} : utilisation d'un seul intrant (ex : productivité du travail)
    \item \textbf{Multifactorielle} : utilisation de plusieurs intrants
    \item \textbf{Globale} : utilisation de tous les intrants
\end{itemize}

\section{Place du chronométrage dans l'étude du travail}

\subsection{Les deux piliers de l'étude du travail}

L'étude du travail repose sur deux piliers complémentaires :

\begin{enumerate}
    \item \textbf{L'étude des méthodes} : consiste à examiner les méthodes existantes et envisagées d'un travail, afin de mettre au point et de faire appliquer des méthodes d'exécution plus efficaces et de diminuer les coûts.
    
    \item \textbf{La mesure du travail} : a pour objectif de quantifier le temps nécessaire à l'exécution d'une tâche selon la méthode retenue.
\end{enumerate}

\subsection{Objectifs de l'étude des méthodes}

L'étude des méthodes vise à :
\begin{itemize}
    \item Éliminer les déplacements inutiles
    \item Réduire les éléments à non valeur ajoutée
    \item Éliminer les mouvements superflus
    \item Combiner des éléments de travail
    \item Réduire la fatigue des opérateurs
    \item Améliorer l'aménagement des lieux de travail
    \item Améliorer la conception des outils et du matériel utilisés
\end{itemize}

\subsection{Les techniques de mesure du travail}

\begin{table}[ht]
\centering
\begin{tabular}{|p{5cm}|p{5cm}|p{5cm}|}
\hline
\multicolumn{3}{|c|}{\textbf{Techniques d'étude des temps}} \\
\hline
\multicolumn{2}{|c|}{\textbf{Techniques par observation directe}} & \textbf{Techniques par mesure indirecte} \\
\cline{1-2}
Chrono-analyse & Mesure par échantillonnage & Prédétermination des temps \\
& & Données standards \\
\hline
\end{tabular}
\caption{Les principales techniques de mesure du travail}
\end{table}

\subsection{Place du chronométrage}

Le chronométrage est la technique de mesure directe la plus utilisée. Il intervient après la définition de la méthode optimale. Il permet de fixer des temps de référence qui serviront de base à :

\begin{itemize}
    \item La gestion de la production
    \item L'évaluation des performances
    \item Le calcul des coûts
    \item La définition des objectifs
\end{itemize}

\begin{center}
\fbox{\begin{minipage}{0.85\textwidth}
\centering
\textbf{Principe fondamental} \\
Le chronométrage mesure \textit{ce qui existe}, \\
l'étude des méthodes détermine \textit{ce qui devrait être}.
\end{minipage}}
\end{center}

\section{Historique et évolution : de Taylor aux temps prédéterminés}

\subsection{Frederick Winslow Taylor (1856-1915)}

Ingénieur américain, Taylor est le promoteur le plus connu de l'Organisation Scientifique du Travail (OST). Il est le premier à systématiser :

\begin{itemize}
    \item La décomposition des tâches en éléments simples
    \item L'utilisation du chronométrage pour mesurer les temps
    \item La définition de temps alloués pour chaque tâche
\end{itemize}

\begin{quote}
\textit{« Il existe toujours une méthode et un outil qui permettent un travail plus rapide et de meilleure qualité que les autres. »} \\
— Frederick Winslow Taylor
\end{quote}

\subsection{Frank et Lillian Gilbreth (années 1920)}

Ce couple de chercheurs a poussé l'analyse plus loin en décomposant chaque geste élémentaire du travailleur en mouvements de base appelés \textbf{therbligs}.

\begin{definition}
Les \textbf{therbligs} sont une décomposition du travail en 17 mouvements élémentaires de base. Le nom est une inversion de « Gilbreth » avec une transposition du « th ».
\end{definition}

Les principaux therbligs incluent :
\begin{itemize}
    \item Chercher, Trouver, Choisir
    \item Saisir, Transporter, Tenir
    \item Positionner, Assembler, Désassembler
    \item Utiliser, Attendre, Reposer
\end{itemize}

\subsection{L'avènement des temps prédéterminés (1940-1970)}

Face aux limites du chronométrage (subjectivité du jugement d'allure, nécessité d'un poste existant), des méthodes de temps prédéterminés apparaissent :

\begin{enumerate}
    \item \textbf{1948 : MTM (Methods-Time Measurement)} \\
    Développé par Maynard, Stegemerten et Schwab. Le MTM attribue des temps standards à des mouvements élémentaires codifiés. L'unité utilisée est le TMU (Time Measurement Unit) = $10^{-5}$ heure.
    
    \item \textbf{1950 : MTS (Methods Time Standards)} \\
    Développé par la General Electric.
    
    \item \textbf{1966 : MODAPTS} \\
    Développé par G.C. Hyde. (Modular Arrangement of Predetermined Time Standards)
    
    \item \textbf{1974 : MOST (Maynard Operation Sequence Technique)} \\
    Créé par Kjell Zandin, cette méthode simplifie et accélère l'analyse des temps. Une heure de cycle peut être chiffrée en une journée d'étude.
    
    \item \textbf{Autres méthodes} : WORK FACTOR (Philips), BTE (Bureau des Temps Élémentaires).
\end{enumerate}

\subsection{Tableau chronologique synthétique}

\begin{table}[ht]
\centering
\begin{tabular}{|c|c|c|}
\hline
\textbf{Année} & \textbf{Auteurs / Organismes} & \textbf{Apport} \\
\hline
1900-1915 & Frederick W. Taylor & Organisation Scientifique du Travail (OST) \\
1920 & Frank et Lillian Gilbreth & Therbligs (décomposition des mouvements) \\
1948 & H.B. Maynard et al. & MTM 1 \\
1950 & General Electric & MTS \\
1965 & H.B. Maynard & MTM 2 \\
1966 & G.C. Hyde & MODAPTS \\
1974 & Kjell Zandin & MOST \\
1975-1984 & Divers & MTM 3, MTM UAS, MTM MEK, MTM SAM \\
\hline
\end{tabular}
\caption{Évolution historique des méthodes de mesure du temps}
\end{table}

\subsection{Évolution contemporaine}

Aujourd'hui, le chronométrage et les temps prédéterminés sont utilisés de manière complémentaire, en fonction :

\begin{itemize}
    \item Du contexte de l'étude
    \item Des délais disponibles
    \item Des objectifs poursuivis
    \item Du niveau de précision requis
\end{itemize}

Les tendances actuelles sont à :
\begin{itemize}
    \item L'intégration des outils numériques (logiciels de capture vidéo, analyses automatisées)
    \item La prise en compte des facteurs ergonomiques dans la fixation des temps
    \item Une approche participative associant les opérateurs à la démarche
    \item La constitution de bases de données internes d'entreprises
\end{itemize}

\subsection{Organisation du cours}

Ce cours est structuré en plusieurs chapitres :

\begin{enumerate}
    \item Introduction à la mesure du travail
    \item Le temps dans l'industrie : unités et typologie
    \item La courbe d'apprentissage et l'évolution des temps
    \item Le chronométrage : méthodologie et mise en œuvre
    \item Le jugement d'allure et la détermination du temps standard
    \item Le dépouillement et l'exploitation des résultats
    \item Les observations instantanées
    \item Les outils de représentation des temps (simogramme, analyse de déroulement)
    \item Les temps prédéterminés (MTM, MOST)
    \item Interface Homme-Machine
\end{enumerate}

\subsection{Synthèse de l'introduction}

\begin{center}
\fbox{\begin{minipage}{0.9\textwidth}
\textbf{Points clés à retenir :}

\begin{itemize}
    \item La mesure du travail détermine le temps nécessaire à l'exécution d'une tâche
    \item Elle se distingue de l'étude des méthodes qui optimise la façon de travailler
    \item Le chronométrage est la technique de mesure directe la plus utilisée
    \item L'histoire de la mesure du travail a connu des étapes clés : Taylor, Gilbreth, MTM, MOST
    \item Les temps prédéterminés permettent de chiffrer des temps avant la création du poste
\end{itemize}
\end{minipage}}
\end{center}

% ==================== FIN DU DOCUMENT ====================
% ==================== PARTIE 1 : FONDAMENTAUX ET PRÉREQUIS ====================
\part*{Partie 1 : Fondamentaux et Prérequis}
\addcontentsline{toc}{part}{Partie 1 : Fondamentaux et Prérequis}

\section{Le Temps dans l’Industrie}

\subsection{Unités de temps}

La mesure du temps dans l'industrie nécessite des unités adaptées aux calculs et aux systèmes informatiques. Deux grandes familles d'unités coexistent.

\subsubsection{Unités classiques (système sexagésimal)}

Ces unités sont issues du système traditionnel de mesure du temps :

\begin{table}[ht]
\centering
\begin{tabular}{|c|c|}
\hline
\textbf{Unité} & \textbf{Symbole} \\
\hline
Mois & m \\
Semaine & sem \\
Jour & j \\
Heure & h \\
Minute & min \\
Seconde & s \\
\hline
\end{tabular}
\caption{Unités classiques de temps}
\end{table}

\textbf{Inconvénient :} Les calculs sont complexes en raison des conversions ($1\ h = 60\ min$, $1\ min = 60\ s$), ce qui les rend peu pratiques pour les opérations arithmétiques fréquentes en gestion de production.

\subsubsection{Unités décimales}

Pour faciliter les calculs, l'industrie utilise des unités décimales basées sur l'heure ou la minute :

\begin{table}[ht]
\centering
\begin{tabular}{|c|c|c|}
\hline
\textbf{Unité} & \textbf{Symbole} & \textbf{Équivalence} \\
\hline
Décitheure & dh & $10^{-1}\ h = 6\ min$ \\
Centiheure & ch & $10^{-2}\ h = 0,6\ min = 36\ s$ \\
Milliheure & mh & $10^{-3}\ h = 0,06\ min = 3,6\ s$ \\
Décimilliheure & dmh & $10^{-4}\ h = 0,006\ min = 0,36\ s$ \\
Centimilliheure & cmh & $10^{-5}\ h = 0,0006\ min = 0,036\ s$ \\
\hline
Déciminute & dmin & $10^{-1}\ min = 6\ s$ \\
Centiminute & cmin & $10^{-2}\ min = 0,6\ s$ \\
\hline
\end{tabular}
\caption{Unités décimales de temps}
\end{table}

\begin{center}
\fbox{\begin{minipage}{0.9\textwidth}
\centering
\textbf{Remarque importante :} Le TMU (Time Measurement Unit), utilisé dans la méthode MTM, correspond à 1 cmh ($10^{-5}$ heure).
\end{minipage}}
\end{center}

\subsubsection{Table de conversion}

\begin{table}[ht]
\centering
\begin{tabular}{|c|c|c|c|c|c|c|}
\hline
\textbf{cmh} & \textbf{dmh} & \textbf{h} & \textbf{cmin} & \textbf{min} & \textbf{cs} & \textbf{s} \\
\hline
100 000 & 10 000 & 1 & 6 000 & 60 & 360 000 & 3 600 \\
\hline
\end{tabular}
\caption{Table de conversion entre les différentes unités}
\end{table}

\subsection{Choix de l'unité}

Le choix de l'unité de temps n'est pas anodin. Il doit être guidé par plusieurs critères.

\subsubsection{Critères de sélection}

\begin{itemize}
    \item \textbf{Précision recherchée} : une tâche courte nécessitera une unité fine (cmh, dmh) alors qu'une tâche longue pourra utiliser la centiheure.
    
    \item \textbf{Domaine étudié} : le temps machine, le temps manuel, le temps d'assemblage n'ont pas les mêmes ordres de grandeur.
    
    \item \textbf{Facilité des calculs} : les unités décimales simplifient les opérations arithmétiques.
    
    \item \textbf{Minimisation des conversions} : il est préférable d'utiliser une seule unité pour l'ensemble de l'étude.
    
    \item \textbf{Technique d'obtention du temps} : le chronométrage direct, la vidéo, les tables MTM imposent parfois l'unité.
    
    \item \textbf{Compatibilité avec les systèmes informatiques} : la GPAO (Gestion de Production Assistée par Ordinateur) peut imposer une unité particulière.
\end{itemize}

\subsubsection{Exemple de choix}

\begin{exemple}
On vise à déterminer la charge d'un opérateur. La précision recherchée est la seconde. La GPAO qui fera les calculs de charge utilise des ch (centiheure). Les temps sont obtenus à partir de la vidéo. Ces temps sont exprimés en h:min:s. L'unité sera donc le ch. Il faudra faire les calculs de conversion pour tous les temps mesurés.
\end{exemple}

\subsection{Types de temps}

Les temps peuvent être classés selon trois dimensions : leur nature, leur fréquence et leur position.

\subsubsection{Nature des temps (relation du temps à la ressource)}

\begin{description}[labelwidth=4cm]
    \item[Temps manuel (Tm)] Temps pendant lequel l'opérateur est actif, réalisant des opérations manuelles.
    \item[Temps technologique (Tt)] Temps pendant lequel la machine travaille sans intervention de l'opérateur.
    \item[Temps technico-manuel (Ttm)] Temps pendant lequel l'opérateur et la machine travaillent simultanément (intervention sur machine en fonctionnement).
\end{description}

\begin{table}[ht]
\centering
\begin{tabular}{|c|c|c|}
\hline
\textbf{Nature du temps} & \textbf{Symbole} & \textbf{Exemple} \\
\hline
Temps manuel & Tm & Prendre une pièce, positionner, évacuer \\
Temps technologique & Tt & Séchage automatique, mise en forme machine \\
Temps technico-manuel & Ttm & Verrouiller un plateau, actionner une pédale \\
\hline
\end{tabular}
\caption{Les trois natures de temps}
\end{table}

\subsubsection{Fréquence des temps (relation du temps à la production)}

\begin{description}[labelwidth=4cm]
    \item[Temps cycliques] Séquences qui se répètent à chaque cycle de production. \\
    \textit{Exemple :} prendre la pièce, la positionner, l'usiner, l'évacuer.
    
    \item[Temps fréquentiels] Séquences qui ne se produisent pas à chaque cycle mais à intervalles réguliers ou aléatoires. \\
    \textit{Exemples :}
    \begin{itemize}
        \item Approvisionnement en matière première (1 fois par lot)
        \item Changement de bobine de fil (après épuisement)
        \item Réglage machine (changement de série)
        \item Contrôle qualité périodique
    \end{itemize}
\end{description}

\subsubsection{Position des temps (relation des temps entre ressources)}

Cette classification concerne l'articulation temporelle entre différentes ressources (opérateurs, machines).

\begin{description}[labelwidth=4cm]
    \item[Temps séquentiels] Les opérations se succèdent dans le temps sans chevauchement.
    
    \item[Temps simultanés] Les opérations se déroulent en parallèle sur différentes ressources.
    
    \item[Temps masqués] Certaines opérations peuvent être réalisées pendant qu'une autre ressource travaille automatiquement, ne contribuant pas à l'allongement du cycle.
\end{description}

\subsection{Application pratique : Simogramme}

Pour visualiser et analyser la position des temps, on utilise un outil graphique appelé \textbf{simogramme}.

\subsubsection{Définition et objectifs}

Le simogramme est une représentation visuelle qui permet de :
\begin{itemize}
    \item Visualiser un cycle de production par type de temps (manuel, technologique, technico-manuel)
    \item Étudier la mise en parallèle de ressources (machines, opérateurs)
    \item Analyser différentes solutions d'organisation
    \item Identifier les opportunités d'optimisation
\end{itemize}

\subsubsection{Étapes de construction d'un simogramme}

\begin{enumerate}
    \item Identifier les séquences ou opérations
    \item Recueillir les durées des opérations
    \item Identifier les ressources mobilisées
    \item Repérer les temps manuels (Tm), technologiques (Tt) et technico-manuels (Ttm)
    \item Placer les ressources en ordonnée
    \item Placer l'échelle des temps en abscisse
\end{enumerate}

\subsubsection{Exemple : Atelier de mise en plis de pantalon}

\begin{table}[ht]
\centering
\begin{tabular}{|l|c|}
\hline
\textbf{Opération} & \textbf{Durée (dmh)} \\
\hline
Prendre et poser pantalon sur la presse & 5 (Tm) \\
Positionner le pantalon sur la presse & 12 (Tm) \\
Verrouiller plateau, actionner pédale vapeur & 8 (Ttm) \\
Mise en forme automatique & 60 (Tt) \\
Actionner pédale de séchage & 5 (Ttm) \\
Séchage automatique & 30 (Tt) \\
Déverrouiller plateau & 5 (Ttm) \\
Évacuer le pantalon & 10 (Tm) \\
\hline
\textbf{Temps total du cycle} & \textbf{135 dmh} \\
\hline
\end{tabular}
\caption{Décomposition des temps pour un cycle de pressage}
\end{table}

\subsubsection{Analyse d'optimisation}

L'analyse du simogramme permet d'envisager des optimisations, comme l'augmentation du nombre de machines confiées à un opérateur.

\begin{itemize}
    \item \textbf{Configuration avec 1 machine} : temps de cycle = 135 dmh
    \item \textbf{Production horaire} : $10\,000 \div 135 = 74$ pantalons/heure
    \item \textbf{Configuration avec 2 machines} : avec un temps de déplacement supplémentaire de 5 dmh entre les machines
    \item \textbf{Nouveau temps de cycle} : 165 dmh
    \item \textbf{Production horaire} : $(10\,000 \div 165) \times 2 = 120$ pantalons/heure
    \item \textbf{Gain} : $120 - 74 = 46$ pantalons/heure (+62\%)
\end{itemize}

\subsection{Synthèse}

\begin{center}
\fbox{\begin{minipage}{0.9\textwidth}
\textbf{Points clés à retenir :}

\begin{itemize}
    \item Les unités de temps doivent être choisies en fonction de la précision et des outils utilisés
    \item Les temps se classent selon leur nature (manuel, technologique, technico-manuel)
    \item La fréquence distingue les temps cycliques des temps fréquentiels
    \item La position des temps (séquentielle, simultanée, masquée) influence l'optimisation
    \item Le simogramme est un outil essentiel pour visualiser et optimiser les cycles
\end{itemize}
\end{minipage}}
\end{center}

% ==================== FIN DE LA PARTIE 1 ====================
% ==================== PARTIE 2 : CONDITIONS PRÉALABLES AU CHRONOMÉTRAGE ====================

\section{Conditions Préalables au Chronométrage}

Avant d'effectuer un chronométrage valable, il est essentiel de s'assurer que trois éléments sont stabilisés : l'opérateur, le poste de travail et le chronométreur lui-même. Sans ces conditions préalables, les mesures obtenues seront non représentatives et inexploitables.

\subsection{Stabilisation de l'opérateur}

La stabilisation de l'opérateur, également appelée stabilisation du poste du point de vue humain, correspond à la régularité de son travail. Un opérateur est considéré comme stabilisé lorsqu'il exécute sa tâche avec constance et maîtrise.

\subsubsection{Caractéristiques d'un opérateur stabilisé}

\begin{itemize}
    \item Les temps d'exécution d'un travail répétitif présentent des écarts négligeables
    \item L'opérateur travaille à une vitesse régulière, aboutissant à un rendement constant
    \item Le mode opératoire est parfaitement maîtrisé et respecté
    \item L'opérateur utilise son matériel avec efficacité et automatisme
    \item La qualité du travail effectué est constante et conforme aux exigences
\end{itemize}

\subsubsection{La courbe d'apprentissage}

Avant d'atteindre un état stabilisé, l'opérateur passe par une phase d'apprentissage durant laquelle ses temps de production diminuent progressivement.

\begin{definition}
La \textbf{courbe d'apprentissage} modélise la réduction du temps nécessaire pour produire une unité à mesure que l'expérience augmente. À chaque doublement du nombre de répétitions de la tâche, sa durée de réalisation diminue d'un pourcentage constant $p$, appelé coefficient d'apprentissage.
\end{definition}

\paragraph{Loi de Wright (1963)}

La durée de réalisation de la tâche au $n^{\text{ème}}$ cycle est donnée par :

\[
T_{n} = T_1 \cdot n^{e} \quad \text{avec} \quad e = \frac{\ln p}{\ln 2}
\]

Où :
\begin{itemize}
    \item $T_{n}$ : durée de réalisation de la tâche au $n^{\text{ème}}$ cycle
    \item $T_1$ : durée de réalisation de la tâche pour la première fois
    \item $p$ : coefficient d'apprentissage (exemple : $p = 0,90$ pour un coefficient de 90\%)
\end{itemize}

\begin{table}[ht]
\centering
\begin{tabular}{|l|c|}
\hline
\textbf{Secteur d'activité} & \textbf{Coefficient d'apprentissage} \\
\hline
Aérospatiale & 70-80 \% \\
Construction navale & 80-85 \% \\
Fabrication électronique & 90-95 \% \\
Soudure & 90 \% \\
Construction de bâtiments & 70-90 \% \\
Transformation des matières premières & 93-96 \% \\
\hline
\end{tabular}
\caption{Coefficients d'apprentissage par secteur}
\end{table}

\paragraph{Loi de Caquot}

Pour la fabrication de grande série, après la phase de mise en route, on atteint un plateau. La loi de Caquot s'applique alors :

\[
T_n = T_i \cdot \left(\frac{Q_n}{Q_i}\right)^{1/4}
\]

Où :
\begin{itemize}
    \item $T_n$ : temps de la $n^{\text{ème}}$ unité
    \item $T_i$ : temps de la $i^{\text{ème}}$ unité
    \item $Q_n$ : rang de la $n^{\text{ème}}$ unité
    \item $Q_i$ : rang de la $i^{\text{ème}}$ unité
    \item $n > i$
\end{itemize}

\begin{exemple}
Une entreprise de BTP réalise 5 chantiers identiques. Le premier chantier dure 20 jours. Le coefficient d'apprentissage de l'équipe est estimé à 92\%.

\begin{align*}
T_2 &= 20 \times 2^{\ln 0,92 / \ln 2} = 18,4 \text{ jours} \\
T_5 &= 20 \times 5^{\ln 0,92 / \ln 2} = 16,48 \text{ jours}
\end{align*}
\end{exemple}

\subsubsection{Le taux d'aléas}

Le taux d'aléas permet de mesurer objectivement l'importance des irrégularités liées au travail et de préciser si le poste est stabilisé.

\[
\text{Taux d'aléas} = \frac{\text{Nombre de valeurs aberrantes}}{\text{Nombre total de mesures}} \times 100
\]

\begin{table}[ht]
\centering
\begin{tabular}{|c|c|}
\hline
\textbf{Taux d'aléas} & \textbf{Interprétation} \\
\hline
Moins de 50 \% & Élément stable \\
Plus de 50 \% & Élément non stable \\
\hline
\end{tabular}
\caption{Degrés de stabilisation}
\end{table}

\begin{exemple}
Pour un élément de travail, après 10 mesures chronométriques, on observe 3 valeurs aberrantes :
\[
\text{Taux d'aléas} = \frac{3}{10} \times 100 = 30\%
\]
L'élément est considéré comme stable.
\end{exemple}

\subsection{Stabilisation du poste de travail}

Le poste de travail doit être dans des conditions normales et stables pour que les mesures soient valables.

\subsubsection{Aménagement du poste}

Avant le chronométrage, il faut vérifier que l'aménagement du poste est définitif et optimisé. Toute modification ultérieure invaliderait les mesures. Les aménagements éventuels doivent être réalisés en collaboration avec :
\begin{itemize}
    \item L'agent de maîtrise du secteur concerné
    \item L'opérateur lui-même
    \item Le mécanicien d'entretien si nécessaire
\end{itemize}

\subsubsection{Alimentation du poste}

Pendant toute la durée du chronométrage, le poste doit être correctement alimenté. L'attente de travail constituerait une perturbation inacceptable.

\subsubsection{Travail correctement présenté et préparé}

Toutes les opérations antérieures doivent être correctement réalisées. Les éléments suivants doivent être vérifiés :
\begin{itemize}
    \item Étiquetage conforme
    \item Préparation des sous-ensembles
    \item Qualité des opérations préalables
\end{itemize}

\subsubsection{Matières premières}

Les opérations peuvent être perturbées par le comportement des matières. Il faut s'assurer que les matières utilisées sont conformes et ne présentent pas de spécificités anormales.

\subsubsection{Machine en bon état de fonctionnement}

\begin{itemize}
    \item Absence de pannes mécaniques
    \item Fonctionnement normal des équipements
    \item Réglages adaptés à la morphologie de l'opérateur
    \item Chaise réglée en harmonie avec la taille de l'opérateur
\end{itemize}

\subsubsection{Grade de qualité bien défini}

Les critères de qualité doivent être clairement définis avec leurs limites de tolérance. Le système de contrôle doit être précisé et compris par tous.

\subsubsection{Absence de dérangements extérieurs}

Pendant le chronométrage :
\begin{itemize}
    \item Le chronométreur ne doit pas être sollicité (téléphone, autres interventions)
    \item L'opérateur ne doit pas être interrompu par sa hiérarchie ou ses collègues
\end{itemize}

\subsubsection{Choix de la plage horaire}

Éviter de chronométrer :
\begin{itemize}
    \item Au début ou à la fin des périodes de travail
    \item À proximité des pauses
    \item Le lundi matin et le samedi
\end{itemize}

\subsection{Stabilisation du chronométreur}

Le chronométreur est un acteur essentiel de la mesure. Sa stabilisation concerne à la fois ses compétences techniques et son jugement.

\subsubsection{Compétences techniques}

\begin{itemize}
    \item \textbf{Précision dans la manipulation} : savoir démarrer et arrêter le chronomètre au bon moment (reconnaissance des "tops")
    \item \textbf{Connaissance des modes opératoires} : être capable de reconnaître d'éventuelles déviances
    \item \textbf{Maîtrise des outils} : utilisation correcte du chronomètre, de la feuille de relevés
    \item \textbf{Aptitude à la décomposition} : savoir découper le travail en séquences pertinentes
\end{itemize}

\subsubsection{Le jugement d'allure (JA)}

Le jugement d'allure est une compétence fondamentale du chronométreur. Il s'agit de la capacité à évaluer objectivement le rythme de travail de l'opérateur observé.

\begin{definition}
Le \textbf{jugement d'allure} est l'action par laquelle un observateur entraîné définit l'allure (le rythme) d'un opérateur par rapport à la représentation mentale qu'il a de celle d'un exécutant placé dans des conditions similaires, l'allure de ce dernier étant prise comme allure de référence (allure 100).
\end{definition}

\subsubsection{Échelles de jugement d'allure}

Plusieurs échelles existent. L'échelle la plus couramment utilisée en Europe est l'échelle BTE (Bureau des Temps Élémentaires) :

\begin{table}[ht]
\centering
\begin{tabular}{|c|c|}
\hline
\textbf{Allure} & \textbf{Interprétation} \\
\hline
80 & Allure minimale tolérable \\
90 & Allure lente mais acceptable \\
100 & Allure normale (référence) \\
110 & Allure rapide \\
120 & Allure maximale (opérateur très motivé et expérimenté) \\
\hline
\end{tabular}
\caption{Échelle BTE du jugement d'allure}
\end{table}

\begin{center}
\fbox{\begin{minipage}{0.85\textwidth}
\centering
\textbf{Conseil pratique :} Il est recommandé d'éviter de chronométrer des opérateurs dont l'allure est inférieure à 80 ou supérieure à 120, car il devient difficile d'apprécier l'allure avec précision. L'idéal se situe entre 80 et 110.
\end{minipage}}
\end{center}

\subsubsection{Tableau de correspondance entre échelles}

\begin{table}[ht]
\centering
\begin{tabular}{|c|c|c|c|}
\hline
\textbf{Bedaux} & \textbf{Maynard} & \textbf{BTE} & \textbf{Marche (km/h)} \\
\hline
75 & 100 & 80 & 4,5 \\
80 & 107 & 85 & 5,0 \\
85 & 113 & 90 & 5,4 \\
90 & 120 & 95 & 5,8 \\
95 & 127 & 100 & 6,2 \\
100 & 133 & 105 & 6,6 \\
105 & 140 & 110 & 7,0 \\
110 & 147 & 115 & 7,4 \\
115 & 153 & 120 & 7,8 \\
\hline
\end{tabular}
\caption{Correspondance entre les principales échelles de jugement d'allure}
\end{table}

\subsubsection{Éléments d'appréciation du jugement d'allure}

Pour évaluer l'allure, le chronométreur doit considérer plusieurs facteurs :

\begin{table}[ht]
\centering
\begin{tabular}{|l|c|c|}
\hline
\textbf{Facteur} & \textbf{Critère} & \textbf{Variation possible} \\
\hline
Habileté & Précision des mouvements & $-0,15$ à $+0,15$ \\
Activité & Vitesse d'exécution & $-0,22$ à $+0,13$ \\
Conditions de travail & Environnement du poste & $-0,07$ à $+0,06$ \\
Régularité & Constance du rythme & $-0,04$ à $+0,04$ \\
\hline
\end{tabular}
\caption{Facteurs d'évaluation du jugement d'allure - Méthode de Levelling}
\end{table}

\subsubsection{Méthode de Levelling}

La méthode de Levelling permet d'affiner le jugement d'allure en corrigeant le temps observé par une combinaison de facteurs.

\begin{exemple}
Temps chronométré = 15 dmh

Correctifs :
\begin{itemize}
    \item Habileté : Excellent = $+0,11$
    \item Activité : Bon = $+0,06$
    \item Conditions de travail : Passable = $-0,03$
    \item Stabilité : Bon = $+0,01$
\end{itemize}

Total des correctifs = $0,15$

Temps corrigé = $15 \times (1 + 0,15) = 17,25$ dmh
\end{exemple}

\subsubsection{Précision du chronométrage (bouclage)}

La précision du chronométreur peut être mesurée par la méthode du bouclage, qui compare le temps total mesuré avec le temps réel écoulé.

\[
\text{Écart} = \frac{\text{Temps total mesuré} - \text{Temps de bouclage}}{\text{Temps de bouclage}} \times 100
\]

\begin{table}[ht]
\centering
\begin{tabular}{|l|c|}
\hline
\textbf{Type de chronométrage} & \textbf{Écart maximum toléré} \\
\hline
Chronométrage avec JA & $\pm 5\%$ \\
Chronométrage sans JA & $\pm 2\%$ \\
\hline
\end{tabular}
\caption{Tolérances de précision pour le bouclage}
\end{table}

\subsection{Qualités professionnelles du chronométreur}

Un bon chronométreur doit posséder les qualités suivantes :

\begin{itemize}
    \item \textbf{L'esprit d'analyse} : savoir découvrir les éléments qui peuvent influencer le travail
    \item \textbf{La souplesse d'esprit} : savoir s'adapter à chaque caractère
    \item \textbf{La patience} : ne jamais s'énerver, recommencer les explications
    \item \textbf{L'honnêteté} : ne pas déformer les résultats
    \item \textbf{La loyauté} : reconnaître ses erreurs, avoir le courage de ses idées
    \item \textbf{L'objectivité} : respecter le climat de confiance avec l'opérateur
\end{itemize}

\subsection{Synthèse des conditions préalables}

\begin{center}
\fbox{\begin{minipage}{0.9\textwidth}
\textbf{Conditions préalables à vérifier avant tout chronométrage :}

\begin{itemize}
    \item[$\checkmark$] L'opérateur a dépassé sa phase d'apprentissage (courbe stabilisée)
    \item[$\checkmark$] Le poste de travail est aménagé, alimenté et fonctionne normalement
    \item[$\checkmark$] Les critères de qualité sont clairement définis et respectés
    \item[$\checkmark$] Le chronométreur maîtrise les techniques de mesure et le jugement d'allure
    \item[$\checkmark$] Les conditions d'observation sont favorables (absence d'interruptions)
    \item[$\checkmark$] La plage horaire choisie est représentative de l'activité normale
\end{itemize}
\end{minipage}}
\end{center}

% ==================== FIN DE LA PARTIE 2 ====================
% ==================== PARTIE 2 : MÉTHODOLOGIE DE LA MESURE ====================
\part*{Partie 2 : Méthodologie de la Mesure}
\addcontentsline{toc}{part}{Partie 2 : Méthodologie de la Mesure}

\section{Les Différents Types de Chronométrage}

Le chronométrage n'est pas une technique unique. Selon l'objectif poursuivi, la précision recherchée et le degré de stabilisation du poste, on distingue plusieurs types de chronométrage. Ces types se différencient principalement par le nombre de relevés effectués et par l'utilisation ou non du jugement d'allure.

\subsection{Classification générale}

On distingue deux grandes familles de chronométrage :

\begin{itemize}
    \item \textbf{Les chronométrages sans jugement d'allure} : utilisés pour le diagnostic, l'étude préliminaire ou lorsque le poste n'est pas encore stabilisé.
    \item \textbf{Les chronométrages avec jugement d'allure} : utilisés pour la fixation des temps de référence, la confirmation ou le contrôle.
\end{itemize}

\begin{table}[ht]
\centering
\begin{tabular}{|p{4.5cm}|p{3.5cm}|p{7cm}|}
\hline
\textbf{Type de chronométrage} & \textbf{Nb de relevés} & \textbf{Objectif principal} \\
\hline
Sans JA - Diagnostic & 2 à 3 par poste & Détecter les postes problématiques \\
Sans JA - Étude & 10 à 15 par élément & Analyser la stabilisation du poste \\
Avec JA - Fixation de tâche & 20 à 30 par élément & Déterminer les temps de référence \\
Avec JA - Confirmation & 20 à 30 par élément & Vérifier les temps existants \\
Avec JA - Contrôle & 20 à 30 par élément & Réviser les temps contestés \\
\hline
\end{tabular}
\caption{Les différents types de chronométrage}
\end{table}

\subsection{Chronométrage sans jugement d'allure}

\subsubsection{Chronométrage de diagnostic}

\begin{definition}
Le \textbf{chronométrage de diagnostic} est une première approche rapide qui consiste à effectuer 2 à 3 relevés chronométriques sur chacun des postes d'une chaîne ou d'un atelier.
\end{definition}

\textbf{Objectifs :}
\begin{itemize}
    \item Détecter le ou les postes susceptibles d'être la cause d'un mauvais fonctionnement
    \item Identifier les goulots d'étranglement
    \item Obtenir une première estimation des temps de cycle
    \item Orienter les investigations ultérieures
\end{itemize}

\textbf{Méthode :}
\begin{itemize}
    \item Relevés rapides sans décomposition fine du travail
    \item Pas de jugement d'allure
    \item Mesures prises sur l'ensemble des postes de la ligne
    \item Durée : quelques heures à une demi-journée
\end{itemize}

\subsubsection{Chronométrage d'étude}

\begin{definition}
Le \textbf{chronométrage d'étude} est une analyse approfondie réalisée sur un poste spécifique pour comprendre les causes des anomalies et y porter remède.
\end{definition}

\textbf{Objectifs :}
\begin{itemize}
    \item Déterminer un temps approximatif pour un élément ou groupe d'éléments de travail
    \item Étudier la stabilisation du poste
    \item Identifier les irrégularités et leurs causes
    \item Préparer les améliorations du mode opératoire
\end{itemize}

\textbf{Méthode :}
\begin{itemize}
    \item 10 à 15 relevés chronométriques par élément ou groupe d'éléments
    \item Décomposition fine du travail en séquences
    \item Pas de jugement d'allure
    \item Calcul du taux d'aléas pour chaque élément
\end{itemize}

\paragraph{Calcul du taux d'aléas}

\[
\text{Taux d'aléas} = \frac{\text{Nombre de mesures aberrantes}}{\text{Nombre total de mesures}} \times 100
\]

\begin{exemple}
Pour un élément "Prendre la pièce", on réalise 10 relevés :

\begin{center}
\begin{tabular}{|c|c|c|c|c|c|c|c|c|c|c|}
\hline
N° & 1 & 2 & 3 & 4 & 5 & 6 & 7 & 8 & 9 & 10 \\
\hline
Temps (cmn) & 18 & 19 & 20 & 21 & 19 & 20 & 18 & 35 & 19 & 20 \\
\hline
\end{tabular}
\end{center}

Valeurs aberrantes : 35 (anomalie détectée) \\
Taux d'aléas = $1/10 \times 100 = 10\%$ \\
L'élément est stable (taux < 50\%). Le temps moyen = 19,4 cmn.

Pour un autre élément "Positionner", avec plusieurs valeurs aberrantes :
Taux d'aléas = $7/10 \times 100 = 70\%$ \\
L'élément n'est pas stable. Il faut rechercher les causes.
\end{exemple}

\subsection{Chronométrage avec jugement d'allure}

\subsubsection{Chronométrage de fixation de tâche}

\begin{definition}
Le \textbf{chronométrage de fixation de tâche} est utilisé lorsque le poste est stabilisé. Il permet de déterminer des temps de référence fiables (temps à l'allure 100) qui serviront de base à la gestion de production.
\end{definition}

\textbf{Objectifs :}
\begin{itemize}
    \item Calculer les capacités prévisionnelles de production
    \item Équilibrer le travail aux postes
    \item Calculer les salaires aux pièces (travail à la tâche)
    \item Compléter le catalogue des temps de l'entreprise
    \item Calculer les coûts de revient de fabrication
\end{itemize}

\textbf{Méthode :}
\begin{itemize}
    \item Minimum 30 relevés chronométriques par élément
    \item Décomposition précise du travail
    \item Jugement d'allure pour chaque relevé
    \item Conversion des temps observés en temps à l'allure 100
    \item Application des coefficients de majoration
\end{itemize}

\paragraph{Déroulement type :}
\begin{enumerate}
    \item Préparation de la feuille de chronométrage avec décomposition du travail
    \item 30 à 50 cycles de mesure avec enregistrement des temps et des allures
    \item Dépouillement : calcul du temps observé moyen pour chaque élément
    \item Correction par le jugement d'allure pour obtenir le temps de base
    \item Application des majorations pour obtenir le temps standard
    \item Validation par la méthode du bouclage (écart < 5\%)
\end{enumerate}

\subsubsection{Chronométrage de confirmation}

\begin{definition}
Le \textbf{chronométrage de confirmation} vise à contrôler et valider les temps déterminés par le chronométrage de fixation de tâche, en les comparant avec les temps réels de production.
\end{definition}

\textbf{Cas de recours :}
\begin{itemize}
    \item \textbf{Demande du chronométreur ou du chef d'atelier} : constat d'écarts importants entre les temps prévus et les temps réels
    \item \textbf{Demande de l'opérateur ou du chef d'équipe} : les temps paraissent trop courts (surcharge)
    \item \textbf{Demande de la direction} : les temps paraissent trop longs (sous-productivité)
\end{itemize}

\textbf{Méthode :}
\begin{enumerate}
    \item Vérification des conditions de travail (identiques à l'étude initiale)
    \item Réalisation de 20 à 30 relevés par élément avec jugement d'allure
    \item Dépouillement et calcul des temps à l'allure 100
    \item Comparaison avec les temps existants
    \item Si écarts significatifs : refaire un chronométrage de fixation de tâche
\end{enumerate}

\subsubsection{Chronométrage de contrôle}

\begin{definition}
Le \textbf{chronométrage de contrôle} est une procédure formelle de révision des temps, souvent réalisée en présence d'un tiers indépendant, pour résoudre des contestations.
\end{definition}

\textbf{Contexte d'utilisation :}
\begin{itemize}
    \item Contestation interne entre services (ex: Bureau des Méthodes et Service Commercial)
    \item Contestation externe entre entreprises (sous-traitance)
    \item Litige entre opérateurs et direction
    \item Révision périodique des temps (longue durée d'accoutumance)
\end{itemize}

\textbf{Méthode :}
\begin{itemize}
    \item 20 à 30 relevés par élément
    \item Jugement d'allure par un chronométreur externe (ingénieur de conseil)
    \item Protocole strict et traçabilité des mesures
    \item Rapport formel avec conclusions
\end{itemize}

\subsection{Récapitulatif des types de chronométrage}

\begin{table}[ht]
\centering
\begin{tabular}{|p{3.5cm}|p{3.5cm}|p{2.5cm}|p{2.5cm}|p{2.5cm}|}
\hline
\textbf{Type} & \textbf{Utilisation} & \textbf{Nb relevés} & \textbf{JA} & \textbf{Précision} \\
\hline
Diagnostic & Détection des goulots & 2-3/postes & Non & Faible \\
Étude & Analyse de stabilisation & 10-15/élément & Non & Moyenne \\
Fixation de tâche & Établissement des temps & 20-30/élément & Oui & Élevée \\
Confirmation & Vérification & 20-30/élément & Oui & Élevée \\
Contrôle & Révision contestée & 20-30/élément & Oui & Très élevée \\
\hline
\end{tabular}
\caption{Récapitulatif des différents types de chronométrage}
\end{table}

\subsection{Choix du type de chronométrage}

Le choix du type de chronométrage dépend de plusieurs facteurs :

\begin{center}
\fbox{\begin{minipage}{0.9\textwidth}
\textbf{Critères de choix :}

\begin{itemize}
    \item[\textbf{1}] \textbf{Objectif de l'étude} : diagnostic rapide, fixation des temps, résolution de litige
    \item[\textbf{2}] \textbf{Stabilisation du poste} : poste nouveau ou mature
    \item[\textbf{3}] \textbf{Ressources disponibles} : temps, compétences du chronométreur
    \item[\textbf{4}] \textbf{Précision requise} : estimation ou valeur contractuelle
    \item[\textbf{5}] \textbf{Contexte relationnel} : climat social, acceptation des mesures
\end{itemize}

\textbf{Recommandation :} Commencer par un diagnostic rapide, puis approfondir avec une étude, et enfin fixer les temps lorsque le poste est stabilisé.
\end{minipage}}
\end{center}

\subsection{Exemple illustratif}

\begin{exemple}
Une entreprise de confection textile souhaite optimiser sa ligne de montage de pantalons.

\begin{enumerate}
    \item \textbf{Chronométrage de diagnostic} (2 heures) \\
    3 relevés sur chaque poste $\rightarrow$ identification du poste "pose de fermeture" comme goulot (temps moyen 2,5 fois supérieur aux autres)
    
    \item \textbf{Chronométrage d'étude} (1 journée) \\
    15 relevés sur le poste "pose de fermeture" $\rightarrow$ taux d'aléas = 65\% $\rightarrow$ problème de formation et d'approvisionnement
    
    \item \textbf{Améliorations} (1 semaine) \\
    Formation de l'opératrice, réorganisation de l'approvisionnement
    
    \item \textbf{Chronométrage de fixation de tâche} (1 journée) \\
    30 relevés avec JA $\rightarrow$ temps standard établi = 1,85 min
    
    \item \textbf{Chronométrage de confirmation} (1 mois plus tard) \\
    Vérification des temps réels $\rightarrow$ écart < 3\% $\rightarrow$ temps validé
\end{enumerate}
\end{exemple}

\subsection{Synthèse}

\begin{center}
\fbox{\begin{minipage}{0.9\textwidth}
\textbf{Points clés à retenir :}

\begin{itemize}
    \item Le choix du type de chronométrage dépend de l'objectif et de la stabilisation du poste
    \item Les chronométrages sans JA sont utilisés pour le diagnostic et l'étude préliminaire
    \item Les chronométrages avec JA sont utilisés pour fixer, confirmer ou contrôler les temps
    \item Le nombre de relevés augmente avec la précision recherchée
    \item Le jugement d'allure est indispensable pour les temps de référence
\end{itemize}
\end{minipage}}
\end{center}

% ==================== FIN DE LA PARTIE 2 (MÉTHODOLOGIE) ====================
% ==================== PARTIE 4 : MATÉRIELS ET PRÉPARATION ====================

\section{Matériels et Préparation}

Avant d'effectuer un chronométrage, il est essentiel de disposer du matériel adéquat et de préparer soigneusement l'étude. Une bonne préparation conditionne la fiabilité des résultats obtenus.

\subsection{Moyens matériels}

\subsubsection{Le chronomètre}

Le chronomètre est l'outil de base qui permet de mesurer le temps. Étant donné que les temps sont utilisés dans de multiples applications (calculs de coûts, planification, etc.), ils feront l'objet d'opérations arithmétiques fréquentes.

\begin{center}
\fbox{\begin{minipage}{0.85\textwidth}
\centering
\textbf{Choix du système de mesure :} \\
Le système sexagésimal (1 h = 60 min, 1 min = 60 s) n'est pas pratique pour les calculs. \\
L'industrie utilise un système décimal où la minute est divisée en centièmes (cmin) ou l'heure en centièmes (ch) ou millièmes (dmh).
\end{minipage}}
\end{center}

\paragraph{Chronomètre électronique (digital)}

Le chronomètre électronique est l'outil le plus utilisé aujourd'hui. Ses avantages sont nombreux :

\begin{itemize}
    \item Affichage sur écran à cristaux liquides
    \item Choix de l'unité de mesure (secondes/100, minutes/1000, DMH)
    \item Fonction de mémorisation (enregistrement de 9 mesures successives)
    \item Lecture possible sans arrêt du chronomètre
    \item Prix abordable
\end{itemize}

\paragraph{Chronomètre mécanique}

Bien que moins utilisé aujourd'hui, le chronomètre mécanique reste un outil fiable :

\begin{itemize}
    \item Graduation selon l'unité choisie (secondes, centièmes de minute, DMH)
    \item Modèle à aiguille rattrapante : deux aiguilles superposées permettant une lecture sans arrêter le chronométrage principal
    \item Nécessite une pratique plus alerte pour la lecture
\end{itemize}

\paragraph{Autres moyens de mesure}

\begin{itemize}
    \item \textbf{Caméra vidéo} : permet de revoir les séquences, utile pour les cycles courts ou complexes
    \item \textbf{Logiciels spécialisés} : capture et analyse automatique des temps
    \item \textbf{Tachymètre} : mesure la vitesse de rotation des machines (RPM)
\end{itemize}

\subsubsection{La planchette de chronométrage}

La planchette est destinée à recevoir la feuille de chronométrage. Elle doit être :
\begin{itemize}
    \item Légère pour ne pas fatiguer le chronométreur
    \item Rigide pour permettre d'écrire facilement
    \item Suffisamment grande pour accueillir les feuilles de relevés
\end{itemize}

\subsubsection{La feuille de relevés}

Il n'existe pas de feuille de chronométrage normalisée, mais une structure type est généralement utilisée :

\begin{table}[ht]
\centering
\begin{tabular}{|p{8cm}|p{6cm}|}
\hline
\textbf{En-tête} & \textbf{Contenu} \\
\hline
Identification & Article, modèle, atelier, opération \\
Caractéristiques & Matière, machine, vitesse (RPM), points/cm \\
Ressources & Opératrice, chronométreur, date \\
\hline
\end{tabular}
\caption{Informations générales de la feuille de chronométrage}
\end{table}

\begin{center}
\fbox{\begin{minipage}{0.9\textwidth}
\centering
\textbf{Exemple d'en-tête de feuille de chronométrage :}

\begin{tabular}{|l|l|}
\hline
\textbf{Matière :} & Tergal laine 45/55 \\
\textbf{Opération :} & Coulisser col \\
\textbf{Modèle :} & Monastir \\
\textbf{Date :} & 7 mai \\
\textbf{Taille :} & 48 \\
\textbf{Machine :} & Piqueuse plate 1A CF + guide B67 \\
\textbf{Opératrice :} & Sonia \\
\textbf{Points/cm :} & 3,5 \\
\textbf{RPM :} & 3840 \\
\textbf{Chronométreur :} & Victor \\
\hline
\end{tabular}
\end{minipage}}
\end{center}

\subsection{Position du chronométreur}

La position du chronométreur est importante pour la qualité des mesures et le confort de l'opérateur.

\subsubsection{Position idéale}

\begin{itemize}
    \item Se tenir \textbf{debout}, un peu en retrait de l'opératrice
    \item Se placer sur la \textbf{gauche} de l'opératrice (pour les droitiers)
    \item Avoir dans le champ visuel : le chronomètre et les mains de l'opératrice
    \item Ne pas gêner le travail ni la circulation dans l'atelier
\end{itemize}

\begin{center}
\fbox{\begin{minipage}{0.85\textwidth}
\centering
\textbf{Conseils pratiques :}
\begin{itemize}
    \item Alterner les prises de temps et le dépouillement pour réduire la fatigue
    \item Éviter de se "cacher" pour chronométrer - privilégier la transparence
    \item Adopter une posture permettant de lire le chronomètre et d'écrire sans mouvements de tête
\end{itemize}
\end{minipage}}
\end{center}

\subsection{Décomposition du travail}

La décomposition du travail en séquences élémentaires est une étape fondamentale. Elle permet d'identifier précisément les temps à mesurer.

\subsubsection{Temps cycliques}

Les temps cycliques sont ceux qui se répètent à chaque cycle de travail. Un cycle est la période qui se déroule depuis le moment où l'opérateur prend un produit jusqu'au moment où il l'évacue pour en reprendre un autre.

\begin{exemple}
Pour une opération de couture de col :
\begin{itemize}
    \item Approvisionner dessous de col
    \item Approvisionner dessus de col et déposer sur dessous
    \item Ajuster et présenter sous pied presseur
    \item Coulisser 3 côtés
    \item Dégager, contrôler, évacuer
\end{itemize}
\end{exemple}

\subsubsection{Temps fréquentiels}

Les temps fréquentiels sont ceux dont la fréquence est indépendante du cycle de travail. Ils ne se produisent pas à chaque cycle.

\begin{table}[ht]
\centering
\begin{tabular}{|l|c|}
\hline
\textbf{Activité} & \textbf{Fréquence} \\
\hline
Approvisionner paquet & 1 fois par lot \\
Délier paquet & 1 fois par lot \\
Changer la couleur des fils & À chaque changement de série \\
Remplacer un rouleau de biais & Après épuisement \\
Enregistrer la production & 1 fois par lot \\
\hline
\end{tabular}
\caption{Exemples de temps fréquentiels}
\end{table}

\subsubsection{Les "tops" de début et fin de séquence}

Pour une décomposition efficace, chaque séquence doit être identifiée par un "top" de début et un "top" de fin clairement identifiables.

\begin{center}
\fbox{\begin{minipage}{0.9\textwidth}
\textbf{Critères pour choisir les "tops" :}

\begin{itemize}
    \item Un geste ample et visible
    \item Un bruit caractéristique (démarrage de machine, coupe-fils)
    \item La pose ou la prise d'un outil (ciseaux posés sur la table)
    \item Un mouvement précis (main qui quitte la pièce)
\end{itemize}
\end{minipage}}
\end{center}

\subsubsection{Nature des temps observés}

\begin{description}[labelwidth=3.5cm]
    \item[Avant opération] Préparation, approvisionnement, mise en place
    \item[Pendant opération] Exécution de la tâche principale
    \item[Après opération] Contrôle, évacuation, rangement
\end{description}

\subsection{Définition du nombre d'observations}

Le nombre de cycles à chronométrer dépend de trois éléments :
\begin{enumerate}
    \item La variabilité des temps observés
    \item L'erreur admissible
    \item L'intervalle ou le degré de confiance souhaité
\end{enumerate}

\subsubsection{Méthode par tableau}

Pour une première approche, on peut utiliser un tableau basé sur la durée du cycle et le volume annuel de production.

\begin{table}[ht]
\centering
\begin{tabular}{|c|c|c|c|}
\hline
\multicolumn{4}{|c|}{\textbf{Nombre minimum de cycles à chronométrer}} \\
\hline
\textbf{Durée du cycle (min)} & \textbf{Activité annuelle < 1 000} & \textbf{1 000 - 10 000} & \textbf{> 10 000} \\
\hline
< 0,12 & 480 & 112 & 120 \\
0,12 - 0,25 & 240 & 60 & 80 \\
0,25 - 0,50 & 120 & 40 & 60 \\
0,50 - 0,75 & 80 & 30 & 50 \\
0,75 - 1,00 & 60 & 25 & 40 \\
1,00 - 2,00 & 40 & 20 & 30 \\
2,00 - 5,00 & 30 & 15 & 20 \\
5,00 - 10,00 & 25 & 12 & 15 \\
10,00 - 20,00 & 20 & 10 & 12 \\
20,00 - 40,00 & 15 & 8 & 10 \\
40,00 - 80,00 & 12 & 6 & 8 \\
80,00 - 120,00 & 10 & 5 & 6 \\
> 120,00 & 8 & 4 & 5 \\
\hline
\end{tabular}
\caption{Nombre de cycles en fonction de la durée et du volume annuel}
\end{table}

\subsubsection{Méthode statistique (formule)}

Pour une précision contrôlée, on utilise une approche statistique. On réalise d'abord $n$ chronométrages préliminaires (10 à 20 mesures) pour estimer la moyenne et l'écart-type.

\begin{definition}
La \textbf{formule statistique} permet de déterminer le nombre d'observations nécessaires pour atteindre un niveau de confiance et une précision donnés.
\end{definition}

\[
n = \left( \frac{t \cdot \sigma}{E \cdot \overline{X}} \right)^2
\]

Où :
\begin{itemize}
    \item $n$ : nombre d'observations nécessaires
    \item $t$ : variable aléatoire de la distribution normale (niveau de confiance)
    \item $\sigma$ : écart-type de l'échantillon préliminaire
    \item $E$ : erreur admissible (en pourcentage de la moyenne)
    \item $\overline{X}$ : moyenne de l'échantillon préliminaire
\end{itemize}

\paragraph{Valeurs typiques de t}

\begin{table}[ht]
\centering
\begin{tabular}{|c|c|}
\hline
\textbf{Intervalle de confiance (\%)} & \textbf{Variable aléatoire t} \\
\hline
90 & 1,65 \\
95 & 1,96 \\
95,5 & 2,00 \\
98 & 2,33 \\
99 & 2,58 \\
\hline
\end{tabular}
\caption{Valeurs de t selon le niveau de confiance}
\end{table}

\paragraph{Calcul de la moyenne et de l'écart-type}

\[
\overline{X} = \frac{\sum_{i=1}^{n} X_i}{n}
\]

\[
\sigma = \sqrt{\frac{\sum_{i=1}^{n} (X_i - \overline{X})^2}{n-1}}
\]

\subsubsection{Exemples d'application}

\begin{exemple}
En utilisant un chronomètre gradué en centièmes de minutes, on effectue initialement 50 lectures dont la moyenne des temps chronométrés est de 6,4 minutes avec un écart-type de 2,1 minutes. La politique de l'entreprise est d'avoir un degré de confiance de 95\%.

\textbf{Calcul avec erreur admissible de ±10\% :}

\[
n = \left( \frac{1,96 \times 2,1}{0,10 \times 6,4} \right)^2 = \left( \frac{4,116}{0,64} \right)^2 = (6,43)^2 \approx 42
\]

\textbf{Calcul avec erreur admissible de ±0,5 minute :}

\[
n = \left( \frac{1,96 \times 2,1}{0,5} \right)^2 = \left( \frac{4,116}{0,5} \right)^2 = (8,232)^2 \approx 68
\]
\end{exemple}

\begin{exemple}
On réalise 10 chronométrages préliminaires sur une séquence de travail :

\begin{center}
\begin{tabular}{|c|c|c|c|c|c|c|c|c|c|c|}
\hline
N° & 1 & 2 & 3 & 4 & 5 & 6 & 7 & 8 & 9 & 10 \\
\hline
Temps (cmin) & 134 & 132 & 130 & 135 & 138 & 133 & 131 & 129 & 136 & 132 \\
\hline
\end{tabular}
\end{center}

Calcul de la moyenne :
\[
\overline{X} = \frac{134 + 132 + 130 + 135 + 138 + 133 + 131 + 129 + 136 + 132}{10} = \frac{1330}{10} = 133 \text{ cmin}
\]

Calcul de l'écart-type :
\[
\sigma = \sqrt{\frac{(134-133)^2 + (132-133)^2 + ... + (132-133)^2}{10-1}} \approx 2,94 \text{ cmin}
\]

Pour un niveau de confiance de 95\% ($t = 1,96$) et une erreur admissible de ±5\% ($E = 0,05$) :

\[
n = \left( \frac{1,96 \times 2,94}{0,05 \times 133} \right)^2 = \left( \frac{5,7624}{6,65} \right)^2 = (0,866)^2 \approx 1
\]

Ce résultat indique que 10 mesures sont déjà suffisantes. Si n > 10, il faudrait compléter avec des mesures supplémentaires.
\end{exemple}

\subsubsection{Procédure pratique}

\begin{enumerate}
    \item Réaliser $n$ chronométrages préliminaires (10 à 20 mesures)
    \item Calculer la moyenne $\overline{X}$ et l'écart-type $\sigma$
    \item Déterminer $n$ avec la formule pour le niveau de confiance et l'erreur souhaités
    \item Si $n$ calculé > mesures initiales, compléter avec des mesures supplémentaires
    \item Vérifier la précision atteinte (erreur relative)
\end{enumerate}

\subsection{Synthèse}

\begin{center}
\fbox{\begin{minipage}{0.9\textwidth}
\textbf{Points clés à retenir :}

\begin{itemize}
    \item Le chronomètre doit être adapté à l'unité de temps choisie
    \item La feuille de relevés doit contenir toutes les informations nécessaires
    \item Le chronométreur doit se placer à gauche de l'opérateur, debout
    \item La décomposition du travail doit utiliser des "tops" clairement identifiables
    \item Le nombre d'observations se détermine par tableau ou par formule statistique
    \item Pour une étude fiable, le nombre d'observations doit être suffisant (10-30 selon le type)
\end{itemize}
\end{minipage}}
\end{center}

% ==================== FIN DE LA PARTIE 4 ====================
% ==================== PARTIE 5 : RELEVÉ DES TEMPS ====================

\section{Relevé des Temps}

Le relevé des temps est l'étape centrale du chronométrage. Elle consiste à mesurer effectivement les durées des différentes séquences de travail selon la méthode préparée. La qualité des relevés conditionne directement la fiabilité des temps obtenus.

\subsection{Procédure de chronométrage}

\subsubsection{Préparation avant le relevé}

Avant de commencer les mesures, le chronométreur doit :
\begin{itemize}
    \item Vérifier que toutes les conditions préalables sont réunies (poste stabilisé, opérateur formé)
    \item S'assurer que la feuille de relevés est correctement préparée
    \item Choisir l'unité de temps et régler le chronomètre en conséquence
    \item Se positionner correctement par rapport à l'opérateur
    \item Expliquer brièvement à l'opérateur le déroulement de l'étude
\end{itemize}

\subsubsection{Chronométrage des temps cycliques}

Les temps cycliques sont mesurés de manière continue sur plusieurs cycles successifs.

\begin{definition}
Le \textbf{chronométrage continu} consiste à mesurer les temps sans arrêter le chronomètre entre les cycles. On lit et enregistre le temps à la fin de chaque séquence, puis on soustrait les valeurs pour obtenir les durées.
\end{definition}

\textbf{Procédure pour le chronométrage continu :}

\begin{enumerate}
    \item Démarrer le chronomètre au début du premier cycle
    \item À la fin de chaque séquence, lire le temps cumulé et l'enregistrer
    \item Continuer sans arrêter le chronomètre pour les cycles suivants
    \item En fin de série, arrêter le chronomètre et noter le temps total
\end{enumerate}

\begin{table}[ht]
\centering
\begin{tabular}{|l|c|c|c|c|c|}
\hline
\textbf{Séquence} & \multicolumn{5}{|c|}{\textbf{Temps cumulés (cmin)}} \\
\cline{2-6}
 & Cycle 1 & Cycle 2 & Cycle 3 & Cycle 4 & Cycle 5 \\
\hline
Prendre caisse & 13 & 46 & 78 & 111 & 143 \\
Enfiler & 21 & 54 & 86 & 119 & 151 \\
Transférer & 35 & 68 & 100 & 133 & 165 \\
Peser & 41 & 74 & 106 & 139 & 171 \\
Palettiser & 54 & 87 & 119 & 152 & 184 \\
Retour & 62 & 95 & 127 & 160 & 192 \\
\hline
\end{tabular}
\caption{Exemple de relevé en chronométrage continu}
\end{table}

\textbf{Calcul des durées :}
\begin{itemize}
    \item Durée d'une séquence = temps cumulé fin - temps cumulé début
    \item Durée du cycle = temps cumulé fin de cycle - temps cumulé début de cycle
\end{itemize}

\subsubsection{Chronométrage des temps fréquentiels}

Les temps fréquentiels sont mesurés lorsqu'ils se produisent, ce qui peut être irrégulier.

\textbf{Procédure pour les temps fréquentiels :}
\begin{itemize}
    \item Les identifier lorsqu'ils surviennent pendant l'observation
    \item Mesurer leur durée comme pour un temps cyclique
    \item Noter la fréquence d'apparition (nombre de fois par lot, par jour, etc.)
    \item Ne pas les intégrer dans les temps cycliques
\end{itemize}

\begin{exemple}
Lors d'une séance de chronométrage sur une piqueuse, on observe :
\begin{itemize}
    \item Changement de canette : survient toutes les 45 minutes
    \item Casse de fil : survient de manière aléatoire (3 fois en 2 heures)
    \item Réglage de tension : 1 fois par changement de série
\end{itemize}
Chacun de ces temps doit être mesuré séparément et enregistré avec sa fréquence.
\end{exemple}

\subsection{Gestion des irrégularités et valeurs aberrantes}

\subsubsection{Définition}

Une irrégularité est un événement anormal qui perturbe le déroulement normal du travail. Une valeur aberrante est une mesure qui s'écarte significativement des autres observations.

\subsubsection{Types d'irrégularités}

\begin{table}[ht]
\centering
\begin{tabular}{|p{5cm}|p{9cm}|}
\hline
\textbf{Type d'irrégularité} & \textbf{Exemples} \\
\hline
Liées à l'opérateur & Hésitation, geste maladroit, arrêt pour boire \\
Liées à la machine & Casse d'aiguille, casse de fil, bourrage \\
Liées à la matière & Défaut de coupe, crans manquants \\
Liées à l'environnement & Demande d'un supérieur, appel téléphonique \\
\hline
\end{tabular}
\caption{Principaux types d'irrégularités}
\end{table}

\subsubsection{Traitement des irrégularités}

\begin{enumerate}
    \item \textbf{Identifier} la cause de l'irrégularité
    \item \textbf{Noter} l'événement dans la colonne "Observations" de la feuille de relevés
    \item \textbf{Entourer} la valeur aberrante sur la feuille de relevés
    \item \textbf{Ne pas éliminer} définitivement avant analyse
\end{enumerate}

\begin{center}
\fbox{\begin{minipage}{0.85\textwidth}
\centering
\textbf{Règle importante :} \\
Ne pas effacer les valeurs aberrantes. Les entourer et les noter comme irrégularités. \\
Elles seront traitées lors du dépouillement (élimination ou conservation selon l'analyse).
\end{minipage}}
\end{center}

\subsubsection{Détection des valeurs aberrantes}

\textbf{Méthode visuelle :} Après quelques mesures, le chronométreur a une idée de la valeur moyenne. Les valeurs très éloignées sont suspectes.

\textbf{Méthode statistique :} Utilisation des cartes de contrôle avec limites à ±2 ou ±3 écarts-types.

\begin{exemple}
Série de temps pour une séquence : 40, 42, 39, 41, 38, 43, 57, 42, 43, 44, 42, 38, 39, 26, 42, 40

\begin{itemize}
    \item Valeur 57 : très supérieure à la moyenne (~40) → aberrant
    \item Valeur 26 : très inférieure à la moyenne → aberrant
\end{itemize}
On entoure ces deux valeurs pour les traiter ultérieurement.
\end{exemple}

\subsubsection{Taux d'aléas}

Le taux d'aléas permet d'évaluer la stabilité du poste :

\[
\text{Taux d'aléas} = \frac{\text{Nombre de valeurs aberrantes}}{\text{Nombre total de mesures}} \times 100
\]

\begin{table}[ht]
\centering
\begin{tabular}{|c|c|}
\hline
\textbf{Taux d'aléas} & \textbf{Interprétation} \\
\hline
< 10\% & Poste très stable \\
10 - 30\% & Poste stable \\
30 - 50\% & Poste moyennement stable \\
> 50\% & Poste instable - amélioration nécessaire \\
\hline
\end{tabular}
\caption{Interprétation du taux d'aléas}
\end{table}

\subsection{Le bouclage (contrôle de précision)}

\subsubsection{Définition}

Le bouclage est une méthode de contrôle de la précision des mesures qui compare le temps total mesuré avec le temps réel écoulé pendant la séance de chronométrage.

\begin{definition}
Le \textbf{bouclage} consiste à mesurer simultanément avec deux chronomètres : l'un pour les relevés séquentiels, l'autre pour le temps total écoulé. L'écart entre les deux indique la qualité du chronométrage.
\end{definition}

\subsubsection{Méthode du bouclage}

\textbf{Matériel nécessaire :}
\begin{itemize}
    \item Un chronomètre pour les relevés séquentiels (mesures des séquences)
    \item Un chronomètre pour le temps de bouclage (temps réel total)
\end{itemize}

\textbf{Procédure :}
\begin{enumerate}
    \item Démarrer les deux chronomètres simultanément au début de la séance
    \item Utiliser le premier chronomètre pour mesurer chaque séquence
    \item Ne pas arrêter le second chronomètre pendant toute la séance
    \item Arrêter les deux chronomètres à la fin de la dernière mesure
    \item Noter le temps de bouclage (temps total réel)
\end{enumerate}

\subsubsection{Calcul de l'écart de bouclage}

\[
\text{Écart} = \frac{\text{Temps total mesuré} - \text{Temps de bouclage}}{\text{Temps de bouclage}} \times 100
\]

Où :
\begin{itemize}
    \item \textbf{Temps total mesuré} = somme des temps relevés + temps des irrégularités
    \item \textbf{Temps de bouclage} = temps réel écoulé entre début et fin de la séance
\end{itemize}

\subsubsection{Tolérances de précision}

\begin{table}[ht]
\centering
\begin{tabular}{|l|c|}
\hline
\textbf{Type de chronométrage} & \textbf{Écart maximum toléré} \\
\hline
Chronométrage avec jugement d'allure & $\pm 5\%$ \\
Chronométrage sans jugement d'allure & $\pm 2\%$ \\
\hline
\end{tabular}
\caption{Tolérances de précision pour le bouclage}
\end{table}

\subsubsection{Exemple de bouclage}

\begin{exemple}
Une séance de chronométrage donne les résultats suivants :

\begin{itemize}
    \item Somme des temps cycliques relevés : 125,50 cmn
    \item Temps des irrégularités notées : 4,20 cmn
    \item Temps de bouclage (chronomètre secondaire) : 130,00 cmn
\end{itemize}

Calcul du temps total mesuré :
\[
\text{Temps total mesuré} = 125,50 + 4,20 = 129,70 \text{ cmn}
\]

Calcul de l'écart :
\[
\text{Écart} = \frac{129,70 - 130,00}{130,00} \times 100 = -0,23\%
\]

L'écart est inférieur à 2\%, le chronométrage est acceptable.
\end{exemple}

\subsubsection{Interprétation des écarts}

\begin{itemize}
    \item \textbf{Écart positif} : le chronométreur a mesuré plus de temps que le temps réel
    \begin{itemize}
        \item Cause possible : délais entre les mesures, temps morts comptés
    \end{itemize}
    \item \textbf{Écart négatif} : le chronométreur a mesuré moins de temps que le temps réel
    \begin{itemize}
        \item Cause possible : séquences oubliées, démarrage tardif du chronomètre
    \end{itemize}
\end{itemize}

\begin{center}
\fbox{\begin{minipage}{0.85\textwidth}
\centering
\textbf{En cas de dépassement des tolérances :}
\begin{itemize}
    \item Vérifier les calculs de somme
    \item Revoir les valeurs entourées (irrégularités)
    \item Si l'écart persiste, refaire la série de mesures
\end{itemize}
\end{minipage}}
\end{center}

\subsection{Enregistrement des résultats}

\subsubsection{Feuille de relevés type}

\begin{table}[ht]
\centering
\begin{tabular}{|l|c|c|c|c|c|c|c|}
\hline
\textbf{Séquence} & \multicolumn{6}{|c|}{\textbf{Temps chronométrés (cmin)}} & \textbf{Observations} \\
\cline{2-7}
 & 1 & 2 & 3 & 4 & 5 & ... & \\
\hline
Prendre pièce & 13 & 14 & 11 & 13 & 14 & ... & \\
\hline
Positionner & 12 & 11 & 13 & 12 & 11 & ... & \\
\hline
Piquer & 45 & 46 & 44 & 45 & 46 & ... & Casse fil cycle 3 \\
\hline
Évacuer & 10 & 11 & 10 & 11 & 10 & ... & \\
\hline
\end{tabular}
\caption{Exemple de feuille de relevés avec observations}
\end{table}

\subsection{Synthèse}

\begin{center}
\fbox{\begin{minipage}{0.9\textwidth}
\textbf{Points clés à retenir :}

\begin{itemize}
    \item Utiliser le chronométrage continu pour les temps cycliques
    \item Mesurer les temps fréquentiels séparément avec leur fréquence
    \item Entourer et noter les valeurs aberrantes sans les effacer
    \item Le bouclage permet de contrôler la précision des mesures
    \item L'écart de bouclage ne doit pas dépasser 2\% (sans JA) ou 5\% (avec JA)
    \item Calculer le taux d'aléas pour évaluer la stabilité du poste
\end{itemize}
\end{minipage}}
\end{center}

% ==================== FIN DE LA PARTIE 5 ====================
% ==================== PARTIE 3 : TRAITEMENT ET ANALYSE DES DONNÉES ====================
\part*{Partie 3 : Traitement et Analyse des Données}
\addcontentsline{toc}{part}{Partie 3 : Traitement et Analyse des Données}

\section{Le Jugement d'Allure (JA)}

Le jugement d'allure est une étape fondamentale du chronométrage avec jugement d'allure. Il permet de corriger les temps observés pour les ramener à une allure de référence, éliminant ainsi les variations dues aux différences individuelles entre opérateurs.

\subsection{Définition et objectif}

\subsubsection{Définition}

\begin{definition}
Le \textbf{jugement d'allure} est l'action par laquelle un observateur entraîné définit l'allure (le rythme) d'un opérateur par rapport à la représentation mentale qu'il a de celle d'un exécutant placé dans des conditions similaires, l'allure de ce dernier étant prise comme allure de référence (allure 100).
\end{definition}

\subsubsection{Objectifs}

Le jugement d'allure poursuit plusieurs objectifs :

\begin{itemize}
    \item \textbf{Neutraliser les différences individuelles} : chaque opérateur travaille à son propre rythme
    \item \textbf{Obtenir un temps de référence universel} : le temps à l'allure 100
    \item \textbf{Assurer l'équité} : le temps alloué ne dépend pas de l'opérateur chronométré
    \item \textbf{Permettre la comparaison} : entre postes, entre ateliers, entre entreprises
\end{itemize}

\begin{center}
\fbox{\begin{minipage}{0.85\textwidth}
\centering
\textbf{Principe fondamental} \\
Le jugement d'allure corrige le temps observé pour le ramener à ce qu'aurait réalisé \\
un opérateur moyen travaillant à une allure normale (allure 100).
\end{minipage}}
\end{center}

\subsection{Échelles de jugement}

Plusieurs échelles de jugement d'allure existent. Le choix de l'échelle doit être fait par l'entreprise et utilisé de manière cohérente.

\subsubsection{Échelle BTE (Bureau des Temps Élémentaires)}

L'échelle la plus couramment utilisée en Europe est l'échelle BTE :

\begin{table}[ht]
\centering
\begin{tabular}{|c|c|}
\hline
\textbf{Allure} & \textbf{Interprétation} \\
\hline
80 & Allure minimale tolérable \\
90 & Allure lente mais acceptable \\
100 & Allure normale (référence) \\
110 & Allure rapide \\
120 & Allure maximale (opérateur très motivé et expérimenté) \\
\hline
\end{tabular}
\caption{Échelle BTE du jugement d'allure}
\end{table}

\begin{center}
\fbox{\begin{minipage}{0.85\textwidth}
\centering
\textbf{Conseil pratique :} \\
Il est recommandé d'éviter de chronométrer des opérateurs dont l'allure est inférieure à 80 ou supérieure à 120, car il devient difficile d'apprécier l'allure avec précision. L'idéal se situe entre 80 et 110.
\end{minipage}}
\end{center}

\subsubsection{Tableau de correspondance entre échelles}

Différentes échelles sont utilisées selon les entreprises et les pays. Le tableau suivant permet de convertir d'une échelle à l'autre :

\begin{table}[ht]
\centering
\begin{tabular}{|c|c|c|c|}
\hline
\textbf{Bedaux} & \textbf{Maynard} & \textbf{BTE} & \textbf{Marche (km/h)} \\
\hline
75 & 100 & 80 & 4,5 \\
80 & 107 & 85 & 5,0 \\
85 & 113 & 90 & 5,4 \\
90 & 120 & 95 & 5,8 \\
95 & 127 & 100 & 6,2 \\
100 & 133 & 105 & 6,6 \\
105 & 140 & 110 & 7,0 \\
110 & 147 & 115 & 7,4 \\
115 & 153 & 120 & 7,8 \\
\hline
\end{tabular}
\caption{Correspondance entre les principales échelles de jugement d'allure}
\end{table}

\subsection{Éléments d'appréciation}

Pour évaluer l'allure, le chronométreur doit considérer plusieurs facteurs qui influencent le rythme de travail.

\subsubsection{L'habileté}

L'habileté représente la dextérité et la précision des mouvements. Elle se manifeste par :

\begin{itemize}
    \item La coordination des mouvements
    \item La précision des gestes
    \item L'absence de mouvements parasites
    \item La maîtrise du geste
\end{itemize}

\subsubsection{L'activité (vitesse)}

L'activité correspond à la vitesse d'exécution des gestes. Elle se caractérise par :

\begin{itemize}
    \item La rapidité des mouvements
    \item La fluidité des enchaînements
    \item La dynamique du travail
\end{itemize}

\subsubsection{La régularité (stabilité)}

La régularité reflète la constance du rythme de travail. Elle se traduit par :

\begin{itemize}
    \item La stabilité des temps entre cycles
    \item L'absence d'hésitations
    \item La constance de l'attention
\end{itemize}

\subsubsection{Les conditions de travail}

Les conditions de travail influencent la performance. Elles incluent :

\begin{itemize}
    \item L'environnement (bruit, température, éclairage)
    \item L'ergonomie du poste
    \item La qualité des outils et équipements
\end{itemize}

\subsection{Méthodes d'évaluation}

\subsubsection{Méthode directe}

Le chronométreur attribue directement un coefficient d'allure à l'opérateur observé. Cette méthode nécessite :

\begin{itemize}
    \item Une expérience importante
    \item Un étalonnage régulier
    \item Une bonne capacité d'observation
\end{itemize}

\subsubsection{Méthode de Levelling}

La méthode de Levelling (ou d'évaluation par niveaux) permet une évaluation plus objective en décomposant l'appréciation en plusieurs facteurs.

\begin{table}[ht]
\centering
\begin{tabular}{|l|c|c|}
\hline
\textbf{Facteur} & \textbf{Critère} & \textbf{Variation possible} \\
\hline
Habileté & Précision des mouvements & $-0,15$ à $+0,15$ \\
Activité & Vitesse d'exécution & $-0,22$ à $+0,13$ \\
Conditions de travail & Environnement du poste & $-0,07$ à $+0,06$ \\
Régularité & Constance du rythme & $-0,04$ à $+0,04$ \\
\hline
\end{tabular}
\caption{Facteurs d'évaluation du jugement d'allure - Méthode de Levelling}
\end{table}

\subsubsection{Grille de notation Levelling détaillée}

\begin{table}[ht]
\centering
\begin{tabular}{|c|c|c|c|c|}
\hline
\textbf{Niveau} & \textbf{Habileté} & \textbf{Activité} & \textbf{Conditions} & \textbf{Régularité} \\
\hline
Parfaite & $+0,15$ & $+0,13$ & $+0,06$ & $+0,04$ \\
Excellente & $+0,11$ & $+0,10$ & $+0,04$ & $+0,03$ \\
Très bonne & $+0,08$ & $+0,08$ & $+0,03$ & $+0,02$ \\
Bonne & $+0,06$ & $+0,06$ & $+0,02$ & $+0,01$ \\
Moyenne & $0$ & $0$ & $0$ & $0$ \\
Passable & $-0,05$ & $-0,04$ & $-0,03$ & $-0,02$ \\
Médiocre & $-0,10$ & $-0,08$ & $-0,07$ & $-0,04$ \\
Mauvaise & $-0,15$ & $-0,12$ & $-0,10$ & $-0,06$ \\
\hline
\end{tabular}
\caption{Grille détaillée de la méthode de Levelling}
\end{table}

\begin{exemple}
\textbf{Application de la méthode de Levelling :}

Temps chronométré = 15 dmh

Correctifs évalués :
\begin{itemize}
    \item Habileté : Excellente = $+0,11$
    \item Activité : Bonne = $+0,06$
    \item Conditions de travail : Passable = $-0,03$
    \item Régularité : Bonne = $+0,01$
\end{itemize}

Total des correctifs = $0,11 + 0,06 - 0,03 + 0,01 = 0,15$

Temps corrigé à l'allure 100 :
\[
T_{100} = 15 \times (1 + 0,15) = 17,25 \text{ dmh}
\]
\end{exemple}

\subsection{Enregistrement du jugement d'allure}

\subsubsection{Moment de l'enregistrement}

Le jugement d'allure doit être enregistré :
\begin{itemize}
    \item \textbf{Pendant le chronométrage} : au fur et à mesure des observations
    \item \textbf{Pour chaque séquence} : l'allure peut varier selon les phases du travail
    \item \textbf{Dès que possible} : l'impression est plus fraîche et plus précise
\end{itemize}

\subsubsection{Enregistrement sur la feuille de relevés}

\begin{table}[ht]
\centering
\begin{tabular}{|l|c|c|c|c|c|c|}
\hline
\textbf{Séquence} & \multicolumn{5}{|c|}{\textbf{Temps chronométrés (cmin)}} & \textbf{JA} \\
\cline{2-6}
 & 1 & 2 & 3 & 4 & 5 & \\
\hline
Prendre pièce & 13 & 14 & 11 & 13 & 14 & 95 \\
Positionner & 12 & 11 & 13 & 12 & 11 & 100 \\
Piquer & 45 & 46 & 44 & 45 & 46 & 110 \\
Évacuer & 10 & 11 & 10 & 11 & 10 & 90 \\
\hline
\end{tabular}
\caption{Enregistrement des jugements d'allure par séquence}
\end{table}

\subsubsection{Variabilité du jugement d'allure}

Il est important de noter que :
\begin{itemize}
    \item L'allure peut varier d'une séquence à l'autre
    \item L'allure peut varier dans le temps (début vs fin de journée)
    \item Chaque séquence doit avoir son propre jugement d'allure
\end{itemize}

\subsection{Conversion du temps observé en temps à l'allure 100}

\subsubsection{Formule de conversion}

\[
T_{100} = T_{\text{observé}} \times \frac{JA}{100}
\]

Où :
\begin{itemize}
    \item $T_{100}$ : temps à l'allure 100 (temps de base)
    \item $T_{\text{observé}}$ : temps chronométré
    \item $JA$ : jugement d'allure (valeur sur l'échelle choisie)
\end{itemize}

\subsubsection{Exemples de conversion}

\begin{exemple}
Pour trois opérateurs différents réalisant la même tâche :

\begin{itemize}
    \item \textbf{Opérateur A} : $T_{\text{obs}} = 15$ cmin, $JA = 80$ \\
    $T_{100} = 15 \times \frac{80}{100} = 12$ cmin
    
    \item \textbf{Opérateur B} : $T_{\text{obs}} = 12$ cmin, $JA = 100$ \\
    $T_{100} = 12 \times \frac{100}{100} = 12$ cmin
    
    \item \textbf{Opérateur C} : $T_{\text{obs}} = 10$ cmin, $JA = 120$ \\
    $T_{100} = 10 \times \frac{120}{100} = 12$ cmin
\end{itemize}

Le temps à l'allure 100 est identique pour les trois opérateurs, ce qui garantit l'équité.
\end{exemple}

\subsection{Étalonnage du chronométreur}

\subsubsection{Nécessité de l'étalonnage}

Pour garantir la fiabilité des jugements d'allure, le chronométreur doit :
\begin{itemize}
    \item S'étalonner régulièrement
    \item Participer à des sessions d'étalonnage avec d'autres chronométreurs
    \item Suivre des formations spécifiques
\end{itemize}

\subsubsection{Méthode d'étalonnage}

L'étalonnage consiste à :
\begin{enumerate}
    \item Chronométrer simultanément avec plusieurs chronométreurs
    \item Comparer les jugements d'allure attribués
    \item Analyser les écarts et ajuster sa perception
    \item Répéter régulièrement l'exercice
\end{enumerate}

\begin{center}
\fbox{\begin{minipage}{0.85\textwidth}
\centering
\textbf{Recommandation :} \\
Un chronométreur doit participer à des sessions d'étalonnage \\
au moins une fois par an pour maintenir la qualité de ses jugements.
\end{minipage}}
\end{center}

\subsection{Qualités du chronométreur pour le jugement d'allure}

Un bon chronométreur doit posséder :

\begin{itemize}
    \item \textbf{Une bonne capacité d'observation} : percevoir les détails du geste
    \item \textbf{Une mémoire visuelle développée} : se souvenir de l'allure de référence
    \item \textbf{Une objectivité} : ne pas être influencé par l'apparence ou le caractère
    \item \textbf{Une expérience variée} : avoir observé de nombreux opérateurs
    \item \textbf{Une honnêteté intellectuelle} : reconnaître ses difficultés
\end{itemize}

\subsection{Synthèse}

\begin{center}
\fbox{\begin{minipage}{0.9\textwidth}
\textbf{Points clés à retenir :}

\begin{itemize}
    \item Le jugement d'allure corrige les différences individuelles entre opérateurs
    \item L'échelle BTE est la plus utilisée (100 = allure normale)
    \item La méthode de Levelling décompose l'évaluation en 4 facteurs
    \item L'allure doit être enregistrée pendant le chronométrage pour chaque séquence
    \item Le temps à l'allure 100 est obtenu par $T_{100} = T_{\text{obs}} \times JA / 100$
    \item L'étalonnage régulier est nécessaire pour maintenir la fiabilité
\end{itemize}
\end{minipage}}
\end{center}

% ==================== FIN DE LA PARTIE 6 ====================
% ==================== PARTIE 7 : LE DÉPOUILLEMENT DES MESURES ====================

\section{Le Dépouillement des Mesures}

Le dépouillement est l'étape qui consiste à transformer les données brutes du chronométrage en temps standard utilisable pour la gestion de production. Cette transformation s'effectue en plusieurs étapes : calcul du temps observé moyen, correction par le jugement d'allure, application des majorations et prise en compte de la fréquence.

\subsection{Calcul du temps observé moyen}

\subsubsection{Définition}

Le temps observé moyen est la moyenne arithmétique des temps mesurés pour une séquence donnée, après élimination des valeurs aberrantes.

\[
\overline{T}_o = \frac{\sum_{i=1}^{n} T_i}{n}
\]

Où :
\begin{itemize}
    \item $\overline{T}_o$ : temps observé moyen
    \item $T_i$ : temps mesurés pour la séquence
    \item $n$ : nombre de mesures valides
\end{itemize}

\subsubsection{Procédure de calcul}

\begin{enumerate}
    \item Additionner tous les temps relevés et conservés pour une colonne
    \item Compter le nombre de relevés acceptés (valeurs aberrantes éliminées)
    \item Diviser le total par le nombre de lectures
    \item Vérifier visuellement que la moyenne n'est pas aberrante
\end{enumerate}

\begin{exemple}
Pour une séquence "Prendre la pièce", on relève les temps suivants (en cmin) :

\begin{center}
\begin{tabular}{|c|c|c|c|c|c|c|c|c|c|}
\hline
13 & 14 & 11 & 13 & 14 & 12 & 13 & 35 & 14 & 13 \\
\hline
\end{tabular}
\end{center}

La valeur 35 est aberrante (casse de fil). On l'élimine.

Somme des temps valides = $13 + 14 + 11 + 13 + 14 + 12 + 13 + 14 + 13 = 117$

Nombre de mesures valides = 9

\[
\overline{T}_o = \frac{117}{9} = 13 \text{ cmin}
\]
\end{exemple}

\subsection{Correction par le JA pour obtenir le temps de base}

\subsubsection{Définition du temps de base}

Le temps de base (ou temps à l'allure 100) est obtenu en corrigeant le temps observé moyen par le jugement d'allure. Il représente le temps qu'aurait réalisé un opérateur moyen travaillant à une allure normale.

\[
T_{100} = \overline{T}_o \times \frac{JA}{100}
\]

\subsubsection{Calcul pour chaque séquence}

Pour chaque séquence, on applique le jugement d'allure qui lui est propre.

\begin{exemple}
D'après l'exemple précédent :

\begin{itemize}
    \item $\overline{T}_o = 13$ cmin
    \item $JA = 95$ (opérateur travaille légèrement moins vite que la normale)
\end{itemize}

\[
T_{100} = 13 \times \frac{95}{100} = 12,35 \text{ cmin}
\]
\end{exemple}

\subsubsection{Cas d'allures différentes par séquence}

\begin{exemple}
Pour une opération de couture de col :

\begin{table}[ht]
\centering
\begin{tabular}{|l|c|c|c|}
\hline
\textbf{Séquence} & $\overline{T}_o$ (cmin) & \textbf{JA} & $T_{100}$ (cmin) \\
\hline
Approvisionner & 29,38 & 105 & 30,85 \\
Piquer & 86,38 & 105 & 90,70 \\
Évacuer & 10,50 & 105 & 11,03 \\
\hline
\end{tabular}
\caption{Calcul du temps de base pour chaque séquence}
\end{table}
\end{exemple}

\subsection{Application des majorations (temps standard)}

\subsubsection{Définition du temps standard}

Le temps standard est le temps de base majoré pour prendre en compte les temps improductifs inévitables (fatigue, besoins personnels, aléas).

\[
T_S = T_{100} \times \text{Facteur de majoration}
\]

\subsubsection{Types de majorations}

\begin{table}[ht]
\centering
\begin{tabular}{|p{4cm}|p{10cm}|}
\hline
\textbf{Type de majoration} & \textbf{Description} \\
\hline
Temps physiologiques & Besoins personnels (toilettes, boisson, collation) \\
Fatigue & Récupération nécessaire après un effort \\
Temps machine & Temps pendant lequel la machine travaille seule \\
Aléas & Événements imprévus (casse de fil, réglages) \\
\hline
\end{tabular}
\caption{Principaux types de majorations}
\end{table}

\subsubsection{Formules de calcul}

\textbf{Majoration exprimée en pourcentage du temps de base :}

\[
T_S = T_{100} \times (1 + M)
\]

Où $M$ est le taux de majoration (exemple : $M = 0,20$ pour 20\%)

\textbf{Majoration exprimée en pourcentage du temps travaillé :}

\[
T_S = T_{100} \times \frac{1}{1 - M}
\]

\begin{exemple}
Pour un temps de base de 1 minute :

\begin{itemize}
    \item Majoration de 20\% du temps de base : $T_S = 1 \times (1 + 0,20) = 1,20$ min
    \item Majoration de 20\% du temps travaillé : $T_S = 1 \times \frac{1}{1 - 0,20} = 1,25$ min
\end{itemize}
\end{exemple}

\subsection{Utilisation de tables de coefficients}

\subsubsection{Méthode traditionnelle (par nature de temps)}

Dans la méthode traditionnelle, on applique des coefficients différents selon la nature du temps (manuel, machine, technico-manuel).

\begin{table}[ht]
\centering
\begin{tabular}{|l|c|}
\hline
\textbf{Nature du temps} & \textbf{Coefficient majorateur} \\
\hline
Temps manuel (Tm) & 1,10 - 1,15 \\
Temps technologique (Tt) & 1,00 - 1,05 \\
Temps technico-manuel (Ttm) & 1,12 - 1,20 \\
\hline
\end{tabular}
\caption{Coefficients majorateurs selon la nature du temps}
\end{table}

\subsubsection{Coefficients pour travaux manuels}

\begin{table}[ht]
\centering
\begin{tabular}{|l|c|}
\hline
\textbf{Type d'opération} & \textbf{Coefficient} \\
\hline
Manipulations (saisir, tenir, déposer, évacuer) & 1,10 \\
Dégagements (dégager pince, couper fils) & 1,10 \\
Contrôles (examiner, compter, vérifier) & 1,10 \\
Préparations (plier, déplier, épingler) & 1,12 \\
Présentations (orienter, maintenir, engager) & 1,12 \\
Compléments d'opérations (dégarnir, cranter) & 1,13 \\
\hline
\end{tabular}
\caption{Coefficients pour travaux manuels}
\end{table}

\subsubsection{Coefficients pour machines à coudre}

\begin{table}[ht]
\centering
\begin{tabular}{|l|c|}
\hline
\textbf{Type de machine} & \textbf{Coefficient} \\
\hline
Piqueuse plate 1 aiguille & 1,20 \\
Piqueuse plate 2 aiguilles & 1,23 \\
Surjeteuse 1 aiguille 2 fils & 1,13 \\
Surjeteuse 2 aiguilles 4 fils & 1,17 \\
Machine à bras déporté & 1,18 \\
Machine à boutonnières & 1,12 \\
\hline
\end{tabular}
\caption{Coefficients pour machines à coudre}
\end{table}

\subsubsection{Méthode simplifiée}

La méthode simplifiée consiste à regrouper les temps par catégorie et à appliquer un coefficient unique basé sur la proportion de temps technologique.

\begin{table}[ht]
\centering
\begin{tabular}{|c|c|}
\hline
\textbf{Proportion de temps technologique} & \textbf{Coefficient} \\
\hline
0 - 25\% & 1,14 \\
26 - 50\% & 1,15 \\
51 - 75\% & 1,16 \\
76 - 100\% & 1,17 \\
\hline
\end{tabular}
\caption{Coefficients majorateurs - Méthode simplifiée}
\end{table}

\subsubsection{Exemple d'application des coefficients}

\begin{exemple}
Pour l'opération de couture de col :

\begin{table}[ht]
\centering
\begin{tabular}{|l|c|c|c|}
\hline
\textbf{Séquence} & $T_{100}$ (cmin) & \textbf{Coefficient} & $T_{\text{majoré}}$ (cmin) \\
\hline
Approvisionner (Tm) & 29,38 & 1,12 & 32,91 \\
Piquer (Ttm) & 90,70 & 1,20 & 108,84 \\
Évacuer (Tm) & 10,50 & 1,10 & 11,55 \\
\hline
\end{tabular}
\caption{Application des coefficients majorateurs}
\end{table}
\end{exemple}

\subsection{Application de la fréquence}

\subsubsection{Définition}

La fréquence est le nombre de fois qu'une séquence se répète pour fabriquer une unité de produit fini.

\begin{table}[ht]
\centering
\begin{tabular}{|c|c|}
\hline
\textbf{Séquence mesurée} & \textbf{Fréquence par produit} \\
\hline
1 col & 1 \\
1 poignet & 2 \\
1 passant (pantalon à 6 passants) & 6 \\
1 poche & 2 \\
\hline
\end{tabular}
\caption{Exemples de fréquences en confection}
\end{table}

\subsubsection{Calcul du temps alloué avec fréquence}

\[
T_{\text{alloué}} = T_{\text{majoré}} \times \text{Fréquence}
\]

\subsubsection{Exemple d'application de la fréquence}

\begin{exemple}
Pour un pantalon nécessitant :

\begin{itemize}
    \item Pose de 2 poches arrière
    \item Pose de 6 passants
\end{itemize}

\begin{table}[ht]
\centering
\begin{tabular}{|l|c|c|c|}
\hline
\textbf{Séquence} & $T_{\text{majoré}}$ (cmin) & \textbf{Fréquence} & $T_{\text{alloué}}$ (cmin) \\
\hline
Pose poche & 45,00 & 2 & 90,00 \\
Pose passant & 8,50 & 6 & 51,00 \\
\hline
\end{tabular}
\caption{Application de la fréquence}
\end{table}
\end{exemple}

\subsection{Calcul complet du temps standard}

\subsubsection{Étapes de calcul}

\begin{enumerate}
    \item Calculer le temps observé moyen ($\overline{T}_o$) pour chaque séquence
    \item Corriger par le JA pour obtenir le temps de base ($T_{100}$)
    \item Appliquer les coefficients majorateurs ($T_{\text{majoré}}$)
    \item Appliquer la fréquence ($T_{\text{alloué}}$)
    \item Somme des temps alloués de toutes les séquences = temps standard total
\end{enumerate}

\subsubsection{Exemple complet}

\begin{exemple}
Dépouillement complet pour l'opération "Coulisser un col" :

\begin{table}[ht]
\centering
\begin{tabular}{|l|c|c|c|c|c|}
\hline
\textbf{Séquence} & $\overline{T}_o$ & \textbf{JA} & $T_{100}$ & \textbf{Coeff} & $T_{\text{maj}}$ \\
\hline
Approvisionner & 29,38 & 105 & 30,85 & 1,12 & 34,55 \\
Piquer & 86,38 & 105 & 90,70 & 1,20 & 108,84 \\
Évacuer & 10,50 & 105 & 11,03 & 1,10 & 12,13 \\
\hline
\multicolumn{5}{|c|}{\textbf{Temps alloué pour 1 col}} & 155,52 cmin \\
\hline
\end{tabular}
\caption{Dépouillement complet pour l'opération "Coulisser un col"}
\end{table}
\end{exemple}

\subsection{Intégration des temps fréquentiels}

Les temps fréquentiels doivent être ajoutés après avoir été :
\begin{itemize}
    \item Mesurés avec leur propre jugement d'allure
    \item Majorés avec les coefficients appropriés
    \item Rammenés à l'unité de produit selon leur fréquence
\end{itemize}

\begin{exemple}
\begin{itemize}
    \item Changement de canette : 30 cmin à l'allure 100, majoration 1,10
    \item Fréquence : 1 changement pour 50 produits
\end{itemize}

Temps alloué par produit = $30 \times 1,10 \div 50 = 0,66$ cmin
\end{exemple}

\subsection{Synthèse}

\begin{center}
\fbox{\begin{minipage}{0.9\textwidth}
\textbf{Points clés à retenir :}

\begin{itemize}
    \item Le temps observé moyen est la moyenne après élimination des valeurs aberrantes
    \item Le temps de base ($T_{100}$) est obtenu par $T_{100} = \overline{T}_o \times JA/100$
    \item Le temps standard intègre les majorations (fatigue, besoins, aléas)
    \item Les coefficients majorateurs varient selon la nature du temps (manuel, machine)
    \item La fréquence permet de ramener les temps à l'unité de produit
    \item Le temps alloué total est la somme des temps alloués de chaque séquence
\end{itemize}
\end{minipage}}
\end{center}

% ==================== FIN DE LA PARTIE 7 ====================
% ==================== PARTIE 4 : APPLICATIONS ET LIMITES ====================
\part*{Partie 4 : Applications et Limites}
\addcontentsline{toc}{part}{Partie 4 : Applications et Limites}

\section{Exploitation des Résultats}

Les temps déterminés par chronométrage et dépouillement sont utilisés dans de nombreuses applications de gestion de production. Leur exploitation permet d'optimiser les ressources, de planifier la production et de maîtriser les coûts.

\subsection{Détermination du temps de cycle et de la production horaire}

\subsubsection{Définition du temps de cycle}

Le temps de cycle est le temps nécessaire pour réaliser une unité de produit sur un poste de travail donné. Il correspond à la somme des temps alloués pour toutes les séquences nécessaires à la fabrication d'une unité.

\[
T_{\text{cycle}} = \sum_{i=1}^{n} T_{\text{alloué},i}
\]

Où $T_{\text{alloué},i}$ est le temps alloué pour chaque séquence après correction et majoration.

\subsubsection{Calcul de la production horaire}

La production horaire théorique est le nombre d'unités qu'un poste peut produire en une heure.

\[
P_{\text{horaire}} = \frac{\text{Temps de base horaire}}{T_{\text{cycle}}}
\]

Avec temps de base horaire = 60 minutes = 3 600 secondes = 6 000 centiminutes selon l'unité choisie.

\begin{exemple}
Pour une opération de pressage de pantalon avec un temps de cycle de 135 dmh :

\[
P_{\text{horaire}} = \frac{10\,000 \text{ dmh}}{135 \text{ dmh}} = 74 \text{ pantalons/heure}
\]
\end{exemple}

\subsubsection{Production horaire avec plusieurs machines}

Lorsqu'un opérateur dessert plusieurs machines, le temps de cycle et la production horaire sont modifiés.

\begin{exemple}
Configuration avec 2 machines pour l'opération de pressage :

\begin{itemize}
    \item Temps de cycle avec 1 machine : 135 dmh
    \item Temps de déplacement entre machines : 5 dmh
    \item Nouveau temps de cycle : 165 dmh
\end{itemize}

\[
P_{\text{horaire}} = \frac{10\,000}{165} \times 2 = 120 \text{ pantalons/heure}
\]

\textbf{Gain :} $120 - 74 = 46$ pantalons/heure (+62\%)
\end{exemple}

\subsection{Calcul des besoins en personnel et des coûts}

\subsubsection{Calcul des besoins en effectifs}

Le nombre d'opérateurs nécessaires pour atteindre une production donnée se calcule par :

\[
N_{\text{opérateurs}} = \frac{T_{\text{cycle}} \times Q_{\text{objectif}}}{T_{\text{disponible}}}
\]

Où :
\begin{itemize}
    \item $T_{\text{cycle}}$ : temps de cycle en heures par unité
    \item $Q_{\text{objectif}}$ : production à réaliser
    \item $T_{\text{disponible}}$ : temps disponible par opérateur
\end{itemize}

\begin{exemple}
Une entreprise doit produire 500 pièces par jour. Le temps de cycle est de 2,5 minutes. Le temps de travail journalier est de 7 heures (420 minutes).

\[
N = \frac{2,5 \times 500}{420} = \frac{1250}{420} \approx 2,98 \text{ opérateurs}
\]

Soit 3 opérateurs nécessaires.
\end{exemple}

\subsubsection{Calcul des coûts de revient}

Le coût de revient d'un produit peut être calculé à partir du temps standard :

\[
C_{\text{revient}} = T_{\text{standard}} \times C_{\text{minute}} + C_{\text{matières}} + C_{\text{frais}}
\]

Où $C_{\text{minute}}$ est le coût d'une minute de travail dans l'entreprise.

\begin{exemple}
\begin{itemize}
    \item Temps standard : 12,5 minutes
    \item Coût minute : 0,50 €
    \item Coût matières : 8,20 €
    \item Frais généraux : 2,30 €
\end{itemize}

\[
C_{\text{revient}} = (12,5 \times 0,50) + 8,20 + 2,30 = 6,25 + 10,50 = 16,75\ €
\]
\end{exemple}

\subsubsection{Calcul de la performance}

La performance d'un opérateur ou d'une équipe se calcule par :

\[
\text{Performance} = \frac{\text{Temps standard produit}}{\text{Temps de présence}} \times 100
\]

\begin{exemple}
Un opérateur a produit 200 pièces en 7 heures (420 minutes). Le temps standard par pièce est de 1,8 minute.

\[
\text{Temps standard produit} = 200 \times 1,8 = 360 \text{ minutes}
\]

\[
\text{Performance} = \frac{360}{420} \times 100 = 85,7\%
\]
\end{exemple}

\subsection{Équilibrage des lignes}

\subsubsection{Définition}

L'équilibrage des lignes consiste à répartir la charge de travail entre les différents postes d'une ligne de production pour minimiser les temps d'attente et maximiser la production.

\subsubsection{Objectifs de l'équilibrage}

\begin{itemize}
    \item Réduire les temps d'attente des opérateurs
    \item Diminuer les stocks en cours
    \item Augmenter la productivité globale
    \item Améliorer la flexibilité
    \item Faciliter la gestion des effectifs
\end{itemize}

\subsubsection{Méthode d'équilibrage}

\begin{enumerate}
    \item Déterminer le temps de cycle de la ligne (poste le plus chargé)
    \item Calculer le temps de cycle théorique souhaité
    \item Regrouper les opérations par poste
    \item Vérifier la charge de chaque poste
    \item Ajuster par regroupement ou décomposition
\end{enumerate}

\begin{exemple}
Une ligne de montage comporte 4 postes avec les temps suivants :

\begin{table}[ht]
\centering
\begin{tabular}{|c|c|}
\hline
\textbf{Poste} & \textbf{Temps (minutes)} \\
\hline
1 & 2,5 \\
2 & 3,0 \\
3 & 2,8 \\
4 & 2,2 \\
\hline
\end{tabular}
\caption{Temps par poste avant équilibrage}
\end{table}

Le temps de cycle est de 3,0 minutes (poste 2). L'inefficacité est de :

\[
\text{Inefficacité} = \frac{\sum \text{Temps postes} - \text{Cycle} \times \text{Nb postes}}{\text{Cycle} \times \text{Nb postes}} \times 100
\]

\[
\text{Inefficacité} = \frac{(2,5+3,0+2,8+2,2) - 3,0 \times 4}{3,0 \times 4} \times 100 = \frac{10,5 - 12}{12} \times 100 = -12,5\%
\]

Après rééquilibrage :

\begin{table}[ht]
\centering
\begin{tabular}{|c|c|}
\hline
\textbf{Poste} & \textbf{Temps (minutes)} \\
\hline
1 & 2,7 \\
2 & 2,8 \\
3 & 2,7 \\
4 & 2,8 \\
\hline
\end{tabular}
\caption{Temps par poste après équilibrage}
\end{table}

Le nouveau temps de cycle est de 2,8 minutes.
\end{exemple}

\subsection{Utilisation du simogramme pour l'équilibrage}

\subsubsection{Définition du simogramme}

Le simogramme est une représentation graphique qui permet de visualiser l'occupation des ressources (opérateurs, machines) au cours d'un cycle de production.

\subsubsection{Objectifs du simogramme}

\begin{itemize}
    \item Visualiser un cycle de production par type de temps (manuel, technologique)
    \item Étudier la mise en parallèle de ressources
    \item Analyser différentes solutions d'organisation
    \item Identifier les opportunités d'optimisation
\end{itemize}

\subsubsection{Étapes de construction d'un simogramme}

\begin{enumerate}
    \item Identifier les séquences ou opérations
    \item Recueillir les durées des opérations
    \item Identifier les ressources mobilisées
    \item Repérer les temps manuels (Tm), technologiques (Tt) et technico-manuels (Ttm)
    \item Placer les ressources en ordonnée
    \item Placer l'échelle des temps en abscisse
\end{enumerate}

\subsubsection{Exemple : Atelier de mise en plis de pantalon}

\begin{table}[ht]
\centering
\begin{tabular}{|l|c|}
\hline
\textbf{Opération} & \textbf{Durée (dmh)} \\
\hline
Prendre et poser pantalon sur la presse & 5 (Tm) \\
Positionner le pantalon sur la presse & 12 (Tm) \\
Verrouiller plateau, actionner pédale vapeur & 8 (Ttm) \\
Mise en forme automatique & 60 (Tt) \\
Actionner pédale de séchage & 5 (Ttm) \\
Séchage automatique & 30 (Tt) \\
Déverrouiller plateau & 5 (Ttm) \\
Évacuer le pantalon & 10 (Tm) \\
\hline
\end{tabular}
\caption{Décomposition des temps pour un cycle de pressage}
\end{table}

\subsubsection{Analyse du simogramme pour l'optimisation}

Le simogramme permet d'identifier les temps masqués (temps pendant lesquels l'opérateur pourrait effectuer d'autres tâches).

\begin{center}
\begin{minipage}{0.9\textwidth}
\begin{verbatim}
Ressource: Opérateur
|-----Tm-----|-----Tm-----|--Ttm--|     |--Ttm--|--Ttm--|-----Tm-----|
|                               |-----Tt-----|-----Tt-----|
0                            60                            120   135 (dmh)

Ressource: Machine
|          |          |--Ttm--|-----Tt-----|--Ttm--|-----Tt-----|--Ttm--|
0                            60                            120   135 (dmh)
\end{verbatim}
\caption{Représentation simplifiée d'un simogramme}
\end{minipage}
\end{center}

\subsubsection{Optimisation par ajout de machines}

L'analyse du simogramme peut conduire à des optimisations comme l'augmentation du nombre de machines confiées à un opérateur.

\begin{table}[ht]
\centering
\begin{tabular}{|l|c|}
\hline
\textbf{Configuration} & \textbf{Production horaire} \\
\hline
1 machine & 74 pantalons/heure \\
2 machines & 120 pantalons/heure \\
\hline
\textbf{Gain} & \textbf{+62\%} \\
\hline
\end{tabular}
\caption{Optimisation par ajout de machines}
\end{table}

\subsubsection{Calcul de l'efficacité de la ligne}

L'efficacité d'une ligne après équilibrage se calcule par :

\[
\text{Efficacité} = \frac{\sum \text{Temps standards}}{\text{Cycle} \times \text{Nombre de postes}} \times 100
\]

\begin{exemple}
Pour une ligne équilibrée avec :
\begin{itemize}
    \item Somme des temps standards : 11,0 minutes
    \item Cycle de ligne : 2,8 minutes
    \item Nombre de postes : 4
\end{itemize}

\[
\text{Efficacité} = \frac{11,0}{2,8 \times 4} \times 100 = \frac{11,0}{11,2} \times 100 = 98,2\%
\]
\end{exemple}

\subsection{Indicateurs de performance clés}

\subsubsection{Taux de rendement synthétique (TRS)}

Le TRS mesure l'efficacité globale d'un équipement :

\[
\text{TRS} = \text{Disponibilité} \times \text{Performance} \times \text{Qualité}
\]

\subsubsection{Temps de traversée (Lead Time)}

Le temps de traversée est le temps total entre le début et la fin d'un processus, incluant les temps d'attente.

\[
LT = \sum \text{Temps de transformation} + \sum \text{Temps d'attente} + \sum \text{Temps de transfert}
\]

\subsubsection{Taux de valeur ajoutée}

\[
\text{Taux VA} = \frac{\sum \text{Temps à valeur ajoutée}}{LT} \times 100
\]

\begin{exemple}
\begin{itemize}
    \item Temps de transformation : 15 minutes
    \item Temps d'attente : 45 minutes
    \item Temps de transfert : 20 minutes
    \item Lead Time total : 80 minutes
\end{itemize}

\[
\text{Taux VA} = \frac{15}{80} \times 100 = 18,75\%
\]
\end{exemple}

\subsection{Synthèse}

\begin{center}
\fbox{\begin{minipage}{0.9\textwidth}
\textbf{Points clés à retenir :}

\begin{itemize}
    \item La production horaire se calcule à partir du temps de cycle
    \item Les besoins en personnel dépendent du temps de cycle et de l'objectif de production
    \item Le coût de revient intègre le temps standard et le coût minute
    \item L'équilibrage des lignes réduit les temps d'attente et augmente la productivité
    \item Le simogramme est un outil essentiel pour visualiser et optimiser les cycles
    \item Le TRS, le Lead Time et le taux de valeur ajoutée sont des indicateurs clés
\end{itemize}
\end{minipage}}
\end{center}

% ==================== FIN DE LA PARTIE 8 ====================
% ==================== PARTIE 9 : LIMITES DU CHRONOMÉTRAGE ET ALTERNATIVES ====================

\section{Limites du Chronométrage et Alternatives}

Malgré son utilité et sa large diffusion, le chronométrage présente des limites importantes qui ont conduit au développement de méthodes alternatives. La connaissance de ces limites est essentielle pour choisir la technique la plus adaptée à chaque situation.

\subsection{Problèmes liés à la subjectivité}

\subsubsection{La subjectivité du jugement d'allure}

Le jugement d'allure est la principale source de subjectivité dans le chronométrage traditionnel.

\begin{definition}
Le \textbf{jugement d'allure} est une estimation subjective de la cadence de travail de l'opérateur observé. Malgré l'entraînement, cette évaluation reste influencée par des facteurs personnels.
\end{definition}

\subsubsection{Facteurs influençant le jugement d'allure}

\begin{table}[ht]
\centering
\begin{tabular}{|p{5cm}|p{9cm}|}
\hline
\textbf{Facteur} & \textbf{Influence sur le jugement} \\
\hline
Expérience du chronométreur & Un débutant a tendance à sous-évaluer ou surévaluer \\
Fatigue & La fatigue peut altérer la perception du rythme \\
Préjugés & L'apparence physique peut influencer le jugement \\
Contexte social & La pression hiérarchique peut biaiser l'évaluation \\
Morphologie & Les gestes amples peuvent sembler plus lents \\
\hline
\end{tabular}
\caption{Facteurs influençant le jugement d'allure}
\end{table}

\subsubsection{Conséquences de la subjectivité}

\begin{itemize}
    \item \textbf{Variabilité entre chronométreurs} : deux chronométreurs peuvent attribuer des JA différents pour le même opérateur
    \item \textbf{Variabilité dans le temps} : un même chronométreur peut varier son jugement selon son état
    \item \textbf{Risque de contestation} : les opérateurs peuvent contester les temps alloués
    \item \textbf{Difficulté de reproductibilité} : il est difficile de reproduire exactement les mêmes conditions
\end{itemize}

\begin{center}
\fbox{\begin{minipage}{0.85\textwidth}
\centering
\textbf{Conséquence majeure} \\
La subjectivité du jugement d'allure peut conduire à des écarts de 15 à 20\% \\
entre différents chronométreurs pour une même observation.
\end{minipage}}
\end{center}

\subsubsection{Les limites de l'étalonnage}

Même avec un étalonnage régulier, certaines difficultés persistent :

\begin{itemize}
    \item L'étalonnage ne peut éliminer totalement la subjectivité
    \item Les sessions d'étalonnage ne reproduisent pas tous les contextes
    \item La variabilité inter-individuelle reste significative
\end{itemize}

\subsection{Nécessité d'un poste existant et stabilisé}

\subsubsection{Contrainte d'existence du poste}

Le chronométrage nécessite que le poste de travail existe physiquement.

\begin{description}[labelwidth=4cm]
    \item[Pour les nouveaux produits] Impossible de chronométrer avant la mise en production
    \item[Pour les réorganisations] Difficulté d'évaluer des configurations non existantes
    \item[Pour les appels d'offres] Nécessité d'estimer des temps sans disposer du poste
\end{description}

\subsubsection{Nécessité de stabilisation}

Le chronométrage exige que le poste soit stabilisé :

\begin{itemize}
    \item L'opérateur doit avoir dépassé sa phase d'apprentissage
    \item Le mode opératoire doit être figé
    \item Les conditions de travail doivent être constantes
\end{itemize}

\begin{center}
\fbox{\begin{minipage}{0.85\textwidth}
\centering
\textbf{Paradoxe du chronométrage} \\
On ne peut chronométrer que ce qui existe déjà, \\
mais on a souvent besoin des temps avant la mise en production.
\end{minipage}}
\end{center}

\subsubsection{Contrainte de représentativité}

Les conditions de chronométrage doivent être représentatives :

\begin{itemize}
    \item L'opérateur chronométré doit être représentatif de la population
    \item Les conditions de travail doivent être normales
    \item Les variations saisonnières ne sont pas toujours prises en compte
\end{itemize}

\subsection{Introduction aux Temps Prédéterminés}

\subsubsection{Définition}

Les temps prédéterminés sont des méthodes qui permettent de déterminer le temps nécessaire à l'exécution d'une tâche sans avoir à observer l'opérateur, en se basant sur des tables de temps standards pour des mouvements élémentaires.

\begin{definition}
Les \textbf{temps prédéterminés} sont des systèmes de durées standardisées attribuées à des mouvements de base, permettant de calculer le temps d'une tâche par somme de ses composantes élémentaires.
\end{definition}

\subsubsection{Avantages des temps prédéterminés}

\begin{table}[ht]
\centering
\begin{tabular}{|p{5cm}|p{9cm}|}
\hline
\textbf{Avantage} & \textbf{Description} \\
\hline
Pas de jugement d'allure & Les temps sont prédéfinis, sans subjectivité \\
Anticipation & Possibilité de calculer avant la création du poste \\
Reproductibilité & Mêmes temps pour tous les utilisateurs \\
Homogénéité & Standards identiques dans toute l'entreprise \\
Formation & Formation structurée et certifiante \\
\hline
\end{tabular}
\caption{Avantages des temps prédéterminés}
\end{table}

\subsubsection{Les principales méthodes de temps prédéterminés}

\begin{enumerate}
    \item \textbf{MTM (Methods-Time Measurement)} - 1948
    \item \textbf{MOST (Maynard Operation Sequence Technique)} - 1974
    \item \textbf{MTS (Methods Time Standards)} - 1950
    \item \textbf{MODAPTS} - 1966
    \item \textbf{WORK FACTOR} - Philips
\end{enumerate}

\subsection{La méthode MTM (Methods-Time Measurement)}

\subsubsection{Historique et origine}

Développé en 1948 par H.B. Maynard, G.J. Stegemerten et J.L. Schwab, le MTM est la première méthode structurée de temps prédéterminés.

\subsubsection{Unité de mesure : le TMU}

\begin{definition}
Le \textbf{TMU (Time Measurement Unit)} est l'unité de base du MTM. \\
1 TMU = 0,00001 heure = 0,0006 minute = 0,036 seconde = 1 cmh.
\end{definition}

\subsubsection{Les mouvements de base du MTM}

\begin{table}[ht]
\centering
\begin{tabular}{|c|c|c|}
\hline
\textbf{Symbole} & \textbf{Mouvement} & \textbf{Description} \\
\hline
R & Atteindre & Déplacer la main vers une destination \\
M & Mouvoir & Déplacer un objet \\
T & Tourner & Rotation de la main ou du poignet \\
AP & Appliquer pression & Appliquer une force \\
G & Saisir & Prendre un objet \\
P & Positionner & Orienter un objet \\
RL & Relâcher & Lâcher un objet \\
D & Désengager & Séparer deux objets \\
ET & Déplacement yeux & Mouvement du regard \\
EF & Focalisation yeux & Mise au point visuelle \\
\hline
\end{tabular}
\caption{Mouvements de base du MTM}
\end{table}

\subsubsection{Les cas du mouvement "Atteindre"}

\begin{table}[ht]
\centering
\begin{tabular}{|c|c|}
\hline
\textbf{Cas} & \textbf{Description} \\
\hline
A & Objet toujours au même endroit ou dans l'autre main \\
B & Emplacement peut varier légèrement d'un cycle à l'autre \\
C & Objet mêlé à d'autres, nécessitant recherche \\
D & Objet très petit ou nécessitant précision \\
E & Vers une position indéfinie (équilibre, dégagement) \\
\hline
\end{tabular}
\caption{Les 5 cas du mouvement "Atteindre" (R)}
\end{table}

\subsubsection{Extrait de table MTM}

\begin{table}[ht]
\centering
\begin{tabular}{|c|c|c|c|c|c|c|}
\hline
\textbf{Distance (cm)} & \textbf{R-A} & \textbf{R-B} & \textbf{R-C} & \textbf{R-D} & \textbf{R-E} \\
\hline
2 & 2,0 & 2,0 & 2,0 & 2,0 & 1,6 \\
4 & 3,3 & 3,3 & 5,2 & 3,3 & 3,0 \\
6 & 4,5 & 4,5 & 6,5 & 4,5 & 3,9 \\
8 & 5,4 & 5,6 & 7,5 & 7,5 & 5,5 \\
10 & 6,0 & 6,6 & 8,4 & 6,4 & 4,9 \\
\hline
\end{tabular}
\caption{Extrait de la table MTM pour le mouvement "Atteindre" (valeurs en TMU)}
\end{table}

\subsubsection{Limites du MTM}

Malgré ses avantages, le MTM présente des limites :

\begin{itemize}
    \item \textbf{Complexité} : nécessite une formation longue (environ 20 jours)
    \item \textbf{Temps d'analyse} : une journée pour analyser 1 minute de cycle
    \item \textbf{Risque d'erreur} : élevé pour un utilisateur occasionnel
    \item \textbf{Coût} : formation et maintenance importantes
\end{itemize}

\subsection{La méthode MOST (Maynard Operation Sequence Technique)}

\subsubsection{Historique et origine}

Développé en 1974 par Kjell Zandin, le MOST est une évolution du MTM visant à simplifier et accélérer l'analyse.

\subsubsection{Principe du MOST}

Le MOST utilise des séquences de base pour décrire les mouvements :

\[
\text{ABG ABP A}
\]

Où :
\begin{itemize}
    \item \textbf{A} : Action (mouvement général)
    \item \textbf{B} : Body (mouvement du corps)
    \item \textbf{G} : Get (saisir)
    \item \textbf{P} : Place (positionner)
\end{itemize}

\subsubsection{Les trois types de séquences MOST}

\begin{table}[ht]
\centering
\begin{tabular}{|c|c|c|}
\hline
\textbf{Type} & \textbf{Formule} & \textbf{Application} \\
\hline
Séquence 1 & ABG ABP A & Déplacement d'objets \\
Séquence 2 & ABG ABP A & Mouvement guidé \\
Séquence 3 & ABG ABP A & Utilisation d'outil à main \\
\hline
\end{tabular}
\caption{Les trois types de séquences MOST}
\end{table}

\subsubsection{Les valeurs clés du Basic MOST}

Les valeurs sont très facilement mémorisables :

\begin{center}
\begin{tabular}{|c|c|}
\hline
\textbf{Valeur (TMU)} & \textbf{Signification} \\
\hline
10 & À portée de main \\
30 & Mouvement limité \\
60 & Mouvement étendu \\
100 & Déplacement plus long \\
160 & Déplacement très long \\
\hline
\end{tabular}
\caption{Valeurs clés du Basic MOST}
\end{center}

\begin{center}
\fbox{\begin{minipage}{0.85\textwidth}
\centering
\textbf{Avantage du MOST} \\
Une seule valeur de 10 TMU remplace environ 136 valeurs de la table MTM1.
\end{minipage}}
\end{center}

\subsubsection{Exemple d'application MOST}

\begin{exemple}
Saisir un marqueur situé à trois pas au sol et le déposer sur une tablette.

\begin{itemize}
    \item Séquence : ABG ABP A
    \item Indices : 6 6 1 6 0 1 0
    \item Somme : 6 + 6 + 1 + 6 + 0 + 1 + 0 = 20
    \item Temps : $20 \times 10 = 200$ TMU
\end{itemize}
\end{exemple}

\subsubsection{Avantages du MOST}

\begin{itemize}
    \item \textbf{Rapidité} : une heure de cycle chiffrée en une journée d'étude
    \item \textbf{Simplicité} : valeurs facilement mémorisables
    \item \textbf{Formation courte} : 6 jours maximum
    \item \textbf{Fiabilité} : précision maîtrisée par regroupements statistiques
\end{itemize}

\subsection{Autres méthodes de temps prédéterminés}

\subsubsection{MTS (Methods Time Standards)}

Développé par la General Electric (1950), le MTS est utilisé notamment par le BTE (Bureau des Temps Élémentaires).

\subsubsection{MODAPTS (Modular Arrangement of Predetermined Time Standards)}

Développé par G.C. Hyde (1966), le MODAPTS a été popularisé par Renault.

\subsubsection{WORK FACTOR}

Développé par Philips, le WORK FACTOR est particulièrement utilisé dans l'industrie électronique.

\subsection{Comparaison chronométrage vs temps prédéterminés}

\begin{table}[ht]
\centering
\begin{tabular}{|p{4.5cm}|p{4.5cm}|p{4.5cm}|}
\hline
\textbf{Caractéristique} & \textbf{Chronométrage} & \textbf{Temps prédéterminés} \\
\hline
Subjectivité & Forte (JA) & Nulle \\
Poste nécessaire & Oui (existant) & Non (peut être virtuel) \\
Formation requise & Quelques jours & Plusieurs jours (MTM: 20 jours) \\
Temps d'analyse & Variable & MTM: 1 min/h; MOST: 1 h/j \\
Précision & Dépend du JA & Standardisée \\
Coût & Modéré & Élevé à l'achat \\
Acceptation sociale & Contestations possibles & Meilleure (neutre) \\
\hline
\end{tabular}
\caption{Comparaison entre chronométrage et temps prédéterminés}
\end{table}

\subsection{Choix de la méthode adaptée}

\subsubsection{Critères de choix}

\begin{itemize}
    \item \textbf{Objectif} : temps de référence, optimisation, devis
    \item \textbf{Contexte} : poste existant ou nouveau produit
    \item \textbf{Ressources} : compétences disponibles, budget
    \item \textbf{Délais} : urgence de la demande
    \item \textbf{Précision requise} : estimation ou valeur contractuelle
\end{itemize}

\subsubsection{Recommandations pratiques}

\begin{center}
\fbox{\begin{minipage}{0.9\textwidth}
\textbf{Recommandations :}

\begin{itemize}
    \item Utiliser le chronométrage pour les postes existants stabilisés
    \item Utiliser le MOST pour les études rapides et les nouveaux produits
    \item Utiliser le MTM pour les analyses très fines et les standards d'entreprise
    \item Combiner les approches : chronométrage pour calibration, temps prédéterminés pour les nouvelles gammes
\end{itemize}
\end{minipage}}
\end{center}

\subsection{Synthèse}

\begin{center}
\fbox{\begin{minipage}{0.9\textwidth}
\textbf{Points clés à retenir :}

\begin{itemize}
    \item Le chronométrage présente des limites : subjectivité du JA, nécessité d'un poste existant
    \item Les temps prédéterminés éliminent la subjectivité et permettent l'anticipation
    \item Le MTM est une méthode fine mais complexe (formation 20 jours)
    \item Le MOST est plus rapide (1 heure de cycle analysée par jour)
    \item Le choix de la méthode dépend du contexte, des objectifs et des ressources
    \item Les approches peuvent être combinées pour optimiser les résultats
\end{itemize}
\end{minipage}}
\end{center}

% ==================== FIN DE LA PARTIE 9 ====================
% ==================== CONCLUSION ====================

\section*{Conclusion}
\addcontentsline{toc}{section}{Conclusion}

\subsection*{Synthèse des bonnes pratiques pour un chronométrage fiable}
\addcontentsline{toc}{subsection}{Synthèse des bonnes pratiques pour un chronométrage fiable}

Au terme de cette étude sur le chronométrage, il est essentiel de synthétiser les bonnes pratiques qui garantissent la fiabilité et la pertinence des mesures obtenues.

\subsubsection*{Avant le chronométrage : la préparation}

\begin{enumerate}
    \item \textbf{Vérifier les conditions préalables}
    \begin{itemize}
        \item L'opérateur a dépassé sa phase d'apprentissage (courbe stabilisée)
        \item Le poste de travail est aménagé, alimenté et fonctionne normalement
        \item Les critères de qualité sont clairement définis et respectés
        \item Le chronométreur maîtrise les techniques de mesure et le jugement d'allure
    \end{itemize}
    
    \item \textbf{Choisir le type de chronométrage adapté}
    \begin{itemize}
        \item Diagnostic rapide : 2-3 relevés sans JA
        \item Analyse de stabilisation : 10-15 relevés sans JA
        \item Fixation des temps : 20-30 relevés avec JA
    \end{itemize}
    
    \item \textbf{Préparer la feuille de relevés}
    \begin{itemize}
        \item Décomposer le travail en séquences avec des "tops" clairs
        \item Distinguer temps cycliques et temps fréquentiels
        \item Prévoir les colonnes pour les relevés et les observations
    \end{itemize}
\end{enumerate}

\subsubsection*{Pendant le chronométrage : l'exécution}

\begin{enumerate}
    \item \textbf{Adopter une position adaptée}
    \begin{itemize}
        \item Se tenir debout, à gauche de l'opérateur
        \item Avoir le chronomètre et les mains de l'opérateur dans le champ visuel
        \item Ne pas gêner le travail
    \end{itemize}
    
    \item \textbf{Effectuer les relevés avec rigueur}
    \begin{itemize}
        \item Utiliser le chronométrage continu pour les temps cycliques
        \item Noter immédiatement les temps et les jugements d'allure
        \item Entourer et noter les valeurs aberrantes sans les effacer
    \end{itemize}
    
    \item \textbf{Assurer le contrôle de précision}
    \begin{itemize}
        \item Utiliser la méthode du bouclage (deux chronomètres)
        \item Vérifier l'écart de bouclage (< 2\% sans JA, < 5\% avec JA)
    \end{itemize}
\end{enumerate}

\subsubsection*{Après le chronométrage : le dépouillement}

\begin{enumerate}
    \item \textbf{Traiter les données avec méthode}
    \begin{itemize}
        \item Calculer le temps observé moyen après élimination des aberrantes
        \item Corriger par le jugement d'allure : $T_{100} = \overline{T}_o \times JA/100$
        \item Appliquer les coefficients majorateurs adaptés
        \item Intégrer la fréquence des séquences
    \end{itemize}
    
    \item \textbf{Valider les résultats}
    \begin{itemize}
        \item Vérifier la cohérence des temps obtenus
        \item Comparer avec les temps historiques si disponibles
        \item Réaliser des chronométrages de confirmation si nécessaire
    \end{itemize}
\end{enumerate}

\begin{center}
\fbox{\begin{minipage}{0.9\textwidth}
\centering
\textbf{Les 10 commandements du chronométreur}

\begin{enumerate}
    \item Prépare soigneusement ton étude avant de commencer
    \item Vérifie la stabilisation du poste et de l'opérateur
    \item Explique ton travail aux opérateurs et gagne leur confiance
    \item Choisis l'unité de temps adaptée à ton étude
    \item Décompose le travail en séquences avec des "tops" clairs
    \item Positionne-toi correctement pour bien observer
    \item Enregistre les temps et les jugements d'allure simultanément
    \item Utilise le bouclage pour contrôler ta précision
    \item Traite les données avec rigueur et sans biais
    \item Valide tes résultats avant de les diffuser
\end{enumerate}
\end{minipage}}
\end{center}

\subsection*{Rôle du chronométrage dans l'amélioration continue de la productivité}
\addcontentsline{toc}{subsection}{Rôle du chronométrage dans l'amélioration continue de la productivité}

\subsubsection*{Le chronométrage comme outil de pilotage}

Le chronométrage n'est pas une fin en soi, mais un outil au service de l'amélioration continue. Il contribue à la productivité à plusieurs niveaux :

\begin{itemize}
    \item \textbf{Objectivation des performances}
    \begin{itemize}
        \item Fournit des données objectives sur les temps de production
        \item Permet de comparer les performances entre opérateurs et entre ateliers
        \item Établit une base de référence pour mesurer les progrès
    \end{itemize}
    
    \item \textbf{Identification des goulots et des gaspillages}
    \begin{itemize}
        \item Le chronométrage de diagnostic révèle les postes problématiques
        \item L'analyse des temps met en évidence les déséquilibres de charge
        \item La décomposition des séquences permet d'identifier les temps sans valeur ajoutée
    \end{itemize}
    
    \item \textbf{Optimisation des méthodes}
    \begin{itemize}
        \item Le chronométrage d'étude permet de tester et valider les améliorations
        \item La comparaison avant/après mesure l'impact des actions
        \item L'analyse des écarts guide les axes de progrès
    \end{itemize}
\end{itemize}

\subsubsection*{Le chronométrage dans la démarche d'amélioration continue}

\begin{table}[ht]
\centering
\begin{tabular}{|p{4cm}|p{10cm}|}
\hline
\textbf{Étape PDCA} & \textbf{Apport du chronométrage} \\
\hline
Plan (Planifier) & Établir l'état initial, fixer des objectifs quantifiés \\
Do (Réaliser) & Former les opérateurs, stabiliser les méthodes \\
Check (Vérifier) & Mesurer les résultats, comparer aux objectifs \\
Act (Agir) & Valider les améliorations, standardiser \\
\hline
\end{tabular}
\caption{Le chronométrage dans la roue de Deming (PDCA)}
\end{table}

\subsubsection*{Vers une approche intégrée}

\begin{center}
\fbox{\begin{minipage}{0.9\textwidth}
\centering
\textbf{Le chronométrage au cœur du Lean Manufacturing}

\begin{itemize}
    \item Élimination des \textbf{Muda} (gaspillages) : le chronométrage révèle les temps improductifs
    \item \textbf{Juste-à-temps} : des temps fiables permettent des plannings précis
    \item \textbf{Standardisation} : le chronométrage fixe les standards de travail
    \item \textbf{Kaizen} : l'amélioration continue s'appuie sur des mesures objectives
    \item \textbf{Management visuel} : les temps standards sont affichés et suivis
\end{itemize}
\end{minipage}}
\end{center}

\subsubsection*{Articulation avec les autres techniques}

Le chronométrage s'inscrit dans un ensemble plus large d'outils d'optimisation :

\begin{table}[ht]
\centering
\begin{tabular}{|p{4.5cm}|p{4.5cm}|p{5cm}|}
\hline
\textbf{Outil} & \textbf{Objectif} & \textbf{Complémentarité} \\
\hline
Simogramme & Optimisation homme-machine & Visualise les temps masqués \\
Observations instantanées & Répartition des activités & Complète l'analyse des temps \\
MTM / MOST & Temps prédéterminés & Anticipe avant création du poste \\
Analyse de déroulement & Cartographie des flux & Identifie les NVA \\
VSM & Cartographie de la valeur & Vision globale du processus \\
\hline
\end{tabular}
\caption{Articulation du chronométrage avec les autres outils}
\end{table}

\subsubsection*{Facteurs clés de succès}

Pour que le chronométrage contribue efficacement à l'amélioration de la productivité, plusieurs facteurs sont essentiels :

\begin{enumerate}
    \item \textbf{Engagement de la direction} : soutien visible et constant
    \item \textbf{Implication des opérateurs} : communication, transparence, participation
    \item \textbf{Formation des chronométreurs} : compétences techniques et relationnelles
    \item \textbf{Utilisation cohérente des résultats} : temps utilisés pour la gestion, pas seulement pour le contrôle
    \item \textbf{Mise à jour régulière} : les temps évoluent, ils doivent être révisés
    \item \textbf{Articulation avec la stratégie} : alignement avec les objectifs de l'entreprise
\end{enumerate}

\subsection*{Perspectives}

L'évolution des technologies ouvre de nouvelles perspectives pour la mesure du travail :

\begin{itemize}
    \item \textbf{Capture vidéo automatisée} : analyse des mouvements par intelligence artificielle
    \item \textbf{Capteurs connectés} : mesure en temps réel des temps opératoires
    \item \textbf{Logiciels intégrés} : couplage avec les ERP et les systèmes de GPAO
    \item \textbf{Analyse prédictive} : anticipation des temps à partir de données historiques
\end{itemize}

\subsection*{Conclusion finale}

\begin{center}
\fbox{\begin{minipage}{0.9\textwidth}
\centering
\textbf{En synthèse}

Le chronométrage reste un outil fondamental de l'organisation industrielle. \\
Sa force réside dans sa capacité à fournir des données objectives \\
pour piloter la production, évaluer les performances et améliorer la productivité.

Cependant, sa pertinence dépend de la rigueur de sa mise en œuvre \\
et de l'intégration de ses résultats dans une démarche globale \\
d'amélioration continue associant tous les acteurs de l'entreprise.

Bien utilisé, le chronométrage n'est pas un outil de contrôle, \\
mais un levier de progrès au service de la performance collective.
\end{minipage}}
\end{center}

\vspace{1cm}

\begin{flushright}
\textit{« Ce qui se mesure s'améliore. »} \\
— Peter Drucker
\end{flushright}

% ==================== FIN DE LA CONCLUSION ====================
% ==================== PARTIE 5 : TECHNIQUES COMPLÉMENTAIRES DE MESURE ====================
\part*{Partie 5 : Techniques Complémentaires de Mesure}
\addcontentsline{toc}{part}{Partie 5 : Techniques Complémentaires de Mesure}

\section{Observations Instantanées}

Les observations instantanées (ou méthode des échantillonnages) constituent une alternative au chronométrage continu, particulièrement adaptée pour analyser la répartition des activités d'une ressource sans avoir à suivre en continu son travail.

\subsection{Définition et objectifs}

\begin{definition}
Les \textbf{observations instantanées} sont une technique de mesure du travail qui consiste à effectuer des observations ponctuelles et aléatoires d'une ressource pour déterminer, par échantillonnage statistique, la proportion de temps qu'elle consacre à différentes activités.
\end{definition}

\textbf{Objectifs principaux :}
\begin{itemize}
    \item Connaître la répartition des activités pour une ressource donnée (humaine ou matérielle)
    \item Éviter l'observation continue longue et coûteuse
    \item Obtenir des données statistiquement représentatives
    \item Diagnostiquer rapidement l'utilisation des ressources
    \item Identifier les temps improductifs et les sources de gaspillage
    \item Servir de base à des décisions d'organisation ou d'investissement
\end{itemize}

\subsection{Différence avec le chronométrage continu}

\begin{table}[ht]
\centering
\begin{tabular}{|p{5cm}|p{4.5cm}|p{4.5cm}|}
\hline
\textbf{Caractéristique} & \textbf{Observations instantanées} & \textbf{Chronométrage continu} \\
\hline
Nature de l'observation & Ponctuelle, aléatoire & Continue, séquentielle \\
Durée d'observation & Courte par passage & Longue (cycles complets) \\
Résultat obtenu & Proportions d'activités & Temps de cycle, temps standard \\
Précision & Statistique & Individuelle \\
Opérateur & Peut s'auto-observer & Observateur externe \\
Application & Répartition d'activités & Temps de production \\
\hline
\end{tabular}
\caption{Comparaison observations instantanées vs chronométrage continu}
\end{table}

\begin{center}
\fbox{\begin{minipage}{0.85\textwidth}
\centering
\textbf{Quand utiliser les observations instantanées ?}
\begin{itemize}
    \item Pour analyser l'activité d'un grand nombre de ressources
    \item Pour identifier les temps improductifs
    \item Lorsque les cycles sont longs ou irréguliers
    \item Pour une première approche avant un chronométrage détaillé
\end{itemize}
\end{minipage}}
\end{center}

\subsection{L'auto-pointage}

L'auto-pointage est une variante des observations instantanées où les opérateurs effectuent eux-mêmes leurs relevés de temps.

\subsubsection{Définition}

\begin{definition}
L'\textbf{auto-pointage} est une méthode de relevé dans laquelle les opérateurs enregistrent eux-mêmes, sur une feuille dédiée, le début et la fin de leurs différentes activités tout au long de la journée.
\end{definition}

\subsubsection{Préparation de l'auto-pointage}

Avant de commencer l'auto-pointage, il faut :
\begin{itemize}
    \item Identifier toutes les activités prévues
    \item Créer un tableau de relevés adapté
    \item Former les opérateurs à la méthode
    \item Expliquer les objectifs et l'utilisation des résultats
\end{itemize}

\begin{table}[ht]
\centering
\begin{tabular}{|p{4cm}|p{6cm}|p{4cm}|}
\hline
\textbf{Horaire de début} & \textbf{Activité} & \textbf{Horaire de fin} \\
\hline
8H00 & Installation poste & 8H15 \\
8H15 & Production série A & 10H30 \\
10H30 & Pause & 10H45 \\
10H45 & Production série A & 12H00 \\
12H00 & Déjeuner & 13H00 \\
13H00 & Réglage machine & 13H30 \\
13H30 & Production série B & 16H00 \\
16H00 & Rangement & 16H30 \\
\hline
\end{tabular}
\caption{Exemple de feuille d'auto-pointage}
\end{table}

\subsubsection{Réalisation de l'auto-pointage}

Pendant la période d'auto-pointage :
\begin{itemize}
    \item L'opérateur remplit le tableau en temps réel
    \item Chaque changement d'activité est enregistré
    \item Les informations doivent être précises et sincères
    \item La durée de l'auto-pointage est définie à l'avance (1 jour, 1 semaine)
\end{itemize}

\subsubsection{Exploitation de l'auto-pointage}

Après la période d'auto-pointage :
\begin{itemize}
    \item Saisie des données dans un tableur
    \item Calcul des durées par activité
    \item Calcul des proportions par activité
    \item Analyse des résultats et identification des axes d'amélioration
\end{itemize}

\subsection{Méthodologie en 6 étapes}

\subsubsection{Étape 1 : Définir les ressources, le niveau de confiance et la marge d'erreur}

\begin{description}[labelwidth=4cm]
    \item[Ressource] La ressource étudiée peut être humaine (opérateur, cariste) ou matérielle (machine, robot). Elle occupe différents états au cours de la journée. Il ne faut oublier aucun état important.
    
    \item[Niveau de confiance] Généralement fixé à 95\%. Cela signifie que dans 95\% des cas, les résultats énoncés seront vrais.
    
    \item[Marge d'erreur absolue E] Demi-intervalle de tolérance. Par exemple, si un opérateur passe entre 52\% et 58\% du temps à régler les machines, on écrira $55\% \pm 3\%$ avec $E = 3\%$.
\end{description}

\begin{exemple}
On souhaite déterminer la répartition des activités d'un cariste. Les états possibles sont :
\begin{itemize}
    \item \textbf{A} : Actif sur son chariot
    \item \textbf{B} : Actif hors chariot (conditionnement)
    \item \textbf{C} : Autre (attente, pause, entretien)
\end{itemize}
Niveau de confiance : $95\%$
Marge d'erreur absolue $E = 10\%$ (maximum)
\end{exemple}

\subsubsection{Étape 2 : Estimer la proportion des états (pré-étude)}

La pré-étude permet d'estimer les proportions initiales. Elle peut être réalisée par :
\begin{itemize}
    \item Auto-pointage préliminaire
    \item Exploitation de relevés existants
    \item Observation continue rapide
    \item Jugement d'experts
\end{itemize}

\begin{exemple}
Après une pré-étude, les proportions estimées sont :
\begin{itemize}
    \item État A (actif sur chariot) : $p_A = 40\%$
    \item État B (actif hors chariot) : $p_B = 25\%$
    \item État C (autre) : $p_C = 35\%$
\end{itemize}
\end{exemple}

\subsubsection{Étape 3 : Calculer le nombre d'observations}

Pour un niveau de confiance de 95\%, on utilise la formule :

\[
n = \frac{1,96^2 \times p \times (1-p)}{E^2}
\]

Où :
\begin{itemize}
    \item $n$ : nombre d'observations nécessaires
    \item $p$ : proportion estimée de l'état considéré
    \item $E$ : marge d'erreur absolue souhaitée
\end{itemize}

Il faut calculer $n$ pour chaque état. La valeur la plus élevée est retenue.

\begin{exemple}
Calcul pour l'état A ($p_A = 40\%$, $E = 10\%$) :
\[
n_A = \frac{1,96^2 \times 0,4 \times 0,6}{0,10^2} = \frac{3,8416 \times 0,24}{0,01} = 92,2 \approx 93
\]

Calcul pour l'état B ($p_B = 25\%$, $E = 10\%$) :
\[
n_B = \frac{1,96^2 \times 0,25 \times 0,75}{0,01} = \frac{3,8416 \times 0,1875}{0,01} = 72,0 \approx 72
\]

Calcul pour l'état C ($p_C = 35\%$, $E = 10\%$) :
\[
n_C = \frac{1,96^2 \times 0,35 \times 0,65}{0,01} = \frac{3,8416 \times 0,2275}{0,01} = 87,4 \approx 88
\]

On retient $n = 93$ observations (la valeur la plus élevée).
\end{exemple}

\subsubsection{Étape 4 : Déterminer la durée des observations}

Pour réaliser les observations, l'observateur doit :
\begin{itemize}
    \item Quitter son lieu de travail habituel
    \item Aller voir chacune des ressources à observer
    \item Noter leur état
    \item Revenir à son lieu de travail
\end{itemize}

\[
\text{Temps d'occupation} = D \times n + \text{Temps de préparation} + \text{Temps d'exploitation}
\]

Où $D$ est la durée d'un parcours complet d'observation.

\begin{exemple}
\begin{itemize}
    \item Durée d'un parcours d'observation : 3 minutes
    \item Nombre d'observations : 93
    \item Temps de préparation : 2 heures
    \item Temps d'exploitation : 3 heures
\end{itemize}

\[
\text{Temps d'occupation} = (3 \times 93) + 120 + 180 = 279 + 300 = 579 \text{ minutes} \approx 9,7 \text{ heures}
\]
\end{exemple}

\textbf{Règle :} Le temps d'occupation de l'observateur doit être inférieur à 30\% de son temps de travail total.

\subsubsection{Étape 5 : Déterminer les horaires d'observation}

Les observations doivent avoir lieu à des horaires définis à l'avance et connus uniquement de l'observateur pour éviter tout biais.

\textbf{Méthode :}
\begin{enumerate}
    \item Définir la plage horaire d'observation (ex: 8h00 - 17h00)
    \item Diviser cette plage en intervalles (ex: 30 minutes)
    \item Générer des nombres aléatoires pour sélectionner les horaires
    \item Utiliser un générateur aléatoire (Excel, table de nombres aléatoires)
\end{enumerate}

\begin{exemple}
Pour 93 observations sur 5 jours (8h00-17h00) :
\begin{itemize}
    \item Plage totale : 9 heures par jour = 540 minutes
    \item Intervalle de tirage : $540 \times 5 / 93 \approx 29$ minutes
    \item Générer 93 horaires aléatoires entre 8h00 et 17h00
\end{itemize}
\end{exemple}

\subsubsection{Étape 6 : Conclure l'étude et valider la marge d'erreur}

À l'issue de la période d'observations, on calcule pour chaque état la proportion observée $p_{\text{obs}}$ et la marge d'erreur réelle :

\[
E_{\text{réelle}} = \sqrt{\frac{1,96^2 \times p_{\text{obs}} \times (1-p_{\text{obs}})}{n}}
\]

L'étude est validée si $E_{\text{réelle}}$ est inférieur ou égal à $E$ fixé initialement.

\begin{exemple}
Après 93 observations, les proportions observées sont :
\begin{itemize}
    \item État A : 38 observations $\rightarrow p_A = 38/93 = 40,9\%$
    \item État B : 22 observations $\rightarrow p_B = 22/93 = 23,7\%$
    \item État C : 33 observations $\rightarrow p_C = 33/93 = 35,5\%$
\end{itemize}

Calcul de la marge d'erreur réelle pour l'état A :
\[
E_A = \sqrt{\frac{1,96^2 \times 0,409 \times (1-0,409)}{93}} = \sqrt{\frac{3,8416 \times 0,242}{93}} = \sqrt{0,00999} \approx 0,10 = 10\%
\]

L'étude est validée ($E_{\text{réelle}} \leq 10\%$).
\end{exemple}

\subsection{Exercice d'application : Étude de la répartition des activités d'un cariste}

\subsubsection{Énoncé}

Une entreprise logistique souhaite analyser l'activité de ses caristes pour optimiser l'organisation du travail. Un cariste peut être dans les états suivants :

\begin{itemize}
    \item \textbf{A} : Déplacement à vide (chariot sans charge)
    \item \textbf{B} : Déplacement en charge (transport de palette)
    \item \textbf{C} : Chargement/déchargement
    \item \textbf{D} : Attente
    \item \textbf{E} : Entretien/repères
\end{itemize}

Une pré-étude par auto-pointage a donné les proportions suivantes :
\begin{itemize}
    \item $p_A = 30\%$
    \item $p_B = 25\%$
    \item $p_C = 20\%$
    \item $p_D = 15\%$
    \item $p_E = 10\%$
\end{itemize}

Le niveau de confiance est fixé à 95\%. La marge d'erreur absolue souhaitée est de 5\%.

\subsubsection{Questions}

\begin{enumerate}
    \item Calculer le nombre d'observations nécessaires pour cette étude.
    \item La durée d'un parcours d'observation est de 5 minutes. Le temps de préparation est estimé à 3 heures et le temps d'exploitation à 4 heures. Calculer le temps d'occupation total de l'observateur.
    \item L'entreprise souhaite réaliser l'étude sur 10 jours de travail (8h00-17h00). Déterminer le nombre d'observations à réaliser par jour.
    \item Après réalisation des observations, les résultats obtenus sont les suivants :
    \begin{itemize}
        \item État A : 85 observations
        \item État B : 72 observations
        \item État C : 58 observations
        \item État D : 42 observations
        \item État E : 28 observations
    \end{itemize}
    Calculer les proportions observées et vérifier si la marge d'erreur est respectée.
\end{enumerate}

\subsubsection{Corrigé}

\textbf{Question 1 : Calcul du nombre d'observations}

Pour chaque état, on applique $n = \frac{1,96^2 \times p \times (1-p)}{E^2}$ avec $E = 0,05$ :

\begin{itemize}
    \item État A : $n_A = \frac{3,8416 \times 0,3 \times 0,7}{0,0025} = \frac{0,8067}{0,0025} = 322,7 \approx 323$
    \item État B : $n_B = \frac{3,8416 \times 0,25 \times 0,75}{0,0025} = \frac{0,7203}{0,0025} = 288,1 \approx 289$
    \item État C : $n_C = \frac{3,8416 \times 0,2 \times 0,8}{0,0025} = \frac{0,6147}{0,0025} = 245,9 \approx 246$
    \item État D : $n_D = \frac{3,8416 \times 0,15 \times 0,85}{0,0025} = \frac{0,4898}{0,0025} = 195,9 \approx 196$
    \item État E : $n_E = \frac{3,8416 \times 0,1 \times 0,9}{0,0025} = \frac{0,3457}{0,0025} = 138,3 \approx 139$
\end{itemize}

\textbf{Nombre total d'observations :} $n = 323$ (la valeur la plus élevée)

\textbf{Question 2 : Temps d'occupation}

\[
\text{Temps d'occupation} = (5 \times 323) + (3 \times 60) + (4 \times 60) = 1615 + 180 + 240 = 2035 \text{ minutes} \approx 33,9 \text{ heures}
\]

\textbf{Question 3 : Observations par jour}

10 jours de travail, 9 heures par jour (540 minutes) :
\[
\text{Observations par jour} = \frac{323}{10} \approx 33 \text{ observations/jour}
\]

\textbf{Question 4 : Validation des résultats}

Nombre total d'observations réalisées : $85 + 72 + 58 + 42 + 28 = 285$

Proportions observées :
\begin{itemize}
    \item État A : $85/285 = 29,8\%$
    \item État B : $72/285 = 25,3\%$
    \item État C : $58/285 = 20,4\%$
    \item État D : $42/285 = 14,7\%$
    \item État E : $28/285 = 9,8\%$
\end{itemize}

Vérification de la marge d'erreur (pour l'état A) :
\[
E_A = \sqrt{\frac{1,96^2 \times 0,298 \times 0,702}{285}} = \sqrt{\frac{3,8416 \times 0,209}{285}} = \sqrt{\frac{0,8029}{285}} = \sqrt{0,00282} \approx 0,053 = 5,3\%
\]

L'étude est validée ($E_{\text{réelle}} \approx 5,3\%$ proche de l'objectif 5\%).

\subsection{Synthèse}

\begin{center}
\fbox{\begin{minipage}{0.9\textwidth}
\textbf{Points clés à retenir :}

\begin{itemize}
    \item Les observations instantanées permettent de connaître la répartition des activités d'une ressource
    \item La méthode est basée sur un échantillonnage statistique avec un niveau de confiance de 95\%
    \item Le nombre d'observations se calcule à partir de la proportion estimée et de la marge d'erreur
    \item Les horaires d'observation doivent être aléatoires pour éviter les biais
    \item L'auto-pointage est une variante où les opérateurs s'auto-observent
    \item La méthode est particulièrement adaptée pour analyser l'utilisation des ressources et identifier les temps improductifs
\end{itemize}
\end{minipage}}
\end{center}

% ==================== FIN DE LA PARTIE 10 ====================
% ==================== PARTIE 11 : OUTILS DE REPRÉSENTATION DES TEMPS ====================

\section{Outils de Représentation des Temps}

La représentation graphique des temps et des flux est essentielle pour visualiser, analyser et optimiser les processus de production. Plusieurs outils complémentaires permettent de cartographier le travail sous différents angles.

\subsection{Analyse de Déroulement}

L'analyse de déroulement est une méthode de relevé de données utilisée lorsqu'on pratique la "marche le long du processus".

\subsubsection{Définition et objectifs}

\begin{definition}
L'\textbf{analyse de déroulement} est une méthode d'observation qui consiste à suivre les étapes de transformation de la matière première en produit fini, en notant les temps, les distances, les poids et en précisant s'il s'agit de valeur ajoutée (VA) ou de non-valeur ajoutée (NVA).
\end{definition}

\textbf{Objectifs :}
\begin{itemize}
    \item Visualiser l'ensemble des étapes d'un processus
    \item Identifier les activités à valeur ajoutée et celles sans valeur ajoutée
    \item Mesurer les temps, distances et quantités
    \item Détecter les gaspillages et les opportunités d'amélioration
    \item Servir de base à l'optimisation des flux
\end{itemize}

\subsubsection{Symboles d'analyse}

L'analyse de déroulement utilise des symboles normalisés pour représenter les différentes activités :

\begin{table}[ht]
\centering
\begin{tabular}{|c|c|c|}
\hline
\textbf{Symbole} & \textbf{Signification} & \textbf{Description} \\
\hline
$\bigcirc$ & Opération & Transformation, usinage, montage, assemblage \\
$\rightarrow$ & Transfert & Déplacement de matière, manutention \\
$\square$ & Contrôle & Vérification, inspection, test qualité \\
$\bigcirc\!\!\!/$ & Attente & Temps d'attente, stockage temporaire \\
$\bigtriangleup$ & Stockage & Stockage permanent, mise en magasin \\
\hline
\end{tabular}
\caption{Symboles de base de l'analyse de déroulement}
\end{table}

\begin{center}
\fbox{\begin{minipage}{0.85\textwidth}
\centering
\textbf{Cinq symboles systémiques complémentaires :}
\begin{itemize}
    \item Ligne pointillée : Limite système (périmètre étudié)
    \item Ligne pleine : Flux (liaison entre étapes)
    \item Y : Source (origine des approvisionnements)
    \item Puits : Destination (point d'aboutissement)
    \item Ordre : Prise de décision, régulation
\end{itemize}
\end{minipage}}
\end{center}

\subsubsection{Indicateurs de performance}

L'analyse de déroulement permet de calculer plusieurs indicateurs clés :

\begin{description}[labelwidth=5cm]
    \item[Efficacité du processus] Proportion d'étapes à valeur ajoutée :
    \[
    \text{Efficacité} = \frac{\text{Nombre d'étapes VA}}{\text{Nombre total d'étapes}}
    \]
    
    \item[Temps de traversée (Lead Time - LT)] Temps total entre le début et la fin du processus :
    \[
    LT = \sum (\text{Temps VA}) + \sum (\text{Temps NVA})
    \]
    
    \item[Efficience du processus] Proportion de temps à valeur ajoutée :
    \[
    \text{Efficience} = \frac{\sum \text{Temps VA}}{LT}
    \]
    
    \item[Indice de tension du flux] Inverse de l'efficience :
    \[
    \text{Indice de tension} = \frac{1}{\text{Efficience}} = \frac{LT}{\sum \text{Temps VA}}
    \]
\end{description}

\begin{exemple}
Pour un processus de fabrication :
\begin{itemize}
    \item Nombre total d'étapes : 25
    \item Étapes à valeur ajoutée : 12
    \item Temps VA total : 45 min
    \item Temps NVA total : 135 min
\end{itemize}

\[
\text{Efficacité} = \frac{12}{25} = 48\%
\]
\[
LT = 45 + 135 = 180 \text{ min}
\]
\[
\text{Efficience} = \frac{45}{180} = 25\%
\]
\[
\text{Indice de tension} = \frac{180}{45} = 4
\]
\end{exemple}

\subsection{Graphique de Déroulement}

Le graphique de déroulement est une représentation visuelle des étapes d'un processus selon différentes perspectives.

\subsubsection{Types de graphiques de déroulement}

\begin{description}[labelwidth=4.5cm]
    \item[Graphique exécutant] Enregistre ce que fait le travailleur (mouvements, manipulations)
    \item[Graphique matière] Enregistre comment la matière est manutentionnée ou transformée
    \item[Graphique matériel] Enregistre comment l'équipement est utilisé (machine, outil)
\end{description}

\subsubsection{Exemple de graphique exécutant}

\begin{table}[ht]
\centering
\begin{tabular}{|c|c|c|c|c|c|}
\hline
\textbf{Étape} & \textbf{Description} & \textbf{Symbole} & \textbf{Temps (s)} & \textbf{Distance (m)} & \textbf{VA/NVA} \\
\hline
1 & Aller vers le placard & $\rightarrow$ & 5 & 3 & NVA \\
2 & Prendre le pot de farine & $\bigcirc$ & 3 & 0 & VA \\
3 & Aller vers le plan de travail & $\rightarrow$ & 5 & 3 & NVA \\
4 & Peser 300g de farine & $\bigcirc$ & 10 & 0 & VA \\
5 & Créer un puits & $\bigcirc$ & 4 & 0 & VA \\
6 & Aller vers le frigo & $\rightarrow$ & 8 & 5 & NVA \\
7 & Prendre des œufs & $\bigcirc$ & 3 & 0 & VA \\
8 & Aller vers le plan de travail & $\rightarrow$ & 8 & 5 & NVA \\
9 & Casser les œufs & $\bigcirc$ & 8 & 0 & VA \\
10 & Mélanger & $\bigcirc$ & 15 & 0 & VA \\
\hline
\end{tabular}
\caption{Exemple de graphique de déroulement exécutant (préparation de pâte)}
\end{table}

\subsection{Matrice de Déroulement}

La matrice de déroulement est un formulaire structuré permettant le recueil systématique d'informations lors de l'observation sur le terrain.

\subsubsection{Structure et informations collectées}

\begin{table}[ht]
\centering
\begin{tabular}{|p{2.5cm}|p{2.5cm}|p{2.5cm}|p{2.5cm}|p{4cm}|}
\hline
\textbf{Distance} & \textbf{Temps} & \textbf{Quantité} & \textbf{Poids} & \textbf{Déroulement} \\
\hline
70 m & 0,3 h & 150 & 300 kg & Sortie magasin matière \\
\hline
& 0,05 h/p & & & vers sciage \\
\hline
& 1,4 h & & & Sciage \\
\hline
& 0,12 h & & & Attente stabilisation \\
\hline
25 m & 0,12 h & 150 & 300 kg & vers fraisage \\
\hline
& 0,08 h/p & & & Fraisage \\
\hline
200 m & 0,9 h & 100 & 200 kg & vers MMT \\
\hline
& 0,5 h/p & & & Mesurage \\
\hline
50 m & 0,25 h & 100 & 200 kg & vers montage \\
\hline
\end{tabular}
\caption{Exemple de matrice de déroulement}
\end{table}

\textbf{Informations typiquement collectées :}
\begin{itemize}
    \item Identification de l'acteur (opérateur, magasinier, machine)
    \item Quantités concernées
    \item Durée de l'étape (mesurée, gammée, estimée)
    \item Fréquence de répétition
    \item Distance parcourue (pour les transferts)
    \item Nature de l'activité (VA/NVA)
\end{itemize}

\subsection{Value Stream Mapping (VSM)}

Le Value Stream Mapping est un outil puissant de cartographie des flux de valeur, largement utilisé dans les démarches Lean.

\subsubsection{Définition et objectifs}

\begin{definition}
Le \textbf{Value Stream Mapping (VSM)} est une représentation graphique qui retrace le flux physique des matières et le flux d'information correspondant tout au long du processus. C'est un outil de visualisation de la création de valeur au sein d'une entreprise.
\end{definition}

\textbf{Objectifs :}
\begin{itemize}
    \item Visualiser et comprendre tous les processus
    \item S'appuyer sur les flux de matériel et d'information
    \item Identifier les gaspillages
    \item Permettre une amélioration globale (non ponctuelle)
    \item Donner à l'ensemble des acteurs un langage commun
    \item Obtenir l'adhésion des acteurs travaillant sur le processus
\end{itemize}

\subsubsection{Flux physique et flux d'information}

\begin{description}[labelwidth=3.5cm]
    \item[Flux physique] Parcours du produit à travers les différentes étapes de transformation, avec les temps de cycle, les stocks intermédiaires, les délais d'attente.
    
    \item[Flux d'information] Circulation des données de planification, des ordres de fabrication, des remontées d'information entre les différents services.
\end{description}

\subsubsection{Vocabulaire : valeur ajoutée et non-valeur ajoutée}

\begin{definition}
La \textbf{valeur ajoutée (VA)} est une activité qui transforme la matière, la prestation ou l'information d'une manière qui répond aux attentes du client et pour laquelle le client est prêt à payer.
\end{definition}

\begin{definition}
La \textbf{non-valeur ajoutée (NVA)} est une activité qui demande du temps, des ressources et de l'espace sans rien apporter au produit ou service. Ces activités doivent être éliminées ou réduites.
\end{definition}

\textbf{Exemples d'activités à valeur ajoutée :}
\begin{itemize}
    \item Usinage, assemblage, transformation
    \item Contrôle qualité (si client l'exige)
    \item Montage, finition
\end{itemize}

\textbf{Exemples d'activités sans valeur ajoutée :}
\begin{itemize}
    \item Manutention, transfert
    \item Attente, stockage
    \item Contrôles multiples, reprises
    \item Déplacements inutiles
\end{itemize}

\subsubsection{Phases de la VSM}

\begin{enumerate}
    \item \textbf{Cartographier l'existant (état actuel)}
    \begin{itemize}
        \item Se rendre sur le terrain
        \item Suivre le produit du début à la fin
        \item Collecter les données : temps de cycle, temps de changement, stocks, délais
        \item Cartographier les flux d'information
        \item Calculer le Lead Time et le temps à valeur ajoutée
    \end{itemize}
    
    \item \textbf{Définir un état futur}
    \begin{itemize}
        \item Identifier les opportunités d'amélioration
        \item Éliminer les gaspillages identifiés
        \item Optimiser les flux
        \item Définir des objectifs de performance
    \end{itemize}
    
    \item \textbf{Construire un plan d'action}
    \begin{itemize}
        \item Définir les actions prioritaires
        \item Attribuer les responsabilités
        \item Fixer les échéances
        \item Suivre la mise en œuvre
    \end{itemize}
\end{enumerate}

\subsubsection{Exemple simplifié de VSM}

\begin{center}
\begin{minipage}{0.9\textwidth}
\begin{verbatim}
Fournisseur → Réception → Stock MP → Découpe → Stock inter → Assemblage → Expédition
    |           |            |          |            |           |
    v           v            v          v            v           v
  Hebdo      1/jour      2 jours     5 min       1 heure     10 min
             1 min       500 pcs    20 min       200 pcs     15 min
\end{verbatim}
\caption{Représentation simplifiée d'une cartographie VSM}
\end{minipage}
\end{center}

\subsection{Exercice pratique : Cartographie d'un processus de production}

\subsubsection{Énoncé}

Une entreprise de fabrication de meubles souhaite optimiser son processus de production de chaises. Les données suivantes ont été collectées :

\begin{table}[ht]
\centering
\begin{tabular}{|l|c|c|c|c|}
\hline
\textbf{Étape} & \textbf{Temps cycle (min)} & \textbf{Temps attente (min)} & \textbf{Stock (pièces)} & \textbf{Nb opérateurs} \\
\hline
Découpe bois & 8 & 0 & 150 & 2 \\
Assemblage piétement & 12 & 45 & 80 & 3 \\
Assemblage assise & 10 & 30 & 60 & 2 \\
Assemblage final & 15 & 60 & 40 & 2 \\
Finition & 20 & 120 & 100 & 4 \\
Emballage & 5 & 0 & 0 & 1 \\
\hline
\end{tabular}
\caption{Données du processus de production}
\end{table}

\textbf{Informations complémentaires :}
\begin{itemize}
    \item La demande client est de 500 chaises par jour (7h de travail)
    \item Le temps de changement de série moyen est de 45 minutes par opération
    \item Le taux de qualité est de 95\% pour toutes les opérations
\end{itemize}

\subsubsection{Questions}

\begin{enumerate}
    \item Calculer le temps de cycle de la ligne et identifier le goulot.
    \item Calculer le Lead Time total du processus.
    \item Calculer le temps total à valeur ajoutée.
    \item Calculer l'efficience du processus.
    \item Proposer au moins 3 axes d'amélioration.
    \item Construire une cartographie VSM simplifiée de l'état actuel.
    \item Définir un objectif pour l'état futur.
\end{enumerate}

\subsubsection{Corrigé}

\textbf{Question 1 : Temps de cycle et goulot}

Le temps de cycle de la ligne est déterminé par l'étape la plus lente :

\begin{itemize}
    \item Découpe bois : 8 min
    \item Assemblage piétement : 12 min
    \item Assemblage assise : 10 min
    \item Assemblage final : 15 min
    \item Finition : 20 min ← \textbf{Goulot}
    \item Emballage : 5 min
\end{itemize}

\textbf{Temps de cycle} = 20 min (poste finition)

Capacité théorique journalière = $7 \times 60 / 20 = 21$ cycles $\times 1$ pièce = 21 chaises/jour. L'objectif de 500 chaises/jour n'est pas atteignable avec cette configuration.

\textbf{Question 2 : Lead Time total}

\[
LT = \sum \text{Temps cycle} + \sum \text{Temps attente}
\]

\[
LT = (8+12+10+15+20+5) + (0+45+30+60+120+0) = 70 + 255 = 325 \text{ minutes}
\]

\textbf{Lead Time total} = 325 minutes (environ 5,4 heures)

\textbf{Question 3 : Temps à valeur ajoutée}

\[
\text{Temps VA} = 8 + 12 + 10 + 15 + 20 + 5 = 70 \text{ minutes}
\]

\textbf{Question 4 : Efficience du processus}

\[
\text{Efficience} = \frac{\text{Temps VA}}{LT} = \frac{70}{325} \times 100 = 21,5\%
\]

\textbf{Question 5 : Axes d'amélioration}

\begin{enumerate}
    \item \textbf{Réduire les temps d'attente} : mettre en place un flux continu (one-piece flow) pour éliminer les stocks intermédiaires
    \item \textbf{Augmenter la capacité du goulot} : ajouter une ressource au poste finition ou optimiser le temps de cycle
    \item \textbf{Réduire les temps de changement} : mettre en place la méthode SMED (Single Minute Exchange of Die)
    \item \textbf{Améliorer la qualité} : réduire le taux de rebut (actuellement 5\%) par des contrôles à la source
    \item \textbf{Réorganiser le layout} : rapprocher les postes pour réduire les distances de transfert
\end{enumerate}

\textbf{Question 6 : Cartographie VSM de l'état actuel}

\begin{center}
\begin{minipage}{0.95\textwidth}
\begin{verbatim}
Flux d'information :
Planification → Ordres de fabrication (journaliers)

Flux matière :
Fournisseur → Découpe (8 min) → Stock 150 → Ass. piétement (12 min) → Stock 80 
→ Ass. assise (10 min) → Stock 60 → Ass. final (15 min) → Stock 40 
→ Finition (20 min) → Stock 100 → Emballage (5 min) → Client

Temps de cycle : 70 min
Lead Time : 325 min
Efficience : 21,5%
\end{verbatim}
\caption{Cartographie VSM - État actuel}
\end{minipage}
\end{center}

\textbf{Question 7 : Objectif pour l'état futur}

\begin{itemize}
    \item \textbf{Lead Time} : réduire de 325 min à 90 min (objectif -72\%)
    \item \textbf{Stock intermédiaire} : réduire à 0 (flux continu)
    \item \textbf{Efficience} : augmenter de 21,5\% à 70\%
    \item \textbf{Temps de cycle} : réduire le goulot de 20 min à 12 min
    \item \textbf{Taux de qualité} : augmenter de 95\% à 98\%
\end{itemize}

\subsection{Synthèse}

\begin{center}
\fbox{\begin{minipage}{0.9\textwidth}
\textbf{Points clés à retenir :}

\begin{itemize}
    \item L'analyse de déroulement utilise des symboles normalisés (opération, transfert, contrôle, attente, stockage)
    \item Les indicateurs clés sont l'efficacité, l'efficience, le Lead Time et l'indice de tension
    \item Le graphique de déroulement peut être décliné selon trois perspectives : exécutant, matière, matériel
    \item La matrice de déroulement structure le recueil des données terrain
    \item Le VSM cartographie les flux physiques et d'information pour identifier les gaspillages
    \item La VSM se déroule en trois phases : état actuel, état futur, plan d'action
    \item La distinction valeur ajoutée / non-valeur ajoutée est fondamentale pour l'optimisation
\end{itemize}
\end{minipage}}
\end{center}

% ==================== FIN DE LA PARTIE 11 ====================
% ==================== PARTIE 6 : TEMPS PRÉDÉTERMINÉS (APPROFONDISSEMENT) ====================
\part*{Partie 6 : Temps Prédéterminés (Approfondissement)}
\addcontentsline{toc}{part}{Partie 6 : Temps Prédéterminés (Approfondissement)}

\section{Méthode MTM (Methods-Time Measurement)}

La méthode MTM est l'une des techniques de temps prédéterminés les plus connues et les plus utilisées dans l'industrie. Elle permet de déterminer le temps nécessaire à l'exécution d'une tâche en analysant ses mouvements élémentaires.

\subsection{Définition et origine}

\begin{definition}
Le \textbf{MTM (Methods-Time Measurement)} est un système de temps prédéterminés qui analyse tous les mouvements élémentaires qu'une opération manuelle demande pour sa réalisation et attribue à chaque mouvement un temps de base prédéterminé à l'aide de tables standardisées.
\end{definition}

\textbf{Origine :}
\begin{itemize}
    \item Développé en 1948 par H.B. Maynard, G.J. Stegemerten et J.L. Schwab
    \item Basé sur les travaux de Frank et Lillian Gilbreth (therbligs)
    \item Première méthode structurée de temps prédéterminés
    \item Aujourd'hui, le MTM est un standard international reconnu
\end{itemize}

\begin{center}
\fbox{\begin{minipage}{0.85\textwidth}
\centering
\textbf{Principe fondamental du MTM} \\
Le MTM élimine la subjectivité du jugement d'allure car les temps sont donnés par des tables. \\
Le temps de base obtenu correspond à celui d'un opérateur qualifié travaillant à allure normale.
\end{minipage}}
\end{center}

\subsection{L'unité de mesure : le TMU (Time Measurement Unit)}

\begin{definition}
Le \textbf{TMU (Time Measurement Unit)} est l'unité de base du MTM. Il permet d'exprimer les temps élémentaires avec une grande précision.
\end{definition}

\textbf{Équivalences du TMU :}
\[
1\ \text{TMU} = 0,00001\ \text{heure} = 0,0006\ \text{minute} = 0,036\ \text{seconde} = 1\ \text{cmh}
\]

\begin{table}[ht]
\centering
\begin{tabular}{|c|c|c|c|}
\hline
\textbf{Unité} & \textbf{TMU} & \textbf{Minutes} & \textbf{Secondes} \\
\hline
1 heure & 100 000 & 60 & 3 600 \\
1 minute & 1 667 & 1 & 60 \\
1 seconde & 27,8 & 0,0167 & 1 \\
\hline
\end{tabular}
\caption{Équivalences du TMU}
\end{table}

\subsection{Les mouvements de base}

Le MTM décompose le travail en mouvements élémentaires codifiés. Chaque mouvement est identifié par un symbole et a un temps prédéterminé selon ses caractéristiques.

\subsubsection{Mouvements des membres supérieurs}

\begin{table}[ht]
\centering
\begin{tabular}{|c|c|c|}
\hline
\textbf{Symbole} & \textbf{Mouvement} & \textbf{Description} \\
\hline
R & Atteindre (Reach) & Déplacer la main vers une destination \\
M & Mouvoir (Move) & Déplacer un objet \\
T & Tourner (Turn) & Rotation de la main ou du poignet \\
AP & Appliquer pression (Apply Pressure) & Appliquer une force \\
G & Saisir (Grasp) & Prendre un objet \\
P & Positionner (Position) & Orienter un objet avec précision \\
RL & Relâcher (Release) & Lâcher un objet \\
D & Désengager (Disengage) & Séparer deux objets assemblés \\
C & Mouvement de manivelle (Crank) & Mouvement circulaire continu \\
\hline
\end{tabular}
\caption{Mouvements des membres supérieurs en MTM}
\end{table}

\subsubsection{Mouvements des membres inférieurs}

\begin{table}[ht]
\centering
\begin{tabular}{|c|c|c|}
\hline
\textbf{Symbole} & \textbf{Mouvement} & \textbf{Description} \\
\hline
LM & Leg Motion & Mouvement de jambe \\
FM & Feet Motion & Mouvement de pied \\
\hline
\end{tabular}
\caption{Mouvements des membres inférieurs}
\end{table}

\subsubsection{Mouvements du corps}

\begin{table}[ht]
\centering
\begin{tabular}{|c|c|c|}
\hline
\textbf{Symbole} & \textbf{Mouvement} & \textbf{Description} \\
\hline
W & Walk & Marche \\
SS & Side Step & Pas de côté \\
TB & Turn Body & Rotation du corps \\
B & Bend & S'incliner (pencher le buste) \\
S & Stoop & Se baisser (genoux fléchis) \\
KOK & Kneel on 1 Knee & S'agenouiller sur un genou \\
KBK & Kneel on Both Knees & S'agenouiller sur deux genoux \\
AB & Arise from Bend & Se relever après s'être incliné \\
AS & Arise from Stoop & Se relever après s'être baissé \\
AKOK & Arise from 1 Knee & Se relever après un genou \\
AKBK & Arise from Both Knees & Se relever après deux genoux \\
SIT & Sit & S'asseoir \\
STD & Stand & Se lever \\
\hline
\end{tabular}
\caption{Mouvements du corps}
\end{table}

\subsubsection{Mouvements visuels}

\begin{table}[ht]
\centering
\begin{tabular}{|c|c|c|}
\hline
\textbf{Symbole} & \textbf{Mouvement} & \textbf{Description} \\
\hline
ET & Eye Travel & Déplacement du regard \\
EF & Eye Focus & Focalisation (mise au point) \\
\hline
\end{tabular}
\caption{Mouvements visuels}
\end{table}

\subsection{Les cas du mouvement "Atteindre" (R)}

Le mouvement "Atteindre" est l'un des plus complexes du MTM. Il se décline en 5 cas selon la situation.

\subsubsection{Description des cas}

\begin{description}[labelwidth=1.5cm]
    \item[R-A] Atteindre un objet toujours placé au même endroit, un objet dans l'autre main, ou un objet sur lequel l'autre main repose.
    
    \item[R-B] Atteindre un objet isolé dont l'emplacement peut varier légèrement d'un cycle à l'autre.
    
    \item[R-C] Atteindre un objet mêlé à d'autres de telle sorte qu'il y ait recherche et sélection.
    
    \item[R-D] Atteindre un objet très petit ou un objet à saisir avec précision ou précaution.
    
    \item[R-E] Atteindre avec la main un endroit peu précis pour assurer l'équilibre du corps, pour évacuer la zone de travail ou pour préparer le mouvement suivant.
\end{description}

\subsection{Les types de mouvements}

Selon l'état de la main au début et à la fin du mouvement, on distingue trois types :

\begin{description}[labelwidth=1.5cm]
    \item[Type I] Mouvement normal : main immobile au début et à la fin du geste.
    \begin{itemize}
        \item Notation : Exemple $R20B$
    \end{itemize}
    
    \item[Type II] Mouvement avec main en mouvement au début ou à la fin.
    \begin{itemize}
        \item Notation : $mR20B$ (main en mouvement au début) ou $R20Bm$ (main en mouvement à la fin)
    \end{itemize}
    
    \item[Type III] Mouvement avec main en mouvement au début et à la fin.
    \begin{itemize}
        \item Notation : $mR20Bm$
    \end{itemize}
\end{description}

\subsection{Extrait des tables MTM}

\subsubsection{Table du mouvement "Atteindre" (R)}

\begin{table}[ht]
\centering
\begin{tabular}{|c|c|c|c|c|c|c|c|}
\hline
\textbf{Distance (cm)} & \textbf{R-A} & \textbf{R-B} & \textbf{R-C} & \textbf{R-D} & \textbf{R-E} & \textbf{mR-A} & \textbf{R-A m} \\
\hline
2 & 2,0 & 2,0 & 2,0 & 2,0 & 1,6 & 1,6 & 1,4 \\
4 & 3,3 & 3,3 & 5,2 & 3,3 & 3,0 & 2,5 & 0,8 \\
6 & 4,5 & 4,5 & 6,5 & 4,5 & 3,9 & 3,0 & 1,5 \\
8 & 5,4 & 5,6 & 7,5 & 7,5 & 5,5 & 3,6 & 2,0 \\
10 & 6,0 & 6,6 & 8,4 & 6,4 & 4,9 & 4,2 & 2,4 \\
12 & 6,4 & 7,4 & 9,1 & 7,1 & 5,2 & 4,8 & 2,6 \\
14 & 6,7 & 8,2 & 9,7 & 7,7 & 5,5 & 5,3 & 2,9 \\
16 & 7,1 & 8,8 & 10,3 & 8,2 & 5,8 & 5,9 & 2,9 \\
18 & 7,4 & 9,4 & 10,8 & 8,7 & 6,1 & 6,5 & 2,9 \\
20 & 7,8 & 9,9 & 11,4 & 9,2 & 6,4 & 7,1 & 2,8 \\
\hline
\end{tabular}
\caption{Extrait de la table MTM pour le mouvement "Atteindre" (valeurs en TMU)}
\caption*{Note : Les valeurs sont données à titre indicatif. Les tables complètes sont plus détaillées.}
\end{table}

\subsubsection{Table du mouvement "Mouvoir" (M)}

\begin{table}[ht]
\centering
\begin{tabular}{|c|c|c|c|}
\hline
\textbf{Distance (cm)} & \textbf{M-A} & \textbf{M-B} & \textbf{M-C} \\
\hline
2 & 2,0 & 2,0 & 2,0 \\
4 & 3,4 & 3,4 & 3,4 \\
6 & 4,5 & 4,5 & 4,8 \\
8 & 5,5 & 5,6 & 6,1 \\
10 & 6,1 & 6,4 & 7,0 \\
15 & 7,9 & 8,6 & 9,6 \\
20 & 9,4 & 10,4 & 11,7 \\
\hline
\end{tabular}
\caption{Extrait de la table MTM pour le mouvement "Mouvoir" (valeurs en TMU)}
\end{table}

\subsection{Avantages et faiblesses du MTM}

\subsubsection{Avantages}

\begin{itemize}
    \item \textbf{Standards homogènes} : mêmes temps pour tous les utilisateurs
    \item \textbf{Pas de jugement d'allure} : élimination de la subjectivité
    \item \textbf{Reproductibilité} : résultats identiques quel que soit l'analyste
    \item \textbf{Anticipation} : calcul possible avant la création du poste
    \item \textbf{Mise à jour facile} : modifications aisées des modes opératoires
    \item \textbf{Acceptation sociale} : méthode neutre et objective
    \item \textbf{Reconnu internationalement} : langage commun entre entreprises
\end{itemize}

\subsubsection{Faiblesses}

\begin{itemize}
    \item \textbf{Complexité} : formation longue (environ 20 jours)
    \item \textbf{Temps d'analyse} : une journée pour analyser 1 minute de cycle
    \item \textbf{Risque d'erreur} : élevé pour un utilisateur occasionnel
    \item \textbf{Coût} : formation et maintenance importantes
    \item \textbf{Nécessité d'analystes certifiés}
    \item \textbf{Précision} : peut être excessive pour certaines applications
\end{itemize}

\begin{center}
\fbox{\begin{minipage}{0.85\textwidth}
\centering
\textbf{Évolution du MTM} \\
Face à ces limites, le MTM a évolué vers des versions simplifiées : \\
MTM-2 (1965), MTM-3, MTM-UAS, MTM-MEK, MTM-SAM, \\
et surtout le MOST (Maynard Operation Sequence Technique) qui permet \\
de chiffrer une heure de cycle en une journée d'étude.
\end{minipage}}
\end{center}

\subsection{Exemples de calcul détaillés}

\subsubsection{Exemple 1 : Saisie d'une pièce sur une table}

\begin{exemple}
\textbf{Description :} L'opérateur atteint une pièce posée à 20 cm sur sa droite, la saisit, puis la dépose sur un poste de travail.

\textbf{Décomposition MTM :}
\begin{enumerate}
    \item Atteindre la pièce : R20B (distance 20 cm, cas B)
    \item Saisir la pièce : G1 (saisie simple)
    \item Mouvoir la pièce : M20B (déplacement de 20 cm)
    \item Relâcher la pièce : RL1 (lâcher simple)
\end{enumerate}

\textbf{Calcul des temps :}
\begin{itemize}
    \item R20B : 9,9 TMU
    \item G1 : 2,0 TMU
    \item M20B : 10,4 TMU
    \item RL1 : 2,0 TMU
\end{itemize}

\textbf{Temps total :} $9,9 + 2,0 + 10,4 + 2,0 = 24,3$ TMU

Conversion en secondes : $24,3 \times 0,036 = 0,875$ secondes
\end{exemple}

\subsubsection{Exemple 2 : Positionnement d'une pièce avec précision}

\begin{exemple}
\textbf{Description :} L'opérateur saisit une pièce et la positionne avec précision dans un emplacement spécifique.

\textbf{Décomposition MTM :}
\begin{enumerate}
    \item Atteindre la pièce à 15 cm : R15C (objet mêlé à d'autres)
    \item Saisir la pièce : G3 (saisie avec sélection)
    \item Mouvoir la pièce : M15C (déplacement avec précision)
    \item Positionner : P2SE (positionnement avec symétrie)
\end{enumerate}

\textbf{Calcul des temps :}
\begin{itemize}
    \item R15C : 9,4 TMU
    \item G3 : 5,6 TMU
    \item M15C : 9,6 TMU
    \item P2SE : 16,2 TMU
\end{itemize}

\textbf{Temps total :} $9,4 + 5,6 + 9,6 + 16,2 = 40,8$ TMU

Conversion en secondes : $40,8 \times 0,036 = 1,469$ secondes
\end{exemple}

\subsection{Exercices d'application MTM}

\subsubsection{Exercice 1 : Saisie d'un outil}

Un opérateur doit saisir un tournevis situé à 25 cm sur sa gauche, dans un support où plusieurs outils sont rangés (cas C). Il doit saisir l'outil avec une prise précise (G3). Calculez le temps de cette séquence.

\subsubsection{Exercice 2 : Assemblage de deux pièces}

Un opérateur assemble deux pièces. Il atteint la pièce A à 30 cm (cas A), la saisit (G1), l'atteint la pièce B à 20 cm (cas A), les assemble avec un mouvement de 15 cm (M15A), puis positionne avec précision (P2SE). Calculez le temps total.

\subsubsection{Exercice 3 : Calcul complet d'assemblage}

\begin{table}[ht]
\centering
\begin{tabular}{|l|c|}
\hline
\textbf{Mouvement} & \textbf{Valeur TMU} \\
\hline
Atteindre pièce (R20A) & ? \\
Saisir pièce (G1) & 2,0 \\
Mouvoir pièce (M25B) & 11,7 \\
Positionner (P1NS) & 5,6 \\
Relâcher (RL1) & 2,0 \\
\hline
\end{tabular}
\caption{Données pour l'exercice 3}
\end{table}

Complétez la valeur manquante dans le tableau en utilisant la table MTM fournie plus haut. Calculez le temps total en TMU puis en secondes.

\subsubsection{Corrigé de l'exercice 1}

\begin{itemize}
    \item R25C : distance 25 cm, cas C (objet mêlé à d'autres)
    \item D'après la table : R25C = 11,9 TMU (valeur extrapolée)
    \item G3 = 5,6 TMU
\end{itemize}

\textbf{Temps total :} $11,9 + 5,6 = 17,5$ TMU

Conversion : $17,5 \times 0,036 = 0,63$ seconde

\subsubsection{Corrigé de l'exercice 2}

\begin{itemize}
    \item R30A = 9,5 TMU (distance 30 cm, cas A)
    \item G1 = 2,0 TMU
    \item R20A = 7,8 TMU
    \item M15A = 7,9 TMU
    \item P2SE = 16,2 TMU
\end{itemize}

\textbf{Temps total :} $9,5 + 2,0 + 7,8 + 7,9 + 16,2 = 43,4$ TMU

Conversion : $43,4 \times 0,036 = 1,562$ secondes

\subsubsection{Corrigé de l'exercice 3}

\begin{itemize}
    \item R20A = 7,8 TMU (d'après la table)
    \item G1 = 2,0 TMU
    \item M25B = 11,7 TMU (donné)
    \item P1NS = 5,6 TMU (donné)
    \item RL1 = 2,0 TMU (donné)
\end{itemize}

\textbf{Temps total :} $7,8 + 2,0 + 11,7 + 5,6 + 2,0 = 29,1$ TMU

Conversion : $29,1 \times 0,036 = 1,048$ secondes

\subsection{Synthèse}

\begin{center}
\fbox{\begin{minipage}{0.9\textwidth}
\textbf{Points clés à retenir :}

\begin{itemize}
    \item Le MTM est une méthode de temps prédéterminés basée sur l'analyse des mouvements élémentaires
    \item L'unité de mesure est le TMU ($1\ \text{TMU} = 0,036\ \text{s}$)
    \item Les mouvements sont codifiés par des symboles : R (atteindre), M (mouvoir), G (saisir), etc.
    \item Le mouvement "Atteindre" a 5 cas selon la situation (A, B, C, D, E)
    \item Les types de mouvements (I, II, III) dépendent de l'état de la main
    \item Le MTM élimine la subjectivité mais est complexe à mettre en œuvre
    \item Le MOST est une évolution simplifiée du MTM
\end{itemize}
\end{minipage}}
\end{center}

% ==================== FIN DE LA PARTIE 13 ====================
% ==================== PARTIE 14 : MÉTHODE MOST ====================

\section{Méthode MOST (Maynard Operation Sequence Technique)}

Le MOST est une évolution du MTM développée pour répondre aux besoins de rapidité et de simplicité des entreprises industrielles. Il permet de chiffrer des temps de cycle beaucoup plus rapidement que le MTM traditionnel.

\subsection{Définition et origine}

\begin{definition}
Le \textbf{MOST (Maynard Operation Sequence Technique)} est un système de temps prédéterminés qui utilise des séquences standardisées de mouvements pour décrire et chiffrer les opérations manuelles. Il regroupe les mouvements élémentaires en blocs fonctionnels.
\end{definition}

\textbf{Origine :}
\begin{itemize}
    \item Développé en 1974 par Kjell Zandin (Suédois) pour le cabinet H.B. Maynard
    \item Premier utilisateur officiel en France : Michelin (années 1980)
    \item Conçu pour répondre aux limites du MTM (complexité, lenteur)
    \item Basé sur des regroupements statistiques de mouvements
\end{itemize}

\begin{center}
\fbox{\begin{minipage}{0.85\textwidth}
\centering
\textbf{Principe fondamental du MOST} \\
Le MOST regroupe les mouvements élémentaires en séquences types, \\
ce qui permet de réduire considérablement le temps d'analyse : \\
une heure de cycle peut être chiffrée en une journée d'étude.
\end{minipage}}
\end{center}

\subsection{Les trois types de séquences}

Le MOST utilise trois types de séquences standardisées pour décrire les opérations.

\subsubsection{Séquence 1 : Déplacement d'objets (ABG ABP A)}

C'est la séquence la plus utilisée. Elle décrit le déplacement d'un objet d'un endroit à un autre.

\begin{center}
\begin{tabular}{|c|c|c|c|c|c|c|}
\hline
\textbf{A} & \textbf{B} & \textbf{G} & \textbf{A} & \textbf{B} & \textbf{P} & \textbf{A} \\
\hline
Action & Body & Get & Action & Body & Place & Action \\
\hline
\end{tabular}
\caption*{Structure de la séquence 1 : ABG ABP A}
\end{center}

\textbf{Signification des paramètres :}
\begin{description}[labelwidth=0.8cm]
    \item[A] Action générale (déplacement, mouvement)
    \item[B] Body (mouvement du corps, marche)
    \item[G] Get (saisie, prise de l'objet)
    \item[P] Place (positionnement, dépôt de l'objet)
\end{description}

\subsubsection{Séquence 2 : Mouvement guidé}

Cette séquence décrit les mouvements où l'objet est contrôlé ou guidé pendant le déplacement.

\begin{center}
\begin{tabular}{|c|c|c|c|c|c|c|}
\hline
\textbf{A} & \textbf{B} & \textbf{G} & \textbf{A} & \textbf{B} & \textbf{P} & \textbf{A} \\
\hline
Action & Body & Get & Action & Body & Place & Action \\
\hline
\end{tabular}
\caption*{Structure de la séquence 2 (identique mais indices spécifiques)}
\end{center}

\subsubsection{Séquence 3 : Utilisation d'outil à main}

Cette séquence décrit l'utilisation d'un outil manuel.

\begin{center}
\begin{tabular}{|c|c|c|c|c|c|c|c|}
\hline
\textbf{A} & \textbf{B} & \textbf{G} & \textbf{A} & \textbf{B} & \textbf{P} & \textbf{A} & \textbf{B} & \textbf{G} & \textbf{A} \\
\hline
Action & Body & Get & Action & Body & Place & Action & Body & Get & Action \\
\hline
\end{tabular}
\caption*{Structure de la séquence 3 : ABG ABP A B G A}
\end{center}

\subsection{Les indices et valeurs clés}

Le MOST utilise des indices pour chaque paramètre. La somme des indices est multipliée par 10 pour obtenir le temps en TMU.

\subsubsection{Valeurs clés (Basic MOST)}

\begin{table}[ht]
\centering
\begin{tabular}{|c|c|c|}
\hline
\textbf{Indice} & \textbf{Valeur (TMU)} & \textbf{Signification} \\
\hline
0 & 0 & Aucun mouvement \\
1 & 10 & Mouvement très court, à portée de main \\
3 & 30 & Mouvement limité \\
6 & 60 & Mouvement étendu \\
10 & 100 & Déplacement plus long \\
16 & 160 & Déplacement très long \\
\hline
\end{tabular}
\caption{Valeurs clés du Basic MOST}
\end{table}

\begin{center}
\fbox{\begin{minipage}{0.85\textwidth}
\centering
\textbf{Remarque importante} \\
Une seule valeur de 10 TMU dans le MOST remplace environ 136 valeurs de la table MTM1. \\
Les valeurs sont très facilement mémorisables.
\end{minipage}}
\end{center}

\subsection{Table MOST simplifiée}

\subsubsection{Indices pour le paramètre A (Action)}

\begin{table}[ht]
\centering
\begin{tabular}{|c|c|c|}
\hline
\textbf{Indice} & \textbf{Description} & \textbf{Exemple} \\
\hline
0 & Pas d'action & Objet déjà en main \\
1 & Mouvement vers l'objet & Atteindre objet à portée \\
3 & Mouvement limité & Atteindre objet à 1-2 pas \\
6 & Mouvement étendu & Atteindre objet à 2-3 pas \\
10 & Déplacement long & Atteindre objet à plus de 3 pas \\
\hline
\end{tabular}
\caption{Indices pour le paramètre A}
\end{table}

\subsubsection{Indices pour le paramètre B (Body)}

\begin{table}[ht]
\centering
\begin{tabular}{|c|c|c|}
\hline
\textbf{Indice} & \textbf{Description} & \textbf{Exemple} \\
\hline
0 & Stationnaire & Pas de mouvement du corps \\
1 & Mouvement limité & Se pencher, tourner \\
3 & Déplacement limité & 1-2 pas \\
6 & Déplacement étendu & 2-3 pas \\
10 & Déplacement long & Plus de 3 pas \\
\hline
\end{tabular}
\caption{Indices pour le paramètre B}
\end{table}

\subsubsection{Indices pour le paramètre G (Get)}

\begin{table}[ht]
\centering
\begin{tabular}{|c|c|c|}
\hline
\textbf{Indice} & \textbf{Description} & \textbf{Exemple} \\
\hline
0 & Pas de saisie & Objet déjà saisi \\
1 & Saisie simple & Saisie d'un objet isolé \\
3 & Saisie avec contrôle & Saisie d'un objet mêlé \\
6 & Saisie avec précision & Saisie d'un petit objet \\
\hline
\end{tabular}
\caption{Indices pour le paramètre G}
\end{table}

\subsubsection{Indices pour le paramètre P (Place)}

\begin{table}[ht]
\centering
\begin{tabular}{|c|c|c|}
\hline
\textbf{Indice} & \textbf{Description} & \textbf{Exemple} \\
\hline
0 & Pas de placement & Objet relâché librement \\
1 & Placement simple & Dépose simple \\
3 & Placement avec ajustement & Positionnement approximatif \\
6 & Placement avec précision & Positionnement précis \\
\hline
\end{tabular}
\caption{Indices pour le paramètre P}
\end{table}

\subsection{Exemples de calcul}

\subsubsection{Exemple 1 : Saisie d'un objet à portée de main}

\begin{exemple}
\textbf{Description :} L'opérateur saisit un marqueur posé sur la table à portée de main et le dépose sur une tablette.

\textbf{Décomposition :}
\begin{itemize}
    \item A : Mouvement vers l'objet (indice 1)
    \item B : Pas de mouvement du corps (indice 0)
    \item G : Saisie simple (indice 1)
    \item A : Mouvement vers la destination (indice 1)
    \item B : Pas de mouvement du corps (indice 0)
    \item P : Placement simple (indice 1)
    \item A : Retour (indice 0)
\end{itemize}

\textbf{Calcul :}
\[
\text{Somme des indices} = 1 + 0 + 1 + 1 + 0 + 1 + 0 = 4
\]
\[
\text{Temps} = 4 \times 10 = 40\ \text{TMU}
\]
\[
\text{Conversion} = 40 \times 0,036 = 1,44\ \text{secondes}
\]
\end{exemple}

\subsubsection{Exemple 2 : Saisie d'un objet à distance}

\begin{exemple}
\textbf{Description :} L'opérateur saisit un marqueur situé à trois pas au sol et le dépose sur une tablette.

\textbf{Décomposition :}
\begin{itemize}
    \item A : Mouvement vers l'objet (indice 6 - étendu)
    \item B : Déplacement (indice 6)
    \item G : Saisie simple (indice 1)
    \item A : Mouvement vers la destination (indice 6)
    \item B : Déplacement retour (indice 0 - pas de déplacement)
    \item P : Placement simple (indice 1)
    \item A : Retour (indice 0)
\end{itemize}

\textbf{Calcul :}
\[
\text{Somme des indices} = 6 + 6 + 1 + 6 + 0 + 1 + 0 = 20
\]
\[
\text{Temps} = 20 \times 10 = 200\ \text{TMU}
\]
\[
\text{Conversion} = 200 \times 0,036 = 7,2\ \text{secondes}
\]
\end{exemple}

\subsubsection{Exemple 3 : Positionnement précis}

\begin{exemple}
\textbf{Description :} L'opérateur saisit une pièce, se déplace de 2 pas, et la positionne avec précision dans un emplacement.

\textbf{Décomposition :}
\begin{itemize}
    \item A : Mouvement vers l'objet (indice 3)
    \item B : Pas de mouvement (indice 0)
    \item G : Saisie simple (indice 1)
    \item A : Mouvement vers la destination (indice 3)
    \item B : Déplacement (indice 3)
    \item P : Placement avec précision (indice 6)
    \item A : Retour (indice 0)
\end{itemize}

\textbf{Calcul :}
\[
\text{Somme des indices} = 3 + 0 + 1 + 3 + 3 + 6 + 0 = 16
\]
\[
\text{Temps} = 16 \times 10 = 160\ \text{TMU}
\]
\[
\text{Conversion} = 160 \times 0,036 = 5,76\ \text{secondes}
\]
\end{exemple}

\subsection{Comparaison MTM vs MOST}

\begin{table}[ht]
\centering
\begin{tabular}{|p{4cm}|p{4.5cm}|p{4.5cm}|}
\hline
\textbf{Caractéristique} & \textbf{MTM} & \textbf{MOST} \\
\hline
Année de création & 1948 & 1974 \\
Unité de mesure & TMU & TMU (x10) \\
Niveau de détail & Très fin (mouvements élémentaires) & Aggloméré (séquences types) \\
Temps d'analyse & 1 jour pour 1 minute de cycle & 1 jour pour 1 heure de cycle \\
Formation requise & 20 jours & 6 jours \\
Nombre de valeurs & > 200 & Environ 30 \\
Complexité & Élevée & Modérée \\
Précision & Très élevée & Élevée (statistique) \\
Utilisation & Études très fines & Études courantes \\
\hline
\end{tabular}
\caption{Comparaison MTM vs MOST}
\end{table}

\begin{center}
\fbox{\begin{minipage}{0.85\textwidth}
\centering
\textbf{Quand utiliser quel outil ?}
\begin{itemize}
    \item MTM : pour les analyses très fines, les postes critiques, les temps de cycle très courts
    \item MOST : pour les études courantes, les nouvelles gammes, les temps de cycle longs
\end{itemize}
\end{minipage}}
\end{center}

\subsection{Avantages et limites}

\subsubsection{Avantages du MOST}

\begin{itemize}
    \item \textbf{Rapidité} : une heure de cycle chiffrée en une journée d'étude
    \item \textbf{Simplicité} : valeurs facilement mémorisables (10, 30, 60, 100, 160)
    \item \textbf{Formation courte} : 6 jours maximum
    \item \textbf{Fiabilité} : précision maîtrisée par regroupements statistiques
    \item \textbf{Reproductibilité} : résultats identiques entre analystes
    \item \textbf{Acceptation} : méthode neutre et objective
    \item \textbf{Évolutivité} : versions adaptées (Basic MOST, Mini MOST, Maxi MOST)
\end{itemize}

\subsubsection{Limites du MOST}

\begin{itemize}
    \item \textbf{Précision moindre} que le MTM pour les cycles très courts
    \item \textbf{Nécessite une formation} (bien que plus courte que MTM)
    \item \textbf{Non adapté} aux mouvements très complexes
    \item \textbf{Nécessite des analystes certifiés} pour une utilisation fiable
\end{itemize}

\subsection{Exercices d'application MOST}

\subsubsection{Exercice 1 : Saisie d'un outil à portée}

Un opérateur saisit un tournevis posé sur la table à portée de main (indice A=1, B=0, G=1), le dépose sur un support à portée (A=1, B=0, P=1). Calculez le temps total.

\subsubsection{Exercice 2 : Déplacement avec marche}

Un opérateur doit se déplacer de 3 pas pour saisir une pièce (A=6, B=6, G=1), revenir à son poste (A=6, B=6, P=0), et déposer la pièce. Calculez le temps total.

\subsubsection{Exercice 3 : Positionnement précis}

Un opérateur saisit une pièce à portée (A=1, B=0, G=1), se déplace d'un pas (A=3, B=3, P=6) et la positionne avec précision. Calculez le temps total.

\subsubsection{Exercice 4 : Séquence complète}

\begin{table}[ht]
\centering
\begin{tabular}{|c|c|c|c|c|c|c|}
\hline
\textbf{A} & \textbf{B} & \textbf{G} & \textbf{A} & \textbf{B} & \textbf{P} & \textbf{A} \\
\hline
3 & 0 & 1 & 6 & 3 & 3 & 0 \\
\hline
\end{tabular}
\caption{Séquence pour l'exercice 4}
\end{table}

Calculez le temps total en TMU puis en secondes pour cette séquence.

\subsubsection{Exercice 5 : Cas industriel}

Dans un atelier de montage, un opérateur doit :
\begin{enumerate}
    \item Saisir une pièce située à 2 pas (indice B=3)
    \item La saisir (G=1)
    \item Revenir à son poste (B=3)
    \item La positionner avec précision (P=6)
\end{enumerate}

Calculez le temps total en TMU et en secondes.

\subsection{Corrigés des exercices}

\subsubsection{Corrigé exercice 1}

\[
\text{Séquence : A=1, B=0, G=1, A=1, B=0, P=1, A=0
\]
\[
\text{Somme} = 1 + 0 + 1 + 1 + 0 + 1 + 0 = 4
\]
\[
\text{Temps} = 4 \times 10 = 40\ \text{TMU}
\]
\[
40 \times 0,036 = 1,44\ \text{secondes}
\]

\subsubsection{Corrigé exercice 2}

\[
\text{Séquence : A=6, B=6, G=1, A=6, B=6, P=0, A=0
\]
\[
\text{Somme} = 6 + 6 + 1 + 6 + 6 + 0 + 0 = 25
\]
\[
\text{Temps} = 25 \times 10 = 250\ \text{TMU}
\]
\[
250 \times 0,036 = 9,0\ \text{secondes}
\]

\subsubsection{Corrigé exercice 3}

\[
\text{Séquence : A=1, B=0, G=1, A=3, B=3, P=6, A=0
\]
\[
\text{Somme} = 1 + 0 + 1 + 3 + 3 + 6 + 0 = 14
\]
\[
\text{Temps} = 14 \times 10 = 140\ \text{TMU}
\]
\[
140 \times 0,036 = 5,04\ \text{secondes}
\]

\subsubsection{Corrigé exercice 4}

\[
\text{Somme} = 3 + 0 + 1 + 6 + 3 + 3 + 0 = 16
\]
\[
\text{Temps} = 16 \times 10 = 160\ \text{TMU}
\]
\[
160 \times 0,036 = 5,76\ \text{secondes}
\]

\subsubsection{Corrigé exercice 5}

\textbf{Décomposition :}
\begin{itemize}
    \item A : Mouvement vers l'objet (indice 3)
    \item B : Déplacement (indice 3)
    \item G : Saisie (indice 1)
    \item A : Mouvement vers destination (indice 3)
    \item B : Déplacement retour (indice 3)
    \item P : Positionnement précis (indice 6)
    \item A : Retour (indice 0)
\end{itemize}

\[
\text{Somme} = 3 + 3 + 1 + 3 + 3 + 6 + 0 = 19
\]
\[
\text{Temps} = 19 \times 10 = 190\ \text{TMU}
\]
\[
190 \times 0,036 = 6,84\ \text{secondes}
\]

\subsection{Synthèse}

\begin{center}
\fbox{\begin{minipage}{0.9\textwidth}
\textbf{Points clés à retenir :}

\begin{itemize}
    \item Le MOST est une évolution du MTM développée par Kjell Zandin en 1974
    \item Il utilise trois types de séquences : déplacement d'objets (ABG ABP A), mouvement guidé, utilisation d'outil
    \item Les indices sont des valeurs clés : 0, 1, 3, 6, 10, 16 (multipliées par 10 pour obtenir les TMU)
    \item Le temps est calculé par : $\text{Somme des indices} \times 10 = \text{TMU}$
    \item Le MOST est beaucoup plus rapide que le MTM (1 heure de cycle par jour d'étude)
    \item La formation MOST est de 6 jours contre 20 jours pour le MTM
    \item Le MOST est adapté aux études courantes, le MTM aux analyses très fines
\end{itemize}
\end{minipage}}
\end{center}

% ==================== FIN DE LA PARTIE 14 ====================
% ==================== PARTIE 7 : EXERCICES ET CAS PRATIQUES ====================
\part*{Partie 7 : Exercices et Cas Pratiques}
\addcontentsline{toc}{part}{Partie 7 : Exercices et Cas Pratiques}

\section{Exercices de Conversion d'Unités}

\subsection{Conversions de base}

\subsubsection*{Énoncé}

Convertir les durées suivantes dans l'unité demandée :

\begin{enumerate}
    \item $2,5$ heures en centièmes d'heure (ch)
    \item $45$ minutes en décimilliheures (dmh)
    \item $3\,600$ secondes en centiminutes (cmin)
    \item $150$ dmh en heures
    \item $8$ minutes en centièmes de minute (cmin)
    \item $0,75$ heure en secondes
    \item $12\,500$ cmh en heures
    \item $1$ heure $15$ minutes $30$ secondes en centièmes d'heure
\end{enumerate}

\subsection{Application industrielle}

\subsubsection*{Énoncé}

Une entreprise utilise des temps exprimés en h:min:s pour ses relevés vidéo, mais sa GPAO travaille en centièmes d'heure (ch). Convertir les temps suivants :

\begin{enumerate}
    \item $0:15:30$ ($0$h $15$min $30$s) en ch
    \item $1:45:20$ en ch
    \item $2:08:45$ en ch
    \item $0:05:15$ en ch
    \item $3:30:00$ en ch
\end{enumerate}

\subsection{Choix de l'unité adaptée}

\subsubsection*{Énoncé}

Pour chaque situation, proposez l'unité de temps la plus adaptée et justifiez votre choix :

\begin{enumerate}
    \item Chronométrage d'une opération manuelle de $3$ secondes
    \item Calcul de charge hebdomadaire pour $50$ opérateurs
    \item Analyse de temps machine pour un cycle de $2$ minutes
    \item Temps de traversée d'un produit sur $3$ semaines
    \item Temps d'assemblage d'un composant électronique ($0,8$ s)
    \item Temps de séchage automatique ($45$ minutes)
\end{enumerate}

\subsection{Corrigés}

\subsubsection*{Corrigé 15.1}

\begin{enumerate}
    \item $2,5\ \text{h} = 2,5 \times 100 = 250\ \text{ch}$
    \item $45\ \text{min} = 45 \times 166,67 = 7\,500\ \text{dmh}$
    \item $3\,600\ \text{s} = 3\,600 \times 1,667 = 6\,000\ \text{cmin}$
    \item $150\ \text{dmh} = 150 \div 10\,000 = 0,015\ \text{h}$
    \item $8\ \text{min} = 800\ \text{cmin}$
    \item $0,75\ \text{h} = 0,75 \times 3\,600 = 2\,700\ \text{s}$
    \item $12\,500\ \text{cmh} = 12\,500 \div 100\,000 = 0,125\ \text{h}$
    \item $1\ \text{h}\ 15\ \text{min}\ 30\ \text{s} = 1 + \frac{15}{60} + \frac{30}{3600} = 1,2583\ \text{h} = 125,83\ \text{ch}$
\end{enumerate}

\subsubsection*{Corrigé 15.2}

\begin{enumerate}
    \item $0:15:30 = 15 + \frac{30}{60} = 15,5\ \text{min} = \frac{15,5}{60} = 0,2583\ \text{h} = 25,83\ \text{ch}$
    \item $1:45:20 = 1 + \frac{45}{60} + \frac{20}{3600} = 1,7556\ \text{h} = 175,56\ \text{ch}$
    \item $2:08:45 = 2 + \frac{8}{60} + \frac{45}{3600} = 2,1458\ \text{h} = 214,58\ \text{ch}$
    \item $0:05:15 = 5 + \frac{15}{60} = 5,25\ \text{min} = \frac{5,25}{60} = 0,0875\ \text{h} = 8,75\ \text{ch}$
    \item $3:30:00 = 3 + \frac{30}{60} = 3,5\ \text{h} = 350\ \text{ch}$
\end{enumerate}

\subsubsection*{Corrigé 15.3}

\begin{enumerate}
    \item Opération manuelle de $3$ s : utiliser la centiminute (cmin) ou le TMU pour plus de précision
    \item Charge hebdomadaire : utiliser l'heure ou la centiheure (ch)
    \item Cycle machine $2$ min : utiliser la centiminute (cmin) ou le dmh
    \item Temps de traversée $3$ semaines : utiliser le jour (j) ou la semaine (sem)
    \item Assemblage $0,8$ s : utiliser le TMU ($1\ \text{TMU} = 0,036\ \text{s}$)
    \item Séchage $45$ min : utiliser la centiheure (ch) ou la minute
\end{enumerate}

\section{Exercices de Calcul du Nombre d'Observations}

\subsection{Application de la formule statistique}

\subsubsection*{Énoncé}

Un chronométreur réalise $15$ mesures préliminaires sur une séquence de travail. Les résultats en centiminutes sont :

\begin{center}
\begin{tabular}{|c|c|c|c|c|c|c|c|c|c|c|c|c|c|c|}
\hline
12 & 13 & 11 & 12 & 14 & 13 & 12 & 11 & 13 & 12 & 14 & 13 & 12 & 11 & 35 \\
\hline
\end{tabular}
\end{center}

\begin{enumerate}
    \item Éliminer les valeurs aberrantes.
    \item Calculer la moyenne et l'écart-type.
    \item Déterminer le nombre d'observations nécessaires pour un niveau de confiance de $95\%$ et une erreur admissible de $\pm 5\%$.
\end{enumerate}

\subsection{Utilisation des tableaux}

\subsubsection*{Énoncé}

Une opération a un temps de cycle de $3,5$ minutes. La production annuelle prévue est de $8\,000$ unités.

\begin{enumerate}
    \item En utilisant le tableau des nombres de cycles, déterminer le nombre minimum de cycles à chronométrer.
    \item Proposer une valeur pour une étude de fixation de tâche.
\end{enumerate}

\subsection{Cas avec différents niveaux de confiance}

\subsubsection*{Énoncé}

Une étude préliminaire donne une moyenne de $25$ secondes avec un écart-type de $3$ secondes.

\begin{enumerate}
    \item Calculer le nombre d'observations nécessaires pour :
    \begin{itemize}
        \item Niveau de confiance $90\%$, erreur $\pm 5\%$
        \item Niveau de confiance $95\%$, erreur $\pm 5\%$
        \item Niveau de confiance $99\%$, erreur $\pm 5\%$
    \end{itemize}
    \item Que constatez-vous ?
\end{enumerate}

\subsection{Corrigés}

\subsubsection*{Corrigé 16.1}

\begin{enumerate}
    \item Valeur aberrante : $35$ (casse de fil). Valeurs valides : 12, 13, 11, 12, 14, 13, 12, 11, 13, 12, 14, 13, 12, 11
    \item $\overline{X} = \frac{12+13+11+12+14+13+12+11+13+12+14+13+12+11}{14} = \frac{173}{14} = 12,36$ cmin
    \item $\sigma = \sqrt{\frac{\sum (X_i - \overline{X})^2}{n-1}} \approx 1,08$ cmin
    \item $n = \left(\frac{1,96 \times 1,08}{0,05 \times 12,36}\right)^2 = \left(\frac{2,1168}{0,618}\right)^2 = (3,425)^2 \approx 12$ observations
\end{enumerate}

\subsubsection*{Corrigé 16.2}

D'après le tableau (durée $2-5$ min, volume $1\,000-10\,000$) :
\begin{itemize}
    \item Nombre minimum : $15$ cycles
    \item Pour fixation de tâche : $20$ à $30$ cycles
\end{itemize}

\subsubsection*{Corrigé 16.3}

Valeurs : $\overline{X} = 25$ s, $\sigma = 3$ s, $E = 5\% = 0,05$

\begin{itemize}
    \item Niveau $90\%$ ($t=1,65$) : $n = \left(\frac{1,65 \times 3}{0,05 \times 25}\right)^2 = \left(\frac{4,95}{1,25}\right)^2 = (3,96)^2 \approx 16$
    \item Niveau $95\%$ ($t=1,96$) : $n = \left(\frac{1,96 \times 3}{1,25}\right)^2 = \left(\frac{5,88}{1,25}\right)^2 = (4,704)^2 \approx 23$
    \item Niveau $99\%$ ($t=2,58$) : $n = \left(\frac{2,58 \times 3}{1,25}\right)^2 = \left(\frac{7,74}{1,25}\right)^2 = (6,192)^2 \approx 39$
\end{itemize}

\textbf{Constat :} Plus le niveau de confiance est élevé, plus le nombre d'observations augmente.

\section{Exercices de Courbe d'Apprentissage}

\subsection{Application de la loi de Wright}

\subsubsection*{Énoncé}

Une entreprise réalise $10$ prototypes. Le premier prototype a demandé $25$ heures de travail. Le coefficient d'apprentissage est estimé à $85\%$.

\begin{enumerate}
    \item Calculer le temps nécessaire pour le $5^{\text{ème}}$ prototype.
    \item Calculer le temps nécessaire pour le $10^{\text{ème}}$ prototype.
    \item Calculer le temps cumulé pour les $10$ prototypes.
\end{enumerate}

\subsection{Application de la loi de Caquot}

\subsubsection*{Énoncé}

Les coûts de fabrication de la $10\,000^{\text{ème}}$ unité sont évalués à $3,12\ \text{\euro}$. Une commande de $40\,000$ unités est reçue.

\begin{enumerate}
    \item Calculer le coût de la $50\,000^{\text{ème}}$ unité.
    \item Calculer le coût de la $20\,000^{\text{ème}}$ unité.
\end{enumerate}

\subsection{Calcul du coefficient d'apprentissage}

\subsubsection*{Énoncé}

Un employé a pris $10$ heures pour exécuter une tâche la première fois, $8$ heures la deuxième fois, $6,5$ heures la quatrième fois.

\begin{enumerate}
    \item Déterminer le coefficient d'apprentissage moyen.
    \item Estimer le temps pour la $3^{\text{ème}}$ exécution.
    \item Estimer le temps pour la $8^{\text{ème}}$ exécution.
\end{enumerate}

\subsection{Estimation des temps cumulés}

\subsubsection*{Énoncé}

Le premier chantier dure $20$ jours. Le coefficient d'apprentissage est de $92\%$.

\begin{enumerate}
    \item Calculer la durée du $2^{\text{ème}}$ chantier.
    \item Calculer la durée du $5^{\text{ème}}$ chantier.
    \item Calculer la durée cumulée pour les $5$ chantiers.
\end{enumerate}

\subsection{Corrigés}

\subsubsection*{Corrigé 17.1}

Loi de Wright : $T_n = T_1 \times n^{\ln p / \ln 2}$

$e = \frac{\ln 0,85}{\ln 2} = \frac{-0,1625}{0,6931} = -0,2345$

\begin{enumerate}
    \item $T_5 = 25 \times 5^{-0,2345} = 25 \times 0,685 = 17,13$ heures
    \item $T_{10} = 25 \times 10^{-0,2345} = 25 \times 0,585 = 14,63$ heures
    \item Temps cumulé : $T_{\text{cum}} = T_1 + T_2 + ... + T_{10}$ (à calculer par table)
\end{enumerate}

\subsubsection*{Corrigé 17.2}

Loi de Caquot : $C_n = C_i \times \left(\frac{Q_n}{Q_i}\right)^{1/4}$

\begin{enumerate}
    \item $C_{50\,000} = 3,12 \times \left(\frac{50\,000}{10\,000}\right)^{0,25} = 3,12 \times 5^{0,25} = 3,12 \times 1,495 = 4,66\ \text{\euro}$
    \item $C_{20\,000} = 3,12 \times \left(\frac{20\,000}{10\,000}\right)^{0,25} = 3,12 \times 2^{0,25} = 3,12 \times 1,189 = 3,71\ \text{\euro}$
\end{enumerate}

\subsubsection*{Corrigé 17.3}

\begin{enumerate}
    \item Coefficient d'apprentissage :
    \begin{itemize}
        \item $T_2/T_1 = 8/10 = 0,80 = 80\%$
        \item $T_4/T_2 = 6,5/8 = 0,8125 = 81\%$
        \item Coefficient moyen : $80,5\%$
    \end{itemize}
    \item $T_3 = T_1 \times 3^{\ln 0,805 / \ln 2} = 10 \times 3^{-0,312} = 10 \times 0,71 = 7,1$ heures
    \item $T_8 = T_4 \times 2^{\ln 0,805 / \ln 2} = 6,5 \times 0,805 = 5,23$ heures
\end{enumerate}

\subsubsection*{Corrigé 17.4}

$e = \frac{\ln 0,92}{\ln 2} = \frac{-0,0834}{0,6931} = -0,1203$

\begin{enumerate}
    \item $T_2 = 20 \times 2^{-0,1203} = 20 \times 0,92 = 18,4$ jours
    \item $T_5 = 20 \times 5^{-0,1203} = 20 \times 0,824 = 16,48$ jours
    \item $T_{\text{cum}} = T_1 + T_2 + T_3 + T_4 + T_5 = 20 + 18,4 + 17,2 + 16,8 + 16,48 = 88,88$ jours
\end{enumerate}

\section{Exercices de Simogramme}

\subsection{Construction d'un simogramme}

\subsubsection*{Énoncé}

Une entreprise de confection dispose des données suivantes pour la mise en plis de pantalon :

\begin{table}[ht]
\centering
\begin{tabular}{|l|c|}
\hline
\textbf{Opération} & \textbf{Durée (dmh)} \\
\hline
Prendre et poser pantalon & 5 (Tm) \\
Positionner le pantalon & 12 (Tm) \\
Verrouiller, actionner vapeur & 8 (Ttm) \\
Mise en forme automatique & 60 (Tt) \\
Actionner séchage & 5 (Ttm) \\
Séchage automatique & 30 (Tt) \\
Déverrouiller & 5 (Ttm) \\
Évacuer le pantalon & 10 (Tm) \\
\hline
\end{tabular}
\caption{Données pour le simogramme}
\end{table}

\begin{enumerate}
    \item Construire le simogramme du poste avec une machine.
    \item Déterminer le temps de cycle.
    \item Calculer la production horaire.
\end{enumerate}

\subsection{Optimisation homme-machine}

\subsubsection*{Énoncé}

À partir des données de l'exercice précédent, on envisage de faire travailler un opérateur sur deux machines. Un temps de déplacement de $5$ dmh est nécessaire entre les deux machines.

\begin{enumerate}
    \item Construire le nouveau simogramme.
    \item Déterminer le nouveau temps de cycle.
    \item Calculer la nouvelle production horaire.
    \item Calculer le gain de productivité.
\end{enumerate}

\subsection{Calcul des productions horaires}

\subsubsection*{Énoncé}

Une ligne de production comporte $5$ postes avec les temps suivants :

\begin{itemize}
    \item Poste 1 : $2,5$ min
    \item Poste 2 : $3,0$ min
    \item Poste 3 : $2,8$ min
    \item Poste 4 : $2,2$ min
    \item Poste 5 : $2,5$ min
\end{itemize}

\begin{enumerate}
    \item Déterminer le temps de cycle de la ligne.
    \item Calculer la production horaire.
    \item Identifier le goulot d'étranglement.
\end{enumerate}

\subsection{Corrigés}

\subsubsection*{Corrigé 18.1}

\begin{enumerate}
    \item Simogramme : représentation graphique avec l'opérateur et la machine en ordonnée, temps en abscisse
    \item Temps de cycle = $5+12+8+60+5+30+5+10 = 135$ dmh
    \item Production horaire = $10\,000 \div 135 = 74$ pantalons/heure
\end{enumerate}

\subsubsection*{Corrigé 18.2}

\begin{enumerate}
    \item Nouveau simogramme : alternance entre les deux machines
    \item Nouveau temps de cycle = $135 + 5 = 140$ dmh (avec déplacement)
    \item Production horaire = $(10\,000 \div 140) \times 2 = 143$ pantalons/heure
    \item Gain = $143 - 74 = 69$ pantalons/heure (+93\%)
\end{enumerate}

\subsubsection*{Corrigé 18.3}

\begin{enumerate}
    \item Temps de cycle = $3,0$ min (poste 2)
    \item Production horaire = $60 \div 3,0 = 20$ unités/heure
    \item Goulot = Poste 2 ($3,0$ min)
\end{enumerate}

\section{Exercices de Chronométrage et Dépouillement}

\subsection{Calcul du temps observé moyen}

\subsubsection*{Énoncé}

Une séquence de travail a été chronométrée $10$ fois. Les résultats en centiminutes sont :

\begin{center}
\begin{tabular}{|c|c|c|c|c|c|c|c|c|c|}
\hline
134 & 132 & 130 & 135 & 138 & 133 & 131 & 129 & 136 & 132 \\
\hline
\end{tabular}
\end{center}

Calculer le temps observé moyen.

\subsection{Correction par le jugement d'allure}

\subsubsection*{Énoncé}

Pour la séquence précédente, le jugement d'allure de l'opérateur est de $105$.

\begin{enumerate}
    \item Calculer le temps de base (temps à l'allure $100$).
    \item Commenter le résultat.
\end{enumerate}

\subsection{Application des coefficients majorateurs}

\subsubsection*{Énoncé}

La séquence étudiée est un temps manuel (Tm). Le coefficient majorateur pour ce type d'opération est de $1,12$.

\begin{enumerate}
    \item Calculer le temps standard après majoration.
    \item Si la séquence était un temps technico-manuel (coefficient $1,20$), quel serait le temps standard ?
\end{enumerate}

\subsection{Intégration de la fréquence}

\subsubsection*{Énoncé}

La séquence chronométrée concerne la pose d'un passant. Chaque pantalon comporte $6$ passants.

\begin{enumerate}
    \item Calculer le temps alloué par pantalon pour cette séquence.
    \item Calculer le temps total pour $1\,000$ pantalons.
\end{enumerate}

\subsection{Calcul complet du temps standard}

\subsubsection*{Énoncé}

Pour une opération d'assemblage, les données suivantes sont disponibles :

\begin{table}[ht]
\centering
\begin{tabular}{|l|c|c|c|c|}
\hline
\textbf{Séquence} & $\overline{T}_o$ (cmin) & \textbf{JA} & \textbf{Coefficient} & \textbf{Fréquence} \\
\hline
Préparation & 45 & 100 & 1,12 & 1 \\
Assemblage & 120 & 95 & 1,20 & 1 \\
Contrôle & 30 & 110 & 1,10 & 1 \\
Évacuation & 25 & 105 & 1,12 & 1 \\
\hline
\end{tabular}
\caption{Données pour le calcul complet}
\end{table}

Calculer le temps standard total.

\subsection{Corrigés}

\subsubsection*{Corrigé 19.1}

\[
\overline{T}_o = \frac{134+132+130+135+138+133+131+129+136+132}{10} = \frac{1330}{10} = 133\ \text{cmin}
\]

\subsubsection*{Corrigé 19.2}

\[
T_{100} = 133 \times \frac{105}{100} = 139,65\ \text{cmin}
\]

L'opérateur travaille plus vite que la normale (JA > 100), donc son temps observé est inférieur au temps de base.

\subsubsection*{Corrigé 19.3}

\begin{enumerate}
    \item $T_S = 139,65 \times 1,12 = 156,41$ cmin
    \item $T_S = 139,65 \times 1,20 = 167,58$ cmin
\end{enumerate}

\subsubsection*{Corrigé 19.4}

\begin{enumerate}
    \item $T_{\text{alloué par pantalon}} = 156,41 \times 6 = 938,46$ cmin
    \item $T_{\text{total}} = 938,46 \times 1\,000 = 938\,460$ cmin = $9\,384,6$ min = $156,41$ heures
\end{enumerate}

\subsubsection*{Corrigé 19.5}

\begin{enumerate}
    \item Préparation : $T_{100} = 45 \times 100/100 = 45$ cmin ; $T_S = 45 \times 1,12 = 50,4$ cmin
    \item Assemblage : $T_{100} = 120 \times 95/100 = 114$ cmin ; $T_S = 114 \times 1,20 = 136,8$ cmin
    \item Contrôle : $T_{100} = 30 \times 110/100 = 33$ cmin ; $T_S = 33 \times 1,10 = 36,3$ cmin
    \item Évacuation : $T_{100} = 25 \times 105/100 = 26,25$ cmin ; $T_S = 26,25 \times 1,12 = 29,4$ cmin
\end{enumerate}

\textbf{Temps standard total :} $50,4 + 136,8 + 36,3 + 29,4 = 252,9$ cmin

\section{Exercices de Jugement d'Allure}

\subsection{Évaluation par méthode directe}

\subsubsection*{Énoncé}

Un opérateur est chronométré sur une séquence. Le temps observé est de $45$ secondes. Le chronométreur estime que l'opérateur travaille à une allure de $90$ (échelle BTE).

\begin{enumerate}
    \item Calculer le temps à l'allure $100$.
    \item Si le même opérateur avait été évalué à $110$, quel aurait été le temps à l'allure $100$ ?
\end{enumerate}

\subsection{Application de la méthode de Levelling}

\subsubsection*{Énoncé}

Un opérateur est chronométré à $18$ dmh. Le chronométreur évalue les facteurs suivants :

\begin{itemize}
    \item Habileté : Excellente ($+0,11$)
    \item Activité : Bonne ($+0,06$)
    \item Conditions de travail : Passable ($-0,03$)
    \item Régularité : Bonne ($+0,01$)
\end{itemize}

\begin{enumerate}
    \item Calculer le temps corrigé par la méthode de Levelling.
    \item Convertir le résultat en secondes.
\end{enumerate}

\subsection{Conversion entre échelles}

\subsubsection*{Énoncé}

Un opérateur est évalué à $95$ sur l'échelle BTE.

\begin{enumerate}
    \item Convertir cette valeur sur l'échelle Bedaux.
    \item Convertir cette valeur sur l'échelle Maynard.
    \item Quelle est l'allure correspondante en km/h ?
\end{enumerate}

\subsection{Corrigés}

\subsubsection*{Corrigé 20.1}

\begin{enumerate}
    \item $T_{100} = 45 \times \frac{90}{100} = 40,5$ secondes
    \item $T_{100} = 45 \times \frac{110}{100} = 49,5$ secondes
\end{enumerate}

\subsubsection*{Corrigé 20.2}

\begin{enumerate}
    \item Total correctifs = $0,11 + 0,06 - 0,03 + 0,01 = 0,15$
    \item $T_{\text{corrigé}} = 18 \times (1 + 0,15) = 20,7$ dmh
    \item $20,7\ \text{dmh} = 20,7 \times 0,36 = 7,45$ secondes
\end{enumerate}

\subsubsection*{Corrigé 20.3}

D'après le tableau de correspondance :
\begin{enumerate}
    \item Échelle Bedaux : $95$ BTE $\rightarrow$ $90$ Bedaux
    \item Échelle Maynard : $95$ BTE $\rightarrow$ $127$ Maynard
    \item Allure en km/h : $95$ BTE $\rightarrow$ $5,8$ km/h
\end{enumerate}

\section{Exercices d'Équilibrage de Lignes}

\subsection{Calcul du temps de cycle}

\subsubsection*{Énoncé}

Une ligne de production comporte $6$ postes avec les temps suivants :

\begin{itemize}
    \item Poste 1 : $2,2$ min
    \item Poste 2 : $2,5$ min
    \item Poste 3 : $3,0$ min
    \item Poste 4 : $2,8$ min
    \item Poste 5 : $2,4$ min
    \item Poste 6 : $2,1$ min
\end{itemize}

\begin{enumerate}
    \item Déterminer le temps de cycle de la ligne.
    \item Calculer la production horaire.
    \item Identifier le goulot.
\end{enumerate}

\subsection{Répartition des charges}

\subsubsection*{Énoncé}

Une ligne doit produire $500$ unités par jour ($7$ heures de travail). Les temps standards des opérations sont :

\begin{itemize}
    \item Opération A : $1,2$ min
    \item Opération B : $1,5$ min
    \item Opération C : $2,0$ min
    \item Opération D : $1,8$ min
    \item Opération E : $1,5$ min
\end{itemize}

\begin{enumerate}
    \item Calculer le temps de cycle nécessaire pour atteindre l'objectif.
    \item Déterminer le nombre théorique de postes nécessaires.
    \item Proposer une répartition des opérations par poste.
\end{enumerate}

\subsection{Calcul de l'efficacité de la ligne}

\subsubsection*{Énoncé}

Après équilibrage, une ligne de $4$ postes a les temps suivants :
\begin{itemize}
    \item Poste 1 : $2,7$ min
    \item Poste 2 : $2,8$ min
    \item Poste 3 : $2,7$ min
    \item Poste 4 : $2,8$ min
\end{itemize}

\begin{enumerate}
    \item Calculer l'efficacité de la ligne.
    \item Calculer l'indice d'équilibrage.
\end{enumerate}

\subsection{Optimisation par regroupement}

\subsubsection*{Énoncé}

Une ligne de $5$ postes a les temps suivants :
\begin{itemize}
    \item Poste 1 : $3,0$ min
    \item Poste 2 : $2,5$ min
    \item Poste 3 : $2,8$ min
    \item Poste 4 : $2,2$ min
    \item Poste 5 : $2,5$ min
\end{itemize}

\begin{enumerate}
    \item Calculer le temps de cycle actuel.
    \item Calculer l'efficacité actuelle.
    \item Proposer une nouvelle répartition pour réduire le temps de cycle à $2,6$ min.
\end{enumerate}

\subsection{Corrigés}

\subsubsection*{Corrigé 21.1}

\begin{enumerate}
    \item Temps de cycle = $3,0$ min (poste 3)
    \item Production horaire = $60 \div 3,0 = 20$ unités/heure
    \item Goulot = Poste 3
\end{enumerate}

\subsubsection*{Corrigé 21.2}

\begin{enumerate}
    \item Temps de cycle nécessaire = $\frac{7 \times 60}{500} = \frac{420}{500} = 0,84$ min/unité
    \item Somme des temps = $1,2+1,5+2,0+1,8+1,5 = 8,0$ min
    \item Nombre de postes = $\frac{8,0}{0,84} \approx 10$ postes
\end{enumerate}

\subsubsection*{Corrigé 21.3}

\begin{enumerate}
    \item Temps de cycle = $2,8$ min
    \item Somme des temps = $2,7+2,8+2,7+2,8 = 11,0$ min
    \item Efficacité = $\frac{11,0}{2,8 \times 4} \times 100 = \frac{11,0}{11,2} \times 100 = 98,2\%$
\end{enumerate}

\subsubsection*{Corrigé 21.4}

\begin{enumerate}
    \item Temps de cycle actuel = $3,0$ min
    \item Somme des temps = $3,0+2,5+2,8+2,2+2,5 = 13,0$ min
    \item Efficacité actuelle = $\frac{13,0}{3,0 \times 5} \times 100 = \frac{13,0}{15} \times 100 = 86,7\%$
\end{enumerate}

% ==================== FIN DE LA PARTIE 7 ====================
% ==================== PARTIE 8 : CAS RÉELS INDUSTRIELS ====================
\part*{Partie 8 : Cas Réels Industriels}
\addcontentsline{toc}{part}{Partie 8 : Cas Réels Industriels}

\section{Cas Réel 1 : Industrie Textile}

\subsection{Contexte : Ligne de confection de pantalons}

\textit{Confection Mode SA} est une entreprise textile spécialisée dans la fabrication de pantalons pour femmes. Elle emploie $150$ personnes et produit $2\,000$ pantalons par jour sur une ligne de $12$ postes.

\textbf{Données clés :}
\begin{itemize}
    \item Production journalière : $2\,000$ pantalons
    \item Nombre de postes : $12$
    \item Temps de travail journalier : $7$ heures ($420$ minutes)
    \item Taux de présence moyen : $92\%$
    \item Coût minute : $0,45$ €
    \item Effectif : $150$ personnes
\end{itemize}

\subsection{Problématique : Disparités de charge et retards}

L'entreprise constate plusieurs problèmes :
\begin{itemize}
    \item Disparités importantes dans les charges de travail hebdomadaires
    \item Retards récurrents sur les délais de livraison
    \item Augmentation des défauts qualité sur certains postes
    \item Stress et fatigue chez les opératrices des postes chargés
    \item Taux d'absentéisme plus élevé sur les postes goulots
\end{itemize}

\subsection{Démarche complète : diagnostic → étude → améliorations → fixation}

\subsubsection{Étape 1 : Chronométrage de diagnostic}

Un premier chronométrage de diagnostic est réalisé sur les $12$ postes avec $3$ relevés par poste :

\begin{table}[ht]
\centering
\begin{tabular}{|c|c|c|c|c|}
\hline
\textbf{Poste} & \textbf{Opération} & \textbf{Relevé 1} & \textbf{Relevé 2} & \textbf{Relevé 3} \\
\hline
1 & Coupe & $2,1$ & $2,2$ & $2,0$ \\
2 & Assemblage dos & $3,5$ & $3,4$ & $3,6$ \\
3 & Pose fermeture & $5,2$ & $5,5$ & $5,3$ \\
4 & Assemblage devant & $3,2$ & $3,3$ & $3,1$ \\
5 & Pose poches & $2,8$ & $2,9$ & $2,7$ \\
6 & Assemblage côtés & $3,4$ & $3,5$ & $3,3$ \\
7 & Pose ceinture & $3,1$ & $3,2$ & $3,0$ \\
8 & Ourlet & $2,5$ & $2,6$ & $2,4$ \\
9 & Contrôle & $1,8$ & $1,9$ & $1,7$ \\
10 & Repassage & $2,2$ & $2,3$ & $2,1$ \\
11 & Emballage & $1,5$ & $1,6$ & $1,4$ \\
12 & Expédition & $1,2$ & $1,3$ & $1,1$ \\
\hline
\end{tabular}
\caption{Résultats du chronométrage de diagnostic (temps en minutes)}
\end{table}

\textbf{Résultat :} Le poste 3 "Pose fermeture" présente un temps moyen de $5,33$ minutes, soit $65\%$ supérieur à la moyenne des autres postes ($3,2$ min). Ce poste est identifié comme goulot.

\subsubsection{Étape 2 : Chronométrage d'étude}

Un chronométrage approfondi est réalisé sur le poste "Pose fermeture" avec $15$ relevés :

\begin{table}[ht]
\centering
\begin{tabular}{|c|c|c|c|c|c|}
\hline
\textbf{N°} & \textbf{Temps} & \textbf{N°} & \textbf{Temps} & \textbf{N°} & \textbf{Temps} \\
\hline
1 & $5,2$ & 6 & $5,8$ & 11 & $5,3$ \\
2 & $5,3$ & 7 & $7,2$ & 12 & $5,4$ \\
3 & $5,1$ & 8 & $5,4$ & 13 & $5,2$ \\
4 & $5,4$ & 9 & $5,3$ & 14 & $6,5$ \\
5 & $5,2$ & 10 & $5,1$ & 15 & $5,3$ \\
\hline
\end{tabular}
\caption{Temps observés pour le poste "Pose fermeture" (en minutes)}
\end{table}

\textbf{Analyse :}
\begin{itemize}
    \item Valeurs aberrantes : $7,2$ (casse de fil) et $6,5$ (approvisionnement défaillant)
    \item Taux d'aléas = $2/15 \times 100 = 13,3\%$ (poste moyennement stable)
    \item Temps observé moyen = $(5,2+5,3+5,1+5,4+5,2+5,8+5,4+5,3+5,1+5,3+5,4+5,2+5,3) / 13 = 5,28$ min
\end{itemize}

\subsubsection{Étape 3 : Identification des causes}

L'observation approfondie révèle :
\begin{itemize}
    \item L'opératrice n'a pas reçu de formation spécifique sur ce poste
    \item L'approvisionnement en fermetures est irrégulier ($2$ à $3$ arrêts par heure)
    \item La posture de travail est inconfortable (chaise non réglée)
    \item Le mode opératoire n'est pas standardisé
    \item Les gestes sont hésitants et comportent des mouvements inutiles
\end{itemize}

\subsubsection{Étape 4 : Plan d'amélioration}

\begin{enumerate}
    \item \textbf{Formation} : $2$ jours de formation sur le mode opératoire standard
    \item \textbf{Approvisionnement} : Mise en place de bacs de proximité (réduction des déplacements)
    \item \textbf{Ergonomie} : Réglage de la chaise et de la hauteur de table
    \item \textbf{Standardisation} : Rédaction d'un mode opératoire illustré
    \item \textbf{Organisation} : Regroupement des fermetures par lots de $50$
\end{enumerate}

\subsubsection{Étape 5 : Chronométrage de fixation de tâche}

Après mise en œuvre des améliorations, un nouveau chronométrage est réalisé ($30$ relevés) :

\begin{table}[ht]
\centering
\begin{tabular}{|l|c|c|c|}
\hline
\textbf{Séquence} & $\overline{T}_o$ (min) & \textbf{JA} & $T_{100}$ (min) \\
\hline
Prendre pantalon & $0,45$ & $100$ & $0,45$ \\
Positionner fermeture & $0,82$ & $105$ & $0,86$ \\
Piquer fermeture & $1,95$ & $100$ & $1,95$ \\
Contrôler & $0,38$ & $95$ & $0,36$ \\
Évacuer & $0,25$ & $100$ & $0,25$ \\
\hline
\textbf{Total} & $3,85$ & & $3,87$ \\
\hline
\end{tabular}
\caption{Temps après amélioration}
\end{table}

\subsubsection{Étape 6 : Application des majorations}

Coefficients majorateurs :
\begin{itemize}
    \item Temps manuels : $1,12$
    \item Temps technico-manuels : $1,18$
\end{itemize}

\[
T_{\text{standard}} = 3,87 \times 1,12 = 4,33 \text{ minutes}
\]

\subsection{Résultats et gains}

\begin{table}[ht]
\centering
\begin{tabular}{|l|c|c|c|}
\hline
\textbf{Indicateur} & \textbf{Avant} & \textbf{Après} & \textbf{Gain} \\
\hline
Temps de cycle (min) & $5,33$ & $4,33$ & $-19\%$ \\
Production horaire & $11,3$ & $13,9$ & $+23\%$ \\
Taux d'aléas & $13,3\%$ & $4,5\%$ & $-66\%$ \\
Coût unitaire (€) & $2,40$ & $1,95$ & $-0,45$ €/pièce \\
Production journalière & $2\,000$ & $2\,460$ & $+460$ pantalons \\
\hline
\end{tabular}
\caption{Résultats de l'optimisation}
\end{table}

\subsection{Questions de réflexion}

\begin{enumerate}
    \item Quelles autres améliorations pourriez-vous proposer pour ce poste ?
    \item Comment généraliser cette démarche aux autres postes de la ligne ?
    \item Quel serait l'impact sur la production journalière si tous les postes étaient optimisés ?
    \item Proposez un nouvel équilibrage de la ligne avec les nouveaux temps.
    \item Calculez le retour sur investissement si les améliorations ont coûté $5\,000$ €.
\end{enumerate}

\section{Cas Réel 2 : Industrie Automobile}

\subsection{Contexte : Poste d'assemblage de portières}

\textit{AutoParts SA} est un équipementier automobile qui fournit des portières complètes pour un constructeur français. L'atelier d'assemblage produit $800$ portières par jour sur une ligne automatisée.

\textbf{Données clés :}
\begin{itemize}
    \item Production journalière : $800$ portières
    \item Temps de cycle actuel : $3,0$ minutes
    \item Nombre d'opérateurs : $4$ par ligne
    \item Temps de travail : $7$ heures $30$ minutes
    \item Coût minute : $0,80$ €
    \item Investissement machine : $150\,000$ €
\end{itemize}

\subsection{Problématique : Goulot d'étranglement}

Le poste d'assemblage des garnitures intérieures est identifié comme goulot :
\begin{itemize}
    \item Temps de cycle : $3,0$ min (les autres postes sont à $2,5$ min)
    \item File d'attente moyenne : $15$ portières
    \item Taux d'occupation : $98\%$
    \item Retards réguliers sur les livraisons
\end{itemize}

\subsection{Analyse par simogramme}

\subsubsection{Décomposition du poste goulot}

\begin{table}[ht]
\centering
\begin{tabular}{|l|c|c|}
\hline
\textbf{Opération} & \textbf{Durée (s)} & \textbf{Nature} \\
\hline
Prise de la portière & $8$ & Tm \\
Positionnement garniture & $12$ & Tm \\
Mise en place & $25$ & Ttm \\
Fixation automatique & $60$ & Tt \\
Contrôle visuel & $10$ & Tm \\
Évacuation & $5$ & Tm \\
\hline
\end{tabular}
\caption{Décomposition du poste goulot}
\end{table}

\subsubsection{Simogramme avant optimisation}

\begin{center}
\begin{minipage}{0.9\textwidth}
\begin{verbatim}
Temps (s)     0    20    40    60    80   100   120
Opérateur   |---Tm---|---Ttm--|     |--Tm--|     |---Tm---|
Machine     |          |---Ttm--|-------Tt-------|--Ttm--|
\end{verbatim}
\caption{Représentation simplifiée du simogramme avant optimisation}
\end{minipage}
\end{center}

\textbf{Analyse :}
\begin{itemize}
    \item Temps manuel total : $8+12+10+5 = 35$ s
    \item Temps machine : $60$ s
    \item Temps technico-manuel : $25$ s
    \item Temps de cycle : $35+25+60 = 120$ s
    \item L'opérateur est inactif pendant $60$ s (attente machine)
\end{itemize}

\subsection{Optimisation par ajout de ressources}

\subsubsection{Solution 1 : Ajout d'une machine}

\begin{table}[ht]
\centering
\begin{tabular}{|l|c|}
\hline
\textbf{Configuration} & \textbf{Temps de cycle} \\
\hline
1 machine & $120$ s \\
2 machines & $35+25 = 60$ s \\
\hline
\end{tabular}
\caption{Impact de l'ajout d'une machine}
\end{table}

\subsubsection{Solution 2 : Délégation des opérations manuelles}

Proposition de décomposer le travail :
\begin{itemize}
    \item Opérateur 1 : Prise portière, positionnement garniture
    \item Machine : Fixation automatique
    \item Opérateur 2 : Contrôle, évacuation
\end{itemize}

Nouveau temps de cycle : $35$ s (opérations en parallèle)

\subsubsection{Simogramme après optimisation}

\begin{center}
\begin{minipage}{0.9\textwidth}
\begin{verbatim}
Temps (s)     0    20    40    60    80   100   120
Opérateur 1 |---Tm---|---Ttm--|                       |
Machine     |          |---Ttm--|-------Tt-------|   |
Opérateur 2 |                       |---Tm---|--Tm--|
\end{verbatim}
\caption{Représentation simplifiée du simogramme après optimisation}
\end{minipage}
\end{center}

\subsection{Résultats et gains}

\begin{table}[ht]
\centering
\begin{tabular}{|l|c|c|c|}
\hline
\textbf{Indicateur} & \textbf{Avant} & \textbf{Après} & \textbf{Gain} \\
\hline
Temps de cycle & $120$ s & $60$ s & $-50\%$ \\
Production horaire & $30$ & $60$ & $+100\%$ \\
Production journalière & $800$ & $1\,600$ & $+800$ portières \\
Taux d'occupation opérateur & $50\%$ & $85\%$ & $+35\%$ \\
Coût unitaire & $4,80$ € & $2,40$ € & $-50\%$ \\
\hline
\end{tabular}
\caption{Résultats de l'optimisation}
\end{table}

\textbf{Investissement :}
\begin{itemize}
    \item Ajout d'une machine : $150\,000$ €
    \item Formation : $5\,000$ €
    \item Retour sur investissement : $150\,000 / (800 \times 2,40) \approx 78$ jours
\end{itemize}

\section{Cas Réel 3 : Industrie Agroalimentaire}

\subsection{Contexte : Ligne de conditionnement de salades}

\textit{FreshSalad SA} est une entreprise agroalimentaire spécialisée dans le conditionnement de salades prêtes à l'emploi. La ligne de conditionnement emploie $25$ opérateurs et produit $10\,000$ sachets par jour.

\textbf{Données clés :}
\begin{itemize}
    \item Production journalière : $10\,000$ sachets
    \item Nombre d'opérateurs : $25$
    \item Temps de travail : $7$ heures
    \item Taux de présence : $95\%$
    \item Coût minute : $0,35$ €
\end{itemize}

\subsection{Problématique : Forte variabilité des performances}

L'entreprise constate :
\begin{itemize}
    \item Une variabilité importante des performances entre opérateurs ($\pm 30\%$)
    \item Des pics de production suivis de ralentissements
    \item Un stress important pendant les "grosses journées"
    \item Des conséquences sur la non-qualité des produits
\end{itemize}

\subsection{Utilisation des observations instantanées}

\subsubsection{Étape 1 : Définition des états}

\begin{table}[ht]
\centering
\begin{tabular}{|c|l|}
\hline
\textbf{État} & \textbf{Description} \\
\hline
A & Conditionnement actif \\
B & Attente matière \\
C & Contrôle qualité \\
D & Réglage machine \\
E & Pause / besoins personnels \\
F & Autre \\
\hline
\end{tabular}
\caption{États observés sur la ligne de conditionnement}
\end{table}

\subsubsection{Étape 2 : Pré-étude par auto-pointage}

Une semaine d'auto-pointage donne les proportions estimées :
\begin{itemize}
    \item État A : $65\%$
    \item État B : $12\%$
    \item État C : $8\%$
    \item État D : $5\%$
    \item État E : $7\%$
    \item État F : $3\%$
\end{itemize}

\subsubsection{Étape 3 : Calcul du nombre d'observations}

Niveau de confiance : $95\%$, marge d'erreur : $5\%$

\[
n = \frac{1,96^2 \times p \times (1-p)}{E^2}
\]

Pour l'état A ($p=0,65$) : $n = \frac{3,8416 \times 0,65 \times 0,35}{0,0025} = \frac{0,874}{0,0025} = 350$ observations

\subsubsection{Étape 4 : Réalisation des observations}

$350$ observations sont réalisées sur $2$ semaines. Résultats :

\begin{table}[ht]
\centering
\begin{tabular}{|c|c|c|}
\hline
\textbf{État} & \textbf{Nb observations} & \textbf{Proportion} \\
\hline
A & $227$ & $64,9\%$ \\
B & $42$ & $12,0\%$ \\
C & $28$ & $8,0\%$ \\
D & $18$ & $5,1\%$ \\
E & $24$ & $6,9\%$ \\
F & $11$ & $3,1\%$ \\
\hline
\end{tabular}
\caption{Résultats des observations instantanées}
\end{table}

\subsection{Identification des causes}

L'analyse des résultats révèle :
\begin{itemize}
    \item $12\%$ de temps d'attente matière (problème d'approvisionnement)
    \item $8\%$ de temps de contrôle (redondant avec contrôle final)
    \item $5\%$ de réglages machines (fréquence élevée)
    \item Variabilité entre opérateurs : $20\%$ à $85\%$ d'activité
\end{itemize}

\subsection{Plan d'amélioration et résultats}

\subsubsection{Actions mises en œuvre}

\begin{enumerate}
    \item \textbf{Approvisionnement} : Mise en place d'un système Kanban
    \item \textbf{Contrôle} : Regroupement des contrôles en fin de ligne
    \item \textbf{Réglages} : Formation aux réglages et maintenance préventive
    \item \textbf{Standardisation} : Création d'un mode opératoire unique
    \item \textbf{Formation} : Formation croisée entre opérateurs
\end{enumerate}

\subsubsection{Résultats après 3 mois}

\begin{table}[ht]
\centering
\begin{tabular}{|l|c|c|c|}
\hline
\textbf{Indicateur} & \textbf{Avant} & \textbf{Après} & \textbf{Gain} \\
\hline
Temps d'attente matière & $12\%$ & $4\%$ & $-67\%$ \\
Temps de contrôle & $8\%$ & $3\%$ & $-62\%$ \\
Activité moyenne & $65\%$ & $82\%$ & $+26\%$ \\
Production horaire & $1\,428$ & $1\,800$ & $+26\%$ \\
Variabilité inter-opérateurs & $\pm 30\%$ & $\pm 10\%$ & $-67\%$ \\
Taux de non-qualité & $3,5\%$ & $1,8\%$ & $-49\%$ \\
\hline
\end{tabular}
\caption{Résultats de l'optimisation}
\end{table}

\section{Cas Réel 4 : Industrie Mécanique}

\subsection{Contexte : Atelier d'usinage}

\textit{MecaPrecision SA} est une entreprise de mécanique de précision. Elle doit lancer la production d'une nouvelle pièce pour l'aéronautique. Le poste d'usinage n'existe pas encore physiquement.

\textbf{Données clés :}
\begin{itemize}
    \item Commande : $5\,000$ pièces
    \item Délai : $3$ mois
    \item Poste à créer : centre d'usinage $5$ axes
    \item Temps machine estimé : $15$ min/pièce
    \item Coût minute machine : $2,50$ €
    \item Coût minute opérateur : $0,60$ €
\end{itemize}

\subsection{Problématique : Définition des temps pour nouveau produit}

L'entreprise doit :
\begin{itemize}
    \item Estimer le temps de cycle avant la mise en production
    \item Déterminer le nombre d'opérateurs nécessaires
    \item Calculer le coût de revient pour le devis
    \item Planifier la production
\end{itemize}

\subsection{Utilisation de la méthode MOST}

\subsubsection{Décomposition du travail}

Le travail est décomposé en séquences MOST :

\begin{table}[ht]
\centering
\begin{tabular}{|l|c|}
\hline
\textbf{Séquence} & \textbf{Durée estimée (TMU)} \\
\hline
Prise de la pièce brute & $40$ \\
Positionnement dans l'étau & $60$ \\
Bridage & $30$ \\
Démarrage machine & $20$ \\
Usinage automatique & $(1\,200)$ \\
Débridage & $30$ \\
Retrait pièce & $40$ \\
Contrôle visuel & $30$ \\
Évacuation & $20$ \\
\hline
\end{tabular}
\caption{Décomposition MOST}
\end{table}

\subsubsection{Calcul du temps standard}

\begin{itemize}
    \item Temps manuel total : $40+60+30+20+30+40+30+20 = 270$ TMU
    \item $270$ TMU $\times 0,036 = 9,72$ secondes $\approx 0,162$ minutes
    \item Temps machine : $15$ minutes
    \item Temps de cycle = $15 + 0,162 = 15,162$ minutes
\end{itemize}

\subsubsection{Application des majorations}

\begin{itemize}
    \item Coefficient majorateur temps manuel : $1,15$
    \item Coefficient majorateur temps machine : $1,02$
\end{itemize}

\[
T_{\text{standard}} = (0,162 \times 1,15) + (15 \times 1,02) = 0,186 + 15,3 = 15,486 \text{ minutes}
\]

\subsection{Comparaison avec chronométrage}

Après mise en production, un chronométrage de confirmation est réalisé :

\begin{table}[ht]
\centering
\begin{tabular}{|l|c|c|}
\hline
\textbf{Indicateur} & \textbf{MOST} & \textbf{Chronométrage} \\
\hline
Temps manuel & $0,162$ min & $0,170$ min \\
Temps machine & $15,0$ min & $14,8$ min \\
Temps standard & $15,486$ min & $15,350$ min \\
Écart & & $+0,9\%$ \\
\hline
\end{tabular}
\caption{Comparaison MOST vs chronométrage}
\end{table}

\subsection{Validation et résultats}

\begin{table}[ht]
\centering
\begin{tabular}{|l|c|}
\hline
\textbf{Indicateur} & \textbf{Valeur} \\
\hline
Temps standard retenu & $15,4$ min \\
Production horaire & $3,9$ pièces/heure \\
Production journalière ($7$h) & $27$ pièces/jour \\
Nombre d'opérateurs nécessaires & $2$ (en alternance) \\
Coût unitaire (main d'œuvre) & $0,154$ € \\
Coût unitaire (machine) & $38,5$ € \\
Coût total de revient & $42,3$ € \\
\hline
\end{tabular}
\caption{Résultats finaux}
\end{table}

\subsection{Leçons apprises}

\begin{itemize}
    \item Le MOST a permis d'estimer les temps avant la création du poste
    \item L'écart avec le chronométrage réel est inférieur à $1\%$
    \item La méthode a été validée par le client pour le devis
    \item La formation des opérateurs a été facilitée par le mode opératoire MOST
    \item Le temps standard sert maintenant de base pour la production en série
\end{itemize}

% ==================== FIN DE LA PARTIE 8 ====================
% ==================== PARTIE 9 : ANNEXES ET SYNTHÈSE ====================
\part*{Partie 9 : Annexes et Synthèse}
\addcontentsline{toc}{part}{Partie 9 : Annexes et Synthèse}

\section{Annexes}

\subsection{Tableaux de conversion des unités}

\subsubsection{Tableau 1 : Conversion des unités classiques vers les unités décimales}

\begin{table}[ht]
\centering
\begin{tabular}{|l|c|c|c|c|}
\hline
\textbf{Unité} & \textbf{En heures} & \textbf{En minutes} & \textbf{En secondes} & \textbf{Facteur} \\
\hline
1 heure (h) & 1 & 60 & 3 600 & 1 \\
1 minute (min) & 0,016667 & 1 & 60 & 1/60 \\
1 seconde (s) & 0,000278 & 0,016667 & 1 & 1/3600 \\
\hline
\end{tabular}
\caption{Conversion des unités classiques}
\end{table}

\subsubsection{Tableau 2 : Unités décimales et leurs équivalences}

\begin{table}[ht]
\centering
\begin{tabular}{|l|c|c|c|c|}
\hline
\textbf{Unité} & \textbf{Symbole} & \textbf{Équivalence} & \textbf{En minutes} & \textbf{En secondes} \\
\hline
Décitheure & dh & $10^{-1}$ h & 6 min & 360 s \\
Centiheure & ch & $10^{-2}$ h & 0,6 min & 36 s \\
Milliheure & mh & $10^{-3}$ h & 0,06 min & 3,6 s \\
Décimilliheure & dmh & $10^{-4}$ h & 0,006 min & 0,36 s \\
Centimilliheure & cmh & $10^{-5}$ h & 0,0006 min & 0,036 s \\
\hline
Déciminute & dmin & $10^{-1}$ min & 0,1 min & 6 s \\
Centiminute & cmin & $10^{-2}$ min & 0,01 min & 0,6 s \\
\hline
\end{tabular}
\caption{Unités décimales et leurs équivalences}
\end{table}

\subsubsection{Tableau 3 : Table de conversion rapide}

\begin{table}[ht]
\centering
\begin{tabular}{|c|c|c|c|c|c|c|}
\hline
\textbf{cmh} & \textbf{dmh} & \textbf{h} & \textbf{cmin} & \textbf{min} & \textbf{cs} & \textbf{s} \\
\hline
100 000 & 10 000 & 1 & 6 000 & 60 & 360 000 & 3 600 \\
50 000 & 5 000 & 0,5 & 3 000 & 30 & 180 000 & 1 800 \\
10 000 & 1 000 & 0,1 & 600 & 6 & 36 000 & 360 \\
5 000 & 500 & 0,05 & 300 & 3 & 18 000 & 180 \\
1 000 & 100 & 0,01 & 60 & 0,6 & 3 600 & 36 \\
500 & 50 & 0,005 & 30 & 0,3 & 1 800 & 18 \\
100 & 10 & 0,001 & 6 & 0,06 & 360 & 3,6 \\
50 & 5 & 0,0005 & 3 & 0,03 & 180 & 1,8 \\
10 & 1 & 0,0001 & 0,6 & 0,006 & 36 & 0,36 \\
\hline
\end{tabular}
\caption{Table de conversion rapide}
\end{table}

\subsection{Tables MTM abrégées}

\subsubsection{Tableau 4 : Mouvement "Atteindre" (R) - Extrait}

\begin{table}[ht]
\centering
\begin{tabular}{|c|c|c|c|c|c|c|c|}
\hline
\textbf{Distance (cm)} & \textbf{R-A} & \textbf{R-B} & \textbf{R-C} & \textbf{R-D} & \textbf{R-E} & \textbf{mR-A} & \textbf{R-A m} \\
\hline
2 & 2,0 & 2,0 & 2,0 & 2,0 & 1,6 & 1,6 & 1,4 \\
4 & 3,3 & 3,3 & 5,2 & 3,3 & 3,0 & 2,5 & 0,8 \\
6 & 4,5 & 4,5 & 6,5 & 4,5 & 3,9 & 3,0 & 1,5 \\
8 & 5,4 & 5,6 & 7,5 & 7,5 & 5,5 & 3,6 & 2,0 \\
10 & 6,0 & 6,6 & 8,4 & 6,4 & 4,9 & 4,2 & 2,4 \\
12 & 6,4 & 7,4 & 9,1 & 7,1 & 5,2 & 4,8 & 2,6 \\
14 & 6,7 & 8,2 & 9,7 & 7,7 & 5,5 & 5,3 & 2,9 \\
16 & 7,1 & 8,8 & 10,3 & 8,2 & 5,8 & 5,9 & 2,9 \\
18 & 7,4 & 9,4 & 10,8 & 8,7 & 6,1 & 6,5 & 2,9 \\
20 & 7,8 & 9,9 & 11,4 & 9,2 & 6,4 & 7,1 & 2,8 \\
22 & 8,1 & 10,5 & 11,9 & 9,7 & 6,8 & 7,6 & 2,9 \\
24 & 8,5 & 11,1 & 12,5 & 10,2 & 7,1 & 8,2 & 2,9 \\
26 & 8,8 & 11,6 & 13,0 & 10,6 & 7,4 & 8,8 & 2,8 \\
28 & 9,2 & 12,2 & 13,6 & 11,1 & 7,7 & 9,4 & 2,8 \\
30 & 9,5 & 12,8 & 14,1 & 11,6 & 8,0 & 9,9 & 2,9 \\
\hline
\end{tabular}
\caption{Table MTM du mouvement "Atteindre" (valeurs en TMU)}
\end{table}

\subsubsection{Tableau 5 : Mouvement "Mouvoir" (M) - Extrait}

\begin{table}[ht]
\centering
\begin{tabular}{|c|c|c|c|}
\hline
\textbf{Distance (cm)} & \textbf{M-A} & \textbf{M-B} & \textbf{M-C} \\
\hline
2 & 2,0 & 2,0 & 2,0 \\
4 & 3,4 & 3,4 & 3,4 \\
6 & 4,5 & 4,5 & 4,8 \\
8 & 5,5 & 5,6 & 6,1 \\
10 & 6,1 & 6,4 & 7,0 \\
12 & 6,6 & 7,1 & 7,8 \\
14 & 7,1 & 7,7 & 8,6 \\
16 & 7,6 & 8,2 & 9,3 \\
18 & 8,0 & 8,7 & 10,0 \\
20 & 8,4 & 9,2 & 10,6 \\
\hline
\end{tabular}
\caption{Table MTM du mouvement "Mouvoir" (valeurs en TMU)}
\end{table}

\subsection{Tables de coefficients majorateurs}

\subsubsection{Tableau 6 : Coefficients pour travaux manuels}

\begin{table}[ht]
\centering
\begin{tabular}{|l|c|}
\hline
\textbf{Type d'opération} & \textbf{Coefficient} \\
\hline
Manipulations (saisir, tenir, déposer, évacuer) & 1,10 \\
Dégagements (dégager pince, couper fils) & 1,10 \\
Contrôles (examiner, compter, vérifier) & 1,10 \\
Préparations (plier, déplier, épingler) & 1,12 \\
Présentations (orienter, maintenir, engager) & 1,12 \\
Compléments d'opérations (dégarnir, cranter) & 1,13 \\
Finition main & 1,12 \\
Réglage main, coupe ciseaux main & 1,15 \\
\hline
\end{tabular}
\caption{Coefficients majorateurs pour travaux manuels}
\end{table}

\subsubsection{Tableau 7 : Coefficients pour machines à coudre}

\begin{table}[ht]
\centering
\begin{tabular}{|l|c|}
\hline
\textbf{Type de machine} & \textbf{Coefficient} \\
\hline
Piqueuse plate 1 aiguille & 1,20 \\
Piqueuse plate 2 aiguilles & 1,23 \\
Piqueuse plate 3 aiguilles & 1,28 \\
Piqueuse plate 4 aiguilles & 1,33 \\
Surjeteuse 1 aiguille 2 fils & 1,13 \\
Surjeteuse 1 aiguille 3 fils & 1,15 \\
Surjeteuse 2 aiguilles 4 fils & 1,17 \\
Surjeteuse 2 aiguilles 5 fils & 1,18 \\
Machine à bras déporté & 1,18 \\
Machine à point d'arrêt & 1,17 \\
Machine à boutonnières & 1,12 \\
Machine à poser boutons & 1,16 \\
\hline
\end{tabular}
\caption{Coefficients majorateurs pour machines à coudre}
\end{table}

\subsubsection{Tableau 8 : Coefficients pour machines d'usinage}

\begin{table}[ht]
\centering
\begin{tabular}{|l|c|}
\hline
\textbf{Type de machine} & \textbf{Coefficient} \\
\hline
Tour conventionnel & 1,15 \\
Fraiseuse conventionnelle & 1,15 \\
Centre d'usinage CNC & 1,05 \\
Tour CNC & 1,05 \\
Rectifieuse & 1,12 \\
Presse plieuse & 1,10 \\
\hline
\end{tabular}
\caption{Coefficients majorateurs pour machines d'usinage}
\end{table}

\subsection{Modèles de feuilles de relevés de chronométrage}

\subsubsection{Modèle 1 : Feuille de chronométrage standard}

\begin{table}[ht]
\centering
\begin{tabular}{|l|l|}
\hline
\multicolumn{2}{|c|}{\textbf{FEUILLE DE CHRONOMÉTRAGE}} \\
\hline
\textbf{Article :} & \\
\textbf{Modèle :} & \\
\textbf{Atelier :} & \\
\textbf{Matière :} & \\
\textbf{Machine :} & \\
\textbf{Opératrice :} & \\
\textbf{Chronométreur :} & \\
\textbf{Date :} & \\
\textbf{Points/cm :} & \textbf{RPM :} \\
\hline
\end{tabular}
\caption{En-tête de la feuille de chronométrage}
\end{table}

\begin{table}[ht]
\centering
\begin{tabular}{|l|c|c|c|c|c|c|c|c|c|c|c|}
\hline
\textbf{Séquence} & \multicolumn{10}{|c|}{\textbf{Temps observés (cmin)}} & \textbf{JA} \\
\cline{2-11}
 & 1 & 2 & 3 & 4 & 5 & 6 & 7 & 8 & 9 & 10 & \\
\hline
& & & & & & & & & & & \\
& & & & & & & & & & & \\
& & & & & & & & & & & \\
& & & & & & & & & & & \\
& & & & & & & & & & & \\
\hline
\end{tabular}
\caption{Corps de la feuille de chronométrage}
\end{table}

\subsection{Modèles de feuilles d'observations instantanées}

\subsubsection{Modèle 2 : Feuille d'observations instantanées}

\begin{table}[ht]
\centering
\begin{tabular}{|l|l|}
\hline
\multicolumn{2}{|c|}{\textbf{FEUILLE D'OBSERVATIONS INSTANTANÉES}} \\
\hline
\textbf{Ressource étudiée :} & \\
\textbf{Période :} & \\
\textbf{Observateur :} & \\
\textbf{Niveau de confiance :} & 95\% \\
\textbf{Marge d'erreur :} & \\
\hline
\end{tabular}
\caption{En-tête de la feuille d'observations instantanées}
\end{table}

\begin{table}[ht]
\centering
\begin{tabular}{|c|c|c|c|c|c|c|}
\hline
\textbf{Date} & \textbf{Horaire} & \textbf{État A} & \textbf{État B} & \textbf{État C} & \textbf{État D} & \textbf{Observations} \\
\hline
& & & & & & \\
& & & & & & \\
& & & & & & \\
& & & & & & \\
& & & & & & \\
\hline
\end{tabular}
\caption{Corps de la feuille d'observations instantanées}
\end{table}

\subsubsection{Modèle 3 : Feuille d'auto-pointage}

\begin{table}[ht]
\centering
\begin{tabular}{|l|l|l|}
\hline
\multicolumn{3}{|c|}{\textbf{FEUILLE D'AUTO-POINTAGE}} \\
\hline
\textbf{Opérateur :} & \textbf{Date :} & \\
\hline
\textbf{Horaire début} & \textbf{Activité} & \textbf{Horaire fin} \\
\hline
& & \\
& & \\
& & \\
& & \\
& & \\
& & \\
\hline
\end{tabular}
\caption{Feuille d'auto-pointage}
\end{table}

\section{Glossaire}

\subsection{Termes techniques et définitions}

\subsubsection{Termes généraux}

\begin{description}[labelwidth=4cm]
    \item[Chronométrage] Technique de mesure directe du temps consistant à observer et mesurer le temps nécessaire à l'exécution d'une tâche.
    
    \item[Jugement d'allure (JA)] Évaluation subjective du rythme de travail d'un opérateur par rapport à une allure de référence (allure 100).
    
    \item[Temps de base] Temps obtenu après correction du temps observé par le jugement d'allure. Correspond au temps à l'allure 100.
    
    \item[Temps standard] Temps de base majoré pour prendre en compte les temps improductifs inévitables (fatigue, besoins personnels, aléas).
    
    \item[Temps observé] Temps mesuré directement lors du chronométrage, sans correction.
    
    \item[Taux d'aléas] Pourcentage de mesures aberrantes dans une série de chronométrages.
\end{description}

\subsubsection{Unités de mesure}

\begin{description}[labelwidth=4cm]
    \item[TMU (Time Measurement Unit)] Unité de base du MTM. $1\ \text{TMU} = 0,00001\ \text{heure} = 0,036\ \text{seconde}$.
    
    \item[DMH (Dix-Millième d'Heure)] $1\ \text{DMH} = 0,0001\ \text{heure} = 0,36\ \text{seconde}$.
    
    \item[Centiheure (ch)] $1\ \text{ch} = 0,01\ \text{heure} = 36\ \text{secondes}$.
    
    \item[Centiminute (cmin)] $1\ \text{cmin} = 0,01\ \text{minute} = 0,6\ \text{seconde}$.
\end{description}

\subsubsection{Outils de représentation}

\begin{description}[labelwidth=4cm]
    \item[Simogramme] Représentation graphique d'un cycle de production permettant de visualiser les temps manuels, technologiques et technico-manuels.
    
    \item[VSM (Value Stream Mapping)] Cartographie des flux de valeur, représentant les flux physiques et d'information d'un processus.
    
    \item[Analyse de déroulement] Méthode d'observation qui suit les étapes de transformation de la matière en produit fini.
\end{description}

\subsubsection{Méthodes de temps prédéterminés}

\begin{description}[labelwidth=4cm]
    \item[MTM (Methods-Time Measurement)] Méthode de temps prédéterminés analysant les mouvements élémentaires avec des tables standardisées.
    
    \item[MOST (Maynard Operation Sequence Technique)] Méthode de temps prédéterminés utilisant des séquences types (ABG ABP A).
    
    \item[Therbligs] Décomposition du travail en $17$ mouvements élémentaires, développée par Frank et Lillian Gilbreth.
    
    \item[Levelling] Méthode d'évaluation du jugement d'allure basée sur l'analyse de quatre facteurs (habileté, activité, conditions, régularité).
\end{description}

\section{Bibliographie et Références}

\subsection{Ouvrages de référence}

\begin{thebibliography}{99}
    \bibitem{taylor} TAYLOR, Frederick Winslow. \textit{The Principles of Scientific Management}. 1911.
    
    \bibitem{gilbreth} GILBRETH, Frank B. et GILBRETH, Lillian M. \textit{Motion Study}. 1911.
    
    \bibitem{maynard} MAYNARD, H.B., STEGEMERTEN, G.J. et SCHWAB, J.L. \textit{Methods-Time Measurement}. McGraw-Hill, 1948.
    
    \bibitem{zandin} ZANDIN, Kjell B. \textit{MOST Work Measurement Systems}. Marcel Dekker, 1990.
    
    \bibitem{kanawaty} KANAWATY, George. \textit{Introduction to Work Study}. International Labour Office, 4th edition, 1992.
    
    \bibitem{freivalds} FREIVALDS, Andris et NIEBEL, Benjamin. \textit{Methods, Standards, and Work Design}. McGraw-Hill, 12th edition, 2009.
    
    \bibitem{afnor} AFNOR. \textit{Étude du travail - Principes généraux}. NF X 06-001.
\end{thebibliography}

\subsection{Articles et publications}

\begin{thebibliography}{99}
    \bibitem{wright} WRIGHT, T.P. "Factors Affecting the Cost of Airplanes". \textit{Journal of the Aeronautical Sciences}, 1936.
    
    \bibitem{caquot} CAQUOT, Albert. "Loi d'évolution des prix". \textit{Revue de Métallurgie}, 1950.
    
    \bibitem{bedaux} BEDAUX, Charles. "The Bedaux System". \textit{Industrial Management}, 1920.
    
    \bibitem{ilomtm} ILO. "Introduction to Work Study". \textit{International Labour Organization}, 4th edition, 1992.
    
    \bibitem{maynard2} MAYNARD, H.B. "The Evolution of Predetermined Time Systems". \textit{Journal of Industrial Engineering}, 1955.
\end{thebibliography}

\subsection{Normes et standards}

\begin{thebibliography}{99}
    \bibitem{iso1} ISO 14000 - Management environnemental.
    
    \bibitem{iso2} ISO 9000 - Management de la qualité.
    
    \bibitem{mtm1} MTM Association. \textit{MTM-1 Standard Data Cards}. MTM Association for Standards and Research.
    
    \bibitem{mtm2} MTM Association. \textit{MTM-2 Standard Data Cards}. MTM Association for Standards and Research.
    
    \bibitem{most} MOST Association. \textit{Basic MOST Standard Data}. Maynard Operation Sequence Technique Association.
    
    \bibitem{afnor1} AFNOR. \textit{NF X 06-001 - Étude du travail - Principes généraux}.
    
    \bibitem{afnor2} AFNOR. \textit{NF X 06-002 - Étude du travail - Chronométrage}.
\end{thebibliography}

\section{Corrigés des Exercices}

\subsection{Corrigés des exercices de conversion (Section 15)}

\subsubsection{Exercice 15.1 - Corrigé}

\begin{enumerate}
    \item $2,5\ \text{h} = 2,5 \times 100 = 250\ \text{ch}$
    \item $45\ \text{min} = 45 \times 166,67 = 7\,500\ \text{dmh}$
    \item $3\,600\ \text{s} = 3\,600 \times 1,667 = 6\,000\ \text{cmin}$
    \item $150\ \text{dmh} = 150 \div 10\,000 = 0,015\ \text{h}$
    \item $8\ \text{min} = 800\ \text{cmin}$
    \item $0,75\ \text{h} = 0,75 \times 3\,600 = 2\,700\ \text{s}$
    \item $12\,500\ \text{cmh} = 12\,500 \div 100\,000 = 0,125\ \text{h}$
    \item $1\ \text{h}\ 15\ \text{min}\ 30\ \text{s} = 1 + 15/60 + 30/3600 = 1,2583\ \text{h} = 125,83\ \text{ch}$
\end{enumerate}

\subsubsection{Exercice 15.2 - Corrigé}

\begin{enumerate}
    \item $0:15:30 = 15 + 30/60 = 15,5\ \text{min} = 15,5/60 = 0,2583\ \text{h} = 25,83\ \text{ch}$
    \item $1:45:20 = 1 + 45/60 + 20/3600 = 1,7556\ \text{h} = 175,56\ \text{ch}$
    \item $2:08:45 = 2 + 8/60 + 45/3600 = 2,1458\ \text{h} = 214,58\ \text{ch}$
\end{enumerate}

\subsection{Corrigés des exercices de courbe d'apprentissage (Section 17)}

\subsubsection{Exercice 17.1 - Corrigé}

$e = \frac{\ln 0,85}{\ln 2} = \frac{-0,1625}{0,6931} = -0,2345$

\begin{enumerate}
    \item $T_5 = 25 \times 5^{-0,2345} = 25 \times 0,685 = 17,13$ heures
    \item $T_{10} = 25 \times 10^{-0,2345} = 25 \times 0,585 = 14,63$ heures
\end{enumerate}

\subsubsection{Exercice 17.3 - Corrigé}

\begin{enumerate}
    \item Coefficient d'apprentissage : $T_2/T_1 = 8/10 = 0,80$ (80\%), $T_4/T_2 = 6,5/8 = 0,8125$ (81\%), coefficient moyen = 80,5\%
    \item $T_3 = 10 \times 3^{\ln 0,805 / \ln 2} = 10 \times 3^{-0,312} = 7,1$ heures
\end{enumerate}

\subsection{Corrigés des exercices de chronométrage (Section 19)}

\subsubsection{Exercice 19.1 - Corrigé}

\[
\overline{T}_o = \frac{134+132+130+135+138+133+131+129+136+132}{10} = \frac{1330}{10} = 133\ \text{cmin}
\]

\subsubsection{Exercice 19.2 - Corrigé}

\[
T_{100} = 133 \times \frac{105}{100} = 139,65\ \text{cmin}
\]

\subsubsection{Exercice 19.5 - Corrigé}

\begin{table}[ht]
\centering
\begin{tabular}{|l|c|c|c|}
\hline
\textbf{Séquence} & $T_{100}$ (cmin) & \textbf{Coeff} & $T_S$ (cmin) \\
\hline
Préparation & $45$ & $1,12$ & $50,4$ \\
Assemblage & $114$ & $1,20$ & $136,8$ \\
Contrôle & $33$ & $1,10$ & $36,3$ \\
Évacuation & $26,25$ & $1,12$ & $29,4$ \\
\hline
\multicolumn{3}{|c|}{\textbf{Total}} & \textbf{252,9 cmin} \\
\hline
\end{tabular}
\caption{Corrigé exercice 19.5}
\end{table}

\subsection{Corrigés des exercices de jugement d'allure (Section 20)}

\subsubsection{Exercice 20.1 - Corrigé}

\begin{enumerate}
    \item $T_{100} = 45 \times 90/100 = 40,5$ secondes
    \item $T_{100} = 45 \times 110/100 = 49,5$ secondes
\end{enumerate}

\subsubsection{Exercice 20.2 - Corrigé}

Total correctifs = $0,11 + 0,06 - 0,03 + 0,01 = 0,15$
\[
T_{\text{corrigé}} = 18 \times (1 + 0,15) = 20,7\ \text{dmh} = 7,45\ \text{secondes}
\]

\subsection{Corrigés des exercices de simogramme (Section 18)}

\subsubsection{Exercice 18.1 - Corrigé}

\begin{itemize}
    \item Temps de cycle = $5+12+8+60+5+30+5+10 = 135$ dmh
    \item Production horaire = $10\,000 \div 135 = 74$ pantalons/heure
\end{itemize}

\subsubsection{Exercice 18.2 - Corrigé}

\begin{itemize}
    \item Nouveau temps de cycle = $135 + 5 = 140$ dmh
    \item Production horaire = $(10\,000 \div 140) \times 2 = 143$ pantalons/heure
    \item Gain = $143 - 74 = 69$ pantalons/heure (+93\%)
\end{itemize}

\subsection{Corrigés des exercices d'équilibrage (Section 21)}

\subsubsection{Exercice 21.1 - Corrigé}

\begin{itemize}
    \item Temps de cycle = $3,0$ min (poste 3)
    \item Production horaire = $60 \div 3,0 = 20$ unités/heure
    \item Goulot = Poste 3
\end{itemize}

\subsubsection{Exercice 21.3 - Corrigé}

\begin{itemize}
    \item Somme des temps = $2,7+2,8+2,7+2,8 = 11,0$ min
    \item Efficacité = $\frac{11,0}{2,8 \times 4} \times 100 = 98,2\%$
\end{itemize}

% ==================== FIN DE LA PARTIE 9 ====================
\documentclass[a4paper,12pt]{article}

% ==================== PACKAGES ====================
\usepackage[french]{babel}      % Pour la langue française
\usepackage[T1]{fontenc}        % Encodage des polices
\usepackage[utf8]{inputenc}     % Encodage des caractères (accents)
\usepackage{geometry}           % Pour gérer les marges
\usepackage{amsmath, amssymb}   % Pour les mathématiques
\usepackage{enumitem}           % Pour personnaliser les listes
\usepackage{titlesec}           % Pour personnaliser les titres
\usepackage{xcolor}             % Pour les couleurs
\usepackage{hyperref}           % Pour les liens internes
\usepackage{amsthm}             % Pour les environnements théorèmes
\usepackage{graphicx}           % Pour les images (optionnel)
\usepackage{array}              % Pour les tableaux avancés
\usepackage{booktabs}           % Pour de meilleurs tableaux
\usepackage{eurosym}            % Pour le symbole euro

% Configuration des listes pour éviter les débordements
\setlist{leftmargin=*, itemsep=2pt, parsep=2pt, labelwidth=2em}

% ==================== ENVIRONNEMENTS PERSONNALISÉS ====================
\newtheorem{definition}{Définition}[section]
\newtheorem{exemple}{Exemple}[section]

% ==================== MISE EN PAGE ====================
\geometry{
    top=2.5cm,
    bottom=2.5cm,
    left=2.5cm,
    right=2.5cm
}

% Personnalisation des titres de sections
\titleformat{\section}
    {\normalfont\Large\bfseries\color{blue!70!black}}
    {\thesection}{1em}{}

\titleformat{\subsection}
    {\normalfont\large\bfseries\color{blue!50!black}}
    {\thesubsection}{1em}{}

\titleformat{\subsubsection}
    {\normalfont\normalsize\bfseries\color{blue!40!black}}
    {\thesubsubsection}{1em}{}

% ==================== DOCUMENT ====================
\begin{document}

% ==================== PAGE DE TITRE ====================
\begin{titlepage}
    \centering
    \vspace*{2cm}
    
    {\Huge \textbf{Étude sur le Chronométrage}}\\[0.5cm]
    {\Large Méthodologie et Applications}\\[1.5cm]
    
    \vfill
    
    \begin{minipage}{0.7\textwidth}
        \centering
        \rule{\textwidth}{0.5pt}\\[0.5cm]
        \textbf{Cours de Mesure du Travail et Facteur Humain}\\
        \vspace{0.3cm}
        \textit{Document de cours - Année académique 2025-2026}
    \end{minipage}
    
    \vfill
    
    {\large \today}
\end{titlepage}

% ==================== SOMMAIRE ====================
\tableofcontents
\newpage

% ==================== PARTIE 1 : INTRODUCTION ====================
\part*{Partie 1 : Introduction}
\addcontentsline{toc}{part}{Partie 1 : Introduction}

\section{Définition et objectifs de la mesure du travail}

\subsection{Définition}

La mesure du travail est une technique fondamentale de l'organisation industrielle. Elle consiste à déterminer le temps nécessaire à un opérateur qualifié pour accomplir une tâche donnée, selon une méthode définie et à un niveau de performance déterminé.

\begin{definition}
La \textbf{mesure du travail} est l'application de techniques visant à déterminer le temps que demande à un ouvrier qualifié l'exécution d'une tâche donnée, à un niveau de rendement bien défini.
\end{definition}

Cette définition soulève trois questions essentielles :
\begin{itemize}
    \item \textbf{Faire quoi ?} (définition précise de la tâche)
    \item \textbf{Avec quoi ?} (moyens et équipements utilisés)
    \item \textbf{Comment ?} (méthode opératoire)
\end{itemize}

\subsection{Objectifs de la mesure du travail}

La mesure du travail poursuit plusieurs objectifs stratégiques et opérationnels :

\begin{itemize}
    \item \textbf{Définir les prix et les coûts} : établir les coûts de revient des produits
    \item \textbf{Motiver les opérateurs} : fixer des objectifs réalistes et équitables
    \item \textbf{Comparer les possibilités de conception} : évaluer différentes solutions techniques
    \item \textbf{Établir les plannings} : planifier la production
    \item \textbf{Contrôler la stabilisation des postes de travail} : suivre la performance
    \item \textbf{Fixer le temps de fabrication} : déterminer le temps nécessaire à l'élaboration d'un produit
    \item \textbf{Étudier, diminuer et éliminer les temps improductifs}
    \item \textbf{Apprécier la performance} : évaluer l'efficacité des opérateurs et des équipes
\end{itemize}

\subsection{Productivité industrielle}

La productivité est le rapport du produit obtenu (production) aux ressources utilisées pour l'obtenir :

\[
\text{Productivité} = \frac{\text{Production}}{\text{Ressources utilisées}}
\]

La mesure de la productivité peut être de trois types :

\begin{itemize}
    \item \textbf{Partielle} : utilisation d'un seul intrant (ex : productivité du travail)
    \item \textbf{Multifactorielle} : utilisation de plusieurs intrants
    \item \textbf{Globale} : utilisation de tous les intrants
\end{itemize}

\section{Place du chronométrage dans l'étude du travail}

\subsection{Les deux piliers de l'étude du travail}

L'étude du travail repose sur deux piliers complémentaires :

\begin{enumerate}
    \item \textbf{L'étude des méthodes} : consiste à examiner les méthodes existantes et envisagées d'un travail, afin de mettre au point et de faire appliquer des méthodes d'exécution plus efficaces et de diminuer les coûts.
    
    \item \textbf{La mesure du travail} : a pour objectif de quantifier le temps nécessaire à l'exécution d'une tâche selon la méthode retenue.
\end{enumerate}

\subsection{Objectifs de l'étude des méthodes}

L'étude des méthodes vise à :
\begin{itemize}
    \item Éliminer les déplacements inutiles
    \item Réduire les éléments à non valeur ajoutée
    \item Éliminer les mouvements superflus
    \item Combiner des éléments de travail
    \item Réduire la fatigue des opérateurs
    \item Améliorer l'aménagement des lieux de travail
    \item Améliorer la conception des outils et du matériel utilisés
\end{itemize}

\subsection{Les techniques de mesure du travail}

\begin{table}[ht]
\centering
\begin{tabular}{|p{5cm}|p{5cm}|p{5cm}|}
\hline
\multicolumn{3}{|c|}{\textbf{Techniques d'étude des temps}} \\
\hline
\multicolumn{2}{|c|}{\textbf{Techniques par observation directe}} & \textbf{Techniques par mesure indirecte} \\
\cline{1-2}
Chrono-analyse & Mesure par échantillonnage & Prédétermination des temps \\
& & Données standards \\
\hline
\end{tabular}
\caption{Les principales techniques de mesure du travail}
\end{table}

\subsection{Place du chronométrage}

Le chronométrage est la technique de mesure directe la plus utilisée. Il intervient après la définition de la méthode optimale. Il permet de fixer des temps de référence qui serviront de base à :

\begin{itemize}
    \item La gestion de la production
    \item L'évaluation des performances
    \item Le calcul des coûts
    \item La définition des objectifs
\end{itemize}

\begin{center}
\fbox{\begin{minipage}{0.85\textwidth}
\centering
\textbf{Principe fondamental} \\
Le chronométrage mesure \textit{ce qui existe}, \\
l'étude des méthodes détermine \textit{ce qui devrait être}.
\end{minipage}}
\end{center}

\section{Historique et évolution : de Taylor aux temps prédéterminés}

\subsection{Frederick Winslow Taylor (1856-1915)}

Ingénieur américain, Taylor est le promoteur le plus connu de l'Organisation Scientifique du Travail (OST). Il est le premier à systématiser :

\begin{itemize}
    \item La décomposition des tâches en éléments simples
    \item L'utilisation du chronométrage pour mesurer les temps
    \item La définition de temps alloués pour chaque tâche
\end{itemize}

\begin{quote}
\textit{« Il existe toujours une méthode et un outil qui permettent un travail plus rapide et de meilleure qualité que les autres. »} \\
— Frederick Winslow Taylor
\end{quote}

\subsection{Frank et Lillian Gilbreth (années 1920)}

Ce couple de chercheurs a poussé l'analyse plus loin en décomposant chaque geste élémentaire du travailleur en mouvements de base appelés \textbf{therbligs}.

\begin{definition}
Les \textbf{therbligs} sont une décomposition du travail en 17 mouvements élémentaires de base. Le nom est une inversion de « Gilbreth » avec une transposition du « th ».
\end{definition}

Les principaux therbligs incluent :
\begin{itemize}
    \item Chercher, Trouver, Choisir
    \item Saisir, Transporter, Tenir
    \item Positionner, Assembler, Désassembler
    \item Utiliser, Attendre, Reposer
\end{itemize}

\subsection{L'avènement des temps prédéterminés (1940-1970)}

Face aux limites du chronométrage (subjectivité du jugement d'allure, nécessité d'un poste existant), des méthodes de temps prédéterminés apparaissent :

\begin{enumerate}
    \item \textbf{1948 : MTM (Methods-Time Measurement)} \\
    Développé par Maynard, Stegemerten et Schwab. Le MTM attribue des temps standards à des mouvements élémentaires codifiés. L'unité utilisée est le TMU (Time Measurement Unit) = $10^{-5}$ heure.
    
    \item \textbf{1950 : MTS (Methods Time Standards)} \\
    Développé par la General Electric.
    
    \item \textbf{1966 : MODAPTS} \\
    Développé par G.C. Hyde. (Modular Arrangement of Predetermined Time Standards)
    
    \item \textbf{1974 : MOST (Maynard Operation Sequence Technique)} \\
    Créé par Kjell Zandin, cette méthode simplifie et accélère l'analyse des temps. Une heure de cycle peut être chiffrée en une journée d'étude.
    
    \item \textbf{Autres méthodes} : WORK FACTOR (Philips), BTE (Bureau des Temps Élémentaires).
\end{enumerate}

\subsection{Tableau chronologique synthétique}

\begin{table}[ht]
\centering
\begin{tabular}{|c|c|c|}
\hline
\textbf{Année} & \textbf{Auteurs / Organismes} & \textbf{Apport} \\
\hline
1900-1915 & Frederick W. Taylor & Organisation Scientifique du Travail (OST) \\
1920 & Frank et Lillian Gilbreth & Therbligs (décomposition des mouvements) \\
1948 & H.B. Maynard et al. & MTM 1 \\
1950 & General Electric & MTS \\
1965 & H.B. Maynard & MTM 2 \\
1966 & G.C. Hyde & MODAPTS \\
1974 & Kjell Zandin & MOST \\
1975-1984 & Divers & MTM 3, MTM UAS, MTM MEK, MTM SAM \\
\hline
\end{tabular}
\caption{Évolution historique des méthodes de mesure du temps}
\end{table}

\subsection{Évolution contemporaine}

Aujourd'hui, le chronométrage et les temps prédéterminés sont utilisés de manière complémentaire, en fonction :

\begin{itemize}
    \item Du contexte de l'étude
    \item Des délais disponibles
    \item Des objectifs poursuivis
    \item Du niveau de précision requis
\end{itemize}

Les tendances actuelles sont à :
\begin{itemize}
    \item L'intégration des outils numériques (logiciels de capture vidéo, analyses automatisées)
    \item La prise en compte des facteurs ergonomiques dans la fixation des temps
    \item Une approche participative associant les opérateurs à la démarche
    \item La constitution de bases de données internes d'entreprises
\end{itemize}

\subsection{Organisation du cours}

Ce cours est structuré en plusieurs chapitres :

\begin{enumerate}
    \item Introduction à la mesure du travail
    \item Le temps dans l'industrie : unités et typologie
    \item La courbe d'apprentissage et l'évolution des temps
    \item Le chronométrage : méthodologie et mise en œuvre
    \item Le jugement d'allure et la détermination du temps standard
    \item Le dépouillement et l'exploitation des résultats
    \item Les observations instantanées
    \item Les outils de représentation des temps (simogramme, analyse de déroulement)
    \item Les temps prédéterminés (MTM, MOST)
    \item Interface Homme-Machine
\end{enumerate}

\subsection{Synthèse de l'introduction}

\begin{center}
\fbox{\begin{minipage}{0.9\textwidth}
\textbf{Points clés à retenir :}

\begin{itemize}
    \item La mesure du travail détermine le temps nécessaire à l'exécution d'une tâche
    \item Elle se distingue de l'étude des méthodes qui optimise la façon de travailler
    \item Le chronométrage est la technique de mesure directe la plus utilisée
    \item L'histoire de la mesure du travail a connu des étapes clés : Taylor, Gilbreth, MTM, MOST
    \item Les temps prédéterminés permettent de chiffrer des temps avant la création du poste
\end{itemize}
\end{minipage}}
\end{center}

% ==================== FIN DE L'INTRODUCTION ====================

% ==================== PARTIE 1 : FONDAMENTAUX ET PRÉREQUIS ====================
\part*{Partie 1 : Fondamentaux et Prérequis}
\addcontentsline{toc}{part}{Partie 1 : Fondamentaux et Prérequis}

\section{Le Temps dans l’Industrie}

\subsection{Unités de temps}

La mesure du temps dans l'industrie nécessite des unités adaptées aux calculs et aux systèmes informatiques. Deux grandes familles d'unités coexistent.

\subsubsection{Unités classiques (système sexagésimal)}

Ces unités sont issues du système traditionnel de mesure du temps :

\begin{table}[ht]
\centering
\begin{tabular}{|c|c|}
\hline
\textbf{Unité} & \textbf{Symbole} \\
\hline
Mois & m \\
Semaine & sem \\
Jour & j \\
Heure & h \\
Minute & min \\
Seconde & s \\
\hline
\end{tabular}
\caption{Unités classiques de temps}
\end{table}

\textbf{Inconvénient :} Les calculs sont complexes en raison des conversions ($1\ h = 60\ min$, $1\ min = 60\ s$), ce qui les rend peu pratiques pour les opérations arithmétiques fréquentes en gestion de production.

\subsubsection{Unités décimales}

Pour faciliter les calculs, l'industrie utilise des unités décimales basées sur l'heure ou la minute :

\begin{table}[ht]
\centering
\begin{tabular}{|c|c|c|}
\hline
\textbf{Unité} & \textbf{Symbole} & \textbf{Équivalence} \\
\hline
Décitheure & dh & $10^{-1}\ h = 6\ min$ \\
Centiheure & ch & $10^{-2}\ h = 0,6\ min = 36\ s$ \\
Milliheure & mh & $10^{-3}\ h = 0,06\ min = 3,6\ s$ \\
Décimilliheure & dmh & $10^{-4}\ h = 0,006\ min = 0,36\ s$ \\
Centimilliheure & cmh & $10^{-5}\ h = 0,0006\ min = 0,036\ s$ \\
\hline
Déciminute & dmin & $10^{-1}\ min = 6\ s$ \\
Centiminute & cmin & $10^{-2}\ min = 0,6\ s$ \\
\hline
\end{tabular}
\caption{Unités décimales de temps}
\end{table}

\begin{center}
\fbox{\begin{minipage}{0.9\textwidth}
\centering
\textbf{Remarque importante :} Le TMU (Time Measurement Unit), utilisé dans la méthode MTM, correspond à 1 cmh ($10^{-5}$ heure).
\end{minipage}}
\end{center}

\subsubsection{Table de conversion}

\begin{table}[ht]
\centering
\begin{tabular}{|c|c|c|c|c|c|c|}
\hline
\textbf{cmh} & \textbf{dmh} & \textbf{h} & \textbf{cmin} & \textbf{min} & \textbf{cs} & \textbf{s} \\
\hline
100 000 & 10 000 & 1 & 6 000 & 60 & 360 000 & 3 600 \\
\hline
\end{tabular}
\caption{Table de conversion entre les différentes unités}
\end{table}

\subsection{Choix de l'unité}

Le choix de l'unité de temps n'est pas anodin. Il doit être guidé par plusieurs critères.

\subsubsection{Critères de sélection}

\begin{itemize}
    \item \textbf{Précision recherchée} : une tâche courte nécessitera une unité fine (cmh, dmh) alors qu'une tâche longue pourra utiliser la centiheure.
    \item \textbf{Domaine étudié} : le temps machine, le temps manuel, le temps d'assemblage n'ont pas les mêmes ordres de grandeur.
    \item \textbf{Facilité des calculs} : les unités décimales simplifient les opérations arithmétiques.
    \item \textbf{Minimisation des conversions} : il est préférable d'utiliser une seule unité pour l'ensemble de l'étude.
    \item \textbf{Technique d'obtention du temps} : le chronométrage direct, la vidéo, les tables MTM imposent parfois l'unité.
    \item \textbf{Compatibilité avec les systèmes informatiques} : la GPAO (Gestion de Production Assistée par Ordinateur) peut imposer une unité particulière.
\end{itemize}

\subsubsection{Exemple de choix}

\begin{exemple}
On vise à déterminer la charge d'un opérateur. La précision recherchée est la seconde. La GPAO qui fera les calculs de charge utilise des ch (centiheure). Les temps sont obtenus à partir de la vidéo. Ces temps sont exprimés en h:min:s. L'unité sera donc le ch. Il faudra faire les calculs de conversion pour tous les temps mesurés.
\end{exemple}

\subsection{Types de temps}

Les temps peuvent être classés selon trois dimensions : leur nature, leur fréquence et leur position.

\subsubsection{Nature des temps (relation du temps à la ressource)}

\begin{description}[labelwidth=4cm]
    \item[Temps manuel (Tm)] Temps pendant lequel l'opérateur est actif, réalisant des opérations manuelles.
    \item[Temps technologique (Tt)] Temps pendant lequel la machine travaille sans intervention de l'opérateur.
    \item[Temps technico-manuel (Ttm)] Temps pendant lequel l'opérateur et la machine travaillent simultanément (intervention sur machine en fonctionnement).
\end{description}

\begin{table}[ht]
\centering
\begin{tabular}{|c|c|c|}
\hline
\textbf{Nature du temps} & \textbf{Symbole} & \textbf{Exemple} \\
\hline
Temps manuel & Tm & Prendre une pièce, positionner, évacuer \\
Temps technologique & Tt & Séchage automatique, mise en forme machine \\
Temps technico-manuel & Ttm & Verrouiller un plateau, actionner une pédale \\
\hline
\end{tabular}
\caption{Les trois natures de temps}
\end{table}

\subsubsection{Fréquence des temps (relation du temps à la production)}

\begin{description}[labelwidth=4cm]
    \item[Temps cycliques] Séquences qui se répètent à chaque cycle de production. \\
    \textit{Exemple :} prendre la pièce, la positionner, l'usiner, l'évacuer.
    
    \item[Temps fréquentiels] Séquences qui ne se produisent pas à chaque cycle mais à intervalles réguliers ou aléatoires. \\
    \textit{Exemples :}
    \begin{itemize}
        \item Approvisionnement en matière première (1 fois par lot)
        \item Changement de bobine de fil (après épuisement)
        \item Réglage machine (changement de série)
        \item Contrôle qualité périodique
    \end{itemize}
\end{description}

\subsubsection{Position des temps (relation des temps entre ressources)}

Cette classification concerne l'articulation temporelle entre différentes ressources (opérateurs, machines).

\begin{description}[labelwidth=4cm]
    \item[Temps séquentiels] Les opérations se succèdent dans le temps sans chevauchement.
    \item[Temps simultanés] Les opérations se déroulent en parallèle sur différentes ressources.
    \item[Temps masqués] Certaines opérations peuvent être réalisées pendant qu'une autre ressource travaille automatiquement, ne contribuant pas à l'allongement du cycle.
\end{description}

\subsection{Application pratique : Simogramme}

Pour visualiser et analyser la position des temps, on utilise un outil graphique appelé \textbf{simogramme}.

\subsubsection{Définition et objectifs}

Le simogramme est une représentation visuelle qui permet de :
\begin{itemize}
    \item Visualiser un cycle de production par type de temps (manuel, technologique, technico-manuel)
    \item Étudier la mise en parallèle de ressources (machines, opérateurs)
    \item Analyser différentes solutions d'organisation
    \item Identifier les opportunités d'optimisation
\end{itemize}

\subsubsection{Étapes de construction d'un simogramme}

\begin{enumerate}
    \item Identifier les séquences ou opérations
    \item Recueillir les durées des opérations
    \item Identifier les ressources mobilisées
    \item Repérer les temps manuels (Tm), technologiques (Tt) et technico-manuels (Ttm)
    \item Placer les ressources en ordonnée
    \item Placer l'échelle des temps en abscisse
\end{enumerate}

\subsubsection{Exemple : Atelier de mise en plis de pantalon}

\begin{table}[ht]
\centering
\begin{tabular}{|l|c|}
\hline
\textbf{Opération} & \textbf{Durée (dmh)} \\
\hline
Prendre et poser pantalon sur la presse & 5 (Tm) \\
Positionner le pantalon sur la presse & 12 (Tm) \\
Verrouiller plateau, actionner pédale vapeur & 8 (Ttm) \\
Mise en forme automatique & 60 (Tt) \\
Actionner pédale de séchage & 5 (Ttm) \\
Séchage automatique & 30 (Tt) \\
Déverrouiller plateau & 5 (Ttm) \\
Évacuer le pantalon & 10 (Tm) \\
\hline
\textbf{Temps total du cycle} & \textbf{135 dmh} \\
\hline
\end{tabular}
\caption{Décomposition des temps pour un cycle de pressage}
\end{table}

\subsubsection{Analyse d'optimisation}

L'analyse du simogramme permet d'envisager des optimisations, comme l'augmentation du nombre de machines confiées à un opérateur.

\begin{itemize}
    \item \textbf{Configuration avec 1 machine} : temps de cycle = 135 dmh
    \item \textbf{Production horaire} : $10\,000 \div 135 = 74$ pantalons/heure
    \item \textbf{Configuration avec 2 machines} : avec un temps de déplacement supplémentaire de 5 dmh entre les machines
    \item \textbf{Nouveau temps de cycle} : 165 dmh
    \item \textbf{Production horaire} : $(10\,000 \div 165) \times 2 = 120$ pantalons/heure
    \item \textbf{Gain} : $120 - 74 = 46$ pantalons/heure (+62\%)
\end{itemize}

\subsection{Synthèse}

\begin{center}
\fbox{\begin{minipage}{0.9\textwidth}
\textbf{Points clés à retenir :}

\begin{itemize}
    \item Les unités de temps doivent être choisies en fonction de la précision et des outils utilisés
    \item Les temps se classent selon leur nature (manuel, technologique, technico-manuel)
    \item La fréquence distingue les temps cycliques des temps fréquentiels
    \item La position des temps (séquentielle, simultanée, masquée) influence l'optimisation
    \item Le simogramme est un outil essentiel pour visualiser et optimiser les cycles
\end{itemize}
\end{minipage}}
\end{center}

% ==================== FIN DE LA PARTIE 1 ====================

% ==================== PARTIE 2 : CONDITIONS PRÉALABLES ====================
\section{Conditions Préalables au Chronométrage}

Avant d'effectuer un chronométrage valable, il est essentiel de s'assurer que trois éléments sont stabilisés : l'opérateur, le poste de travail et le chronométreur lui-même.

\subsection{Stabilisation de l'opérateur}

La stabilisation de l'opérateur correspond à la régularité de son travail.

\subsubsection{Caractéristiques d'un opérateur stabilisé}

\begin{itemize}
    \item Les temps d'exécution d'un travail répétitif présentent des écarts négligeables
    \item L'opérateur travaille à une vitesse régulière, aboutissant à un rendement constant
    \item Le mode opératoire est parfaitement maîtrisé et respecté
    \item L'opérateur utilise son matériel avec efficacité et automatisme
    \item La qualité du travail effectué est constante et conforme aux exigences
\end{itemize}

\subsubsection{La courbe d'apprentissage}

\begin{definition}
La \textbf{courbe d'apprentissage} modélise la réduction du temps nécessaire pour produire une unité à mesure que l'expérience augmente. À chaque doublement du nombre de répétitions de la tâche, sa durée de réalisation diminue d'un pourcentage constant $p$, appelé coefficient d'apprentissage.
\end{definition}

\paragraph{Loi de Wright (1963)}

\[
T_{n} = T_1 \cdot n^{e} \quad \text{avec} \quad e = \frac{\ln p}{\ln 2}
\]

\paragraph{Loi de Caquot}

\[
T_n = T_i \cdot \left(\frac{Q_n}{Q_i}\right)^{1/4}
\]

\subsection{Stabilisation du poste de travail}

Le poste de travail doit être dans des conditions normales et stables.

\subsection{Stabilisation du chronométreur}

Le chronométreur doit maîtriser les techniques de mesure et le jugement d'allure.

\subsubsection{Le jugement d'allure (JA)}

\begin{definition}
Le \textbf{jugement d'allure} est l'action par laquelle un observateur entraîné définit l'allure d'un opérateur par rapport à une allure de référence (allure 100).
\end{definition}

\subsection{Synthèse des conditions préalables}

\begin{center}
\fbox{\begin{minipage}{0.9\textwidth}
\textbf{Conditions préalables à vérifier avant tout chronométrage :}

\begin{itemize}
    \item[$\checkmark$] L'opérateur a dépassé sa phase d'apprentissage (courbe stabilisée)
    \item[$\checkmark$] Le poste de travail est aménagé, alimenté et fonctionne normalement
    \item[$\checkmark$] Les critères de qualité sont clairement définis et respectés
    \item[$\checkmark$] Le chronométreur maîtrise les techniques de mesure et le jugement d'allure
    \item[$\checkmark$] Les conditions d'observation sont favorables (absence d'interruptions)
    \item[$\checkmark$] La plage horaire choisie est représentative de l'activité normale
\end{itemize}
\end{minipage}}
\end{center}

% ==================== FIN DE LA PARTIE 2 ====================

\end{document}